\part{Examples from holomorphic Poisson structures: Serre-type vanishing, duality, the surface formula}
\setcounter{section}{0}\renewcommand{\thesection}{I.\arabic{section}}
\noindent\textit{Part I of the series; originally published as a separate note (\url{https://hylt-math.github.io/2026/11/22/generalized-geometry-poisson/}). Within this paper, Part-internal references use the numbering I.$\ast$.$\ast$.}\par
%--------------------------------------------------------------
\section{Notation}\label{p1:sec:notation}
%--------------------------------------------------------------
Overly general symbols (wedge product, differentials, etc.) are omitted.

\begin{longtable}{>{\raggedright\arraybackslash}p{3.4cm}>{\raggedright\arraybackslash}p{11.6cm}}
\toprule Symbol & Meaning\\ \midrule\endhead
$X$ & compact complex manifold, $\dim_\C X=n$ ($n$ is even in the log-symplectic case)\\
$T,\ T^*,\ \Omega^p$ & holomorphic tangent bundle, holomorphic cotangent bundle, sheaf of holomorphic $p$-forms. $K_X=\Omega^n$\\
$\bT=T\oplus T^*$ & generalized tangent bundle. Pairing $\langle X+\xi,Y+\eta\rangle=\tfrac12(\xi(Y)+\eta(X))$, Courant (Dorfman) bracket $[X+\xi,Y+\eta]=[X,Y]+\mathcal L_X\eta-\iota_Yd\xi$\\
$\Jc$ & generalized complex structure: $\Jc^2=-1$, preserves the pairing, and its $\pm i$ eigenbundles are closed under the Courant bracket\\
$L$ & eigenbundle of $\Jc$ (maximal isotropic subbundle of $\bT_\C$, a complex Lie algebroid). $L^*\cong\Lb$\\
$\sigma$ & holomorphic Poisson structure: $\sigma\in H^0(X,\wedge^2T)$, $[\sigma,\sigma]=0$ (Schouten bracket)\\
$\sigma^\sh$ & $T^*\to T$, $\xi\mapsto\sigma(\xi,\cdot)$\\
$\Jc_\sigma$ & generalized complex structure determined by $\sigma$ ($\beta$-transform). $L_\sigma=T^{0,1}\oplus\{\sigma^\sh\xi+\xi:\xi\in T^{*1,0}\}$\\
$A=(T^*)_\sigma$ & holomorphic Lie algebroid: anchor $\sigma^\sh$, bracket $[a,b]_\sigma=\mathcal L_{\sigma^\sh a}b-\mathcal L_{\sigma^\sh b}a-d\sigma(a,b)$\\
$D=\sum D_i$ & degeneracy divisor of $\sigma$ (reduced; simple normal crossings, or only ordinary double points for surfaces)\\
$T_D=T_X(-\log D)$,\ $B=\Omega^1(\log D)$ & logarithmic tangent and cotangent bundles. $T_D$ is a Lie algebroid, $B=T_D^*$\\
$U$ & $X\setminus D$\\
log-symplectic & $\sigma\in H^0(\wedge^2T_D)$ and $\sigma^\sh:B\to T_D$ is an isomorphism\\
$L$ (as a line bundle) & to be distinguished from the eigenbundle of $\bT_\C$ by context. The line bundle is often written $E$, $L_0$, $\mathcal O(m)$, etc.\ (\S\ref{p1:sec:setting})\\
Poisson module & holomorphic vector bundle $(E,\nabla)$ with a flat $A$-connection $\nabla$. The same thing as a generalized holomorphic (GH) bundle (Prop.~\ref{p1:prop:ident})\\
$V$ & holomorphic vector field defined by $\nabla_\beta e=\beta(V)e$ for a local frame $e$ (Poisson vector field, Lemma~\ref{p1:lem:flat})\\
$r_i$ & residue (a constant) of the logarithmic connection along $D_i$\\
$C(A;E)$ & $(\wedge^bT\otimes E,\ d_E=[\sigma,\cdot]+V\wedge)$ (Chevalley--Eilenberg complex of $A$)\\
$\HH^k(E)$ & $H^k(X;C(A;E))$: generalized Dolbeault cohomology (Prop.~\ref{p1:prop:ident})\\
$E_1^{a,b}$ & $H^a(X,\wedge^bT\otimes E)$: $E_1$ term of the spectral sequence converging to $\HH^{a+b}$ ($a$ is called the transversal degree)\\
$\mathrm{DR}_{\log}(E)$ & logarithmic de Rham complex $(\Omega^\bullet(\log D)\otimes E,\nabla^{\log})$\\
$Q$ & quotient complex $C(A;E)/C(T_D\text{-subcomplex})$, supported on $D$\\
$\mu$ & (surfaces) number of ordinary double points of $D$\\
$c,\ \rho,\ b_1$ & number $c$ of irreducible components of $D$, dimension $\rho$ of the span of $[D_i]\in H^2(X;\C)$, $b_1=c-\rho$\\
$j_1:U\hookrightarrow X',\ j_2:X'\hookrightarrow X$ & open embeddings (Thm.~\ref{p1:thm:main2})\\
$L_U$ & local system on $U$ of $\nabla$-flat sections (rank 1, monodromy $e^{2\pi ir_i}$ around $D_i$)\\
$P_E(t)$ & $\sum_k\dim\HH^k(E)\,t^k$: Poincar\'e polynomial\\
$[\,\cdot\,]$ (tables) & shorthand for $[\dim\HH^0,\dim\HH^1,\dots,\dim\HH^{2n}]$\\
\bottomrule
\end{longtable}

%--------------------------------------------------------------
\section{The classical three treasures: statements, hypotheses, and where they enter the proofs}\label{p1:sec:classical}
%--------------------------------------------------------------
To make the target of the generalization clear, we recall the standard statements of the classical versions and the skeletons of their proofs. The content is standard textbook material and is not new.

\begin{thm}[Hirzebruch--Riemann--Roch]\tC\label{p1:thm:HRR}
Let $M$ be a compact complex manifold and $E$ a holomorphic vector bundle. Then
$\chi(M,E)=\sum_q(-1)^q\dim H^q(M,\mathcal O(E))=\int_M\ch(E)\Td(T^{1,0}M).$
\src{Hirzebruch, \emph{Neue topologische Methoden in der algebraischen Geometrie}, Springer 1956 (projective algebraic case); Atiyah--Singer, Ann.\ Math.\ 87 (1968) 484--530, 531--545, 546--604 (general compact complex manifolds, index of the Dolbeault complex). Textbook: Huybrechts, \emph{Complex Geometry}, Springer 2005, \S5.1.}
\end{thm}
\begin{hyptab}
Compactness & The cohomology is finite-dimensional. & The Dolbeault complex is elliptic, hence Fredholm, so the index is defined.\\\midrule
Holomorphic structure & Condition $\dbar_E^2=0$ for a complex to exist. & Definition of the complex $(\Omega^{0,\bullet}(E),\dbar_E)$.\\\midrule
(Form of the conclusion) $\ch(E)$ & The index is the integral of topological data of the principal symbol. & Atiyah--Singer index theorem. $\Td$ arises from the symbol, and $E$ enters not through its rank but through $\ch(E)$.
\end{hyptab}

\begin{thm}[Serre's vanishing theorem]\tC\label{p1:thm:Serre}
Let $X$ be a projective scheme over a field, $\mathcal O(1)$ a very ample line bundle and $\mathcal F$ a coherent sheaf. Then there is $m_0$ such that for $m\ge m_0$: (a) $H^q(X,\mathcal F(m))=0\ (q>0)$; (b) $\mathcal F(m)$ is generated by global sections.
\src{Serre, \emph{Faisceaux alg\'ebriques coh\'erents}, Ann.\ Math.\ 61 (1955) 197--278; modern proof in Hartshorne, \emph{Algebraic Geometry}, GTM 52, Springer 1977, Theorem III.5.2.}
\end{thm}
\begin{hyptab}
Coherence & Finite generation gives a surjection $\bigoplus\mathcal O(-N_i)\twoheadrightarrow\mathcal F$. & To reduce to the known computation on $\mathbb P^N$.\\\midrule
Ampleness & Twisting $\mathcal F(m)$ by large $m$ kills higher cohomology. & It is the conclusion itself.\\\midrule
Projectivity & $X$ embeds as a closed subscheme of $\mathbb P^N$, and $H^q(\mathbb P^N,\mathcal O(m))$ can be computed directly by the Cech complex. & When running the descending induction on $q$ (vanishing for $q>\dim X$) on the long exact sequence for the kernel $\mathcal K$ of the surjection.
\end{hyptab}

\begin{thm}[Kodaira--Akizuki--Nakano vanishing]\tC\label{p1:thm:KAN}
Let $M$ be a compact K\"ahler manifold and $E$ a positive holomorphic Hermitian line bundle. If $p+q>n$ then $H^q(M,\Omega^p\otimes E)=0$. In particular $H^q(M,K_M\otimes E)=0\ (q>0)$.
\src{Kodaira, Proc.\ Nat.\ Acad.\ Sci.\ USA 39 (1953) 1268--1273; Akizuki--Nakano, Proc.\ Japan Acad.\ 30 (1954) 266--272. Textbooks: Griffiths--Harris, \emph{Principles of Algebraic Geometry}, Wiley 1978, Ch.~1 \S2; Huybrechts \S5.2.}
\end{thm}
\begin{hyptab}
Compactness & (i) Hodge theory ($H^q\cong\ker\Delta_{\dbar_E}$); (ii) no boundary terms in integration by parts. & Existence of harmonic representatives, and the integral estimate below.\\\midrule
K\"ahler & The K\"ahler identity $[\Lambda,\dbar]=-i\partial^*$. & Derivation of the Bochner--Kodaira--Nakano identity $\Delta_{\dbar_E}=\Delta_{\partial_E}+[i\Theta_E,\Lambda]$.\\\midrule
Positive ($i\Theta_E>0$) & The curvature term $[i\Theta_E,\Lambda]$ is positive definite for $p+q>n$. & For $s\in\ker\Delta_{\dbar_E}$, in $0=\|\partial_Es\|^2+\|\partial_E^*s\|^2+\langle[i\Theta_E,\Lambda]s,s\rangle$ the last term is $>0$ unless $s=0$.
\end{hyptab}

\begin{thm}[Kodaira's embedding theorem]\tC\label{p1:thm:Kod}
Let $M$ be a compact complex manifold and $L$ a positive holomorphic line bundle. Then $M$ is projective, and for sufficiently large $m$ the map $M\to\mathbb P(H^0(M,L^m)^*)$ defined by $|L^m|$ is an embedding.
\src{Kodaira, Ann.\ Math.\ 60 (1954) 28--48; Griffiths--Harris Ch.~1 \S4; Huybrechts \S5.3.}
\end{thm}
\begin{hyptab}
$L$ positive & The curvature form $i\Theta_L$ is a K\"ahler form, so $M$ is K\"ahler. & Ensures the situation where Kodaira vanishing can be used.\\\midrule
Compactness & $H^0(M,L^m)$ is finite-dimensional, so the map to projective space is defined. & Construction of the map.\\\midrule
Enough holomorphic sections & The standard proof blows up a point $x$ to get $\tilde M$ with exceptional divisor $E_x$ and deduces surjectivity of $H^0(L^m)\to L^m\otimes\mathcal O/\mathfrak m_x^2$ from vanishing (Kodaira vanishing) of $H^1(\tilde M,\pi^*L^m\otimes\mathcal O(-2E_x))$. & Separation of points and of tangent vectors. The point is that no dimension estimate from Riemann--Roch is needed.
\end{hyptab}

\begin{rem}\tK\ (Hints for the generalization)
Classically, four items share the work: (1) Hodge theory (ellipticity), (2) K\"ahler identities $\Rightarrow$ BKN identity, (3) positivity of curvature, (4) abundance of global sections. In generalized complex geometry, (1) holds as the elliptic complex $(\wedge L^*\otimes E,d_L)$ (Prop.~\ref{p1:prop:ellip}). The line-bundle-coefficient versions of (2) and (3) have not yet been written in the generalized K\"ahler framework. (4) does not hold in the classical form, because generalized holomorphic sections become Casimir sections of a Poisson module (\S\ref{p1:sec:ghfun}). The present note studies these points in the simplest non-trivial examples (type-jumping $\beta$-transforms).
\end{rem}

%--------------------------------------------------------------
\section{Basics of generalized complex geometry}\label{p1:sec:gc}
%--------------------------------------------------------------
\subsection{Definitions}
\begin{defn}[Generalized complex structure]\tC\label{p1:def:GC}
A map $\Jc:\bT\to\bT$ with $\Jc^2=-1$ preserving the pairing is called a generalized almost complex structure. If the $+i$-eigenbundle $L\subset\bT_\C$ is closed under the Courant bracket, $\Jc$ is called integrable (a generalized complex structure).
\src{Hitchin, \emph{Generalized Calabi--Yau manifolds}, Q.\ J.\ Math.\ 54 (2003) 281--308, arXiv:math/0209099; Gualtieri, \emph{Generalized complex geometry}, Ann.\ Math.\ 174 (2011) 75--123, arXiv:math/0703298, \S3 (type Def.~3.5, canonical bundle Def.~3.7, generalized Calabi--Yau Def.~3.17, Picard group Def.~3.18; numbers checked in the source); Gualtieri's thesis arXiv:math/0401221.}
\end{defn}
\begin{ex}\tC
Complex structure $I$: $L=T^{0,1}\oplus T^{*1,0}$ (type $n$). Symplectic structure $\omega$: $L=\{X-i\omega(X)\}$ (type $0$).
\src{Hitchin 2003 (introduction); Gualtieri 2011, \S3 (examples of complex and symplectic structures; no specific example numbers could be confirmed; Def.~4.1 in \S4 is the definition of a regular point, so the ``Examples 4.1--4.2'' cited in v1 do not exist) \tQ.}
\end{ex}
\begin{defn}[Type]\tC
$\mathrm{type}(\Jc)_p=\mathrm{codim}_\C\,\pi_T(L_p)$ ($\pi_T:\bT\to T$). Points where the type is locally constant are called regular points. The type can jump in general.
\src{Gualtieri 2011, \S3, and Abouzaid--Boyarchenko, \emph{Local structure of generalized complex manifolds}, J.\ Symplectic Geom.\ 4 (2006) 43--62, arXiv:math/0511049.}
\end{defn}
\begin{prop}[Elliptic complex]\tC\label{p1:prop:ellip}
$L$ is a complex Lie algebroid (anchor $\pi_T|_L$), and via $L^*\cong\Lb$ the exterior algebra $\wedge^\bullet L^*$ carries the Lie algebroid differential $d_L$. The complex $(\wedge^\bullet L^*,d_L)$ is elliptic, and on a compact manifold the generalized Dolbeault cohomology $H^k(\wedge L^*,d_L)$ is finite-dimensional. Using the spinor bundle $K$, under $\wedge^{n-k}\Lb\cdot K\cong U^k$ the differential $d_L$ corresponds to $\dbar:U^k\to U^{k-1}$.
\src{Gualtieri 2011, Prop.~3.13 ($(\wedge L^*,d_L)$ is elliptic), Cor.~3.14 (finite-dimensional cohomology when compact). In the proof in \S\ref{p1:sec:duality}, for $\Jc_\sigma$ ellipticity is verified directly from the principal symbol (\tP).}
\end{prop}

\subsection{The $\beta$-transform $\Jc_\sigma$ and generalized holomorphic bundles}\label{p1:sec:setting}
Let $X$ be a complex manifold and $\sigma$ a holomorphic Poisson structure.

\begin{prop}\tC\label{p1:prop:Jsigma}
(a) $L_\sigma=T^{0,1}\oplus\{\sigma^\sh\xi+\xi:\xi\in T^{*1,0}\}$ is a maximal isotropic subbundle of $\bT_\C$ and is integrable. The corresponding generalized complex structure $\Jc_\sigma$ has type $n-\mathrm{rk}\,\sigma^\sh_p$ at $p\in X$ ($\mathrm{rk}$ is the complex rank). In particular it has type 0 (symplectic type) where $\sigma$ is non-degenerate, and type $n$ (complex type) where $\sigma$ vanishes.
(b) $(\Xi,\xi)\mapsto(\Xi+\sigma^\sh\xi,\xi)$ is an isomorphism of complex Lie algebroids $T^{0,1}\bowtie(T^{1,0})^*_\sigma\to L_\sigma$ (the left side is a matched pair, written $T^{0,1}\bowtie A^{1,0}$ in LGSX).
\src{(a) Gualtieri, \emph{Branes on Poisson varieties}, arXiv:0710.2719, Prop.~5; Laurent-Gengoux--Sti\'enon--Xu (LGSX below), \emph{Holomorphic Poisson manifolds and holomorphic Lie algebroids}, IMRN 2008, Art.\ ID rnn088, arXiv:0707.4253, Thm.~2.7 and Rem.~2.8 (there written $\Jc_\pi$, $\pi=\pi_R+i\pi_I$). The type formula is the author's reorganization, consistent with Gualtieri's examples ($\mathbb{CP}^2$ etc.) (\tQ). (b) LGSX Thm.~4.21 and Lemma~4.22 (the factor 4 in $\Jc_{4\pi}$ comes from the normalization of $\pi=\pi_R+i\pi_I$, LGSX Prop.~3.7). Gualtieri's paper and LGSX differ in the choice of sign of eigenvalues ($\ker(\Jc\mp i)$), but the two are exchanged by complex conjugation, so the statements below are unaffected.}
\end{prop}

\begin{defn}[Generalized holomorphic bundle]\tC
A complex vector bundle $E$ equipped with a flat $L$-connection $\dbar_E:\Gamma(E)\to\Gamma(L^*\otimes E)$ ($\dbar_E^2=0$).
\src{Gualtieri 2011, \S3 (generalized holomorphic bundles are introduced in prose after Thm.~3.15, with no numbered definition; the complex case is Example 3.19; Def.~3.18 is the definition of the Picard group); Gualtieri, \emph{Branes on Poisson varieties}, \S4.1.}
\end{defn}
\begin{defn}[Poisson module]\tC
A holomorphic vector bundle $E$ with a flat $A$-connection $\nabla:\Gamma(E)\to\Gamma(A^*\otimes E)$, i.e.\ satisfying $\nabla_{\xi}(fs)=f\nabla_\xi s+(\sigma^\sh\xi)(f)s$ and $[\nabla_a,\nabla_b]=\nabla_{[a,b]_\sigma}$ ($a,b$ holomorphic 1-forms). Here $A=(T^*)_\sigma$.
\src{LGSX, Def.~4.15 and Lemma~4.16; Gualtieri, \emph{Branes on Poisson varieties}, \S2--4.}
\end{defn}

\begin{prop}[Generalized holomorphic bundles $=$ Poisson modules, and the Dolbeault model]\label{p1:prop:ident}
Let $(X,\sigma)$ be as above.
\begin{enumerate}
\item[(a)] A generalized holomorphic bundle with respect to $\Jc_\sigma$ is the same thing as a holomorphic vector bundle with a Poisson module structure. \tC\ Gualtieri, \emph{Branes on Poisson varieties}, Prop.~8; LGSX Lemma~4.16. A direct proof is given below (\tP). For line bundles, by Prop.~9 of the same paper, the obstruction is a class $\sigma\cdot\alpha$ in $H^1(T^{1,0})$ together with a class in $H^2_\sigma$, and when structures exist they form an affine space over $H^1_\sigma(X)=\HH^1(\mathcal O)$. \tC
\item[(b)] Since $L^*=T^{*0,1}\oplus T^{1,0}$, we have $\Gamma(\wedge^kL^*\otimes E)=\bigoplus_{a+b=k}A^{0,a}(\wedge^bT\otimes E)$ and $d_L=\dbar+(-1)^ad_A$, where $d_A=[\sigma,\cdot]+V\wedge$ acts on the coefficients in a holomorphic frame. Hence
\[H^k_{GC}(E)=H^k\bigl(\Gamma(\wedge^\bullet L^*\otimes E),d_L\bigr)=\mathbb H^k\bigl(X;(\wedge^\bullet T\otimes E,\ [\sigma,\cdot]+V\wedge)\bigr)=:\HH^k(E).\]
\tC\ This is LGSX Thm.~4.19 ($H^*(A,E)=H^*(T^{0,1}\bowtie A^{1,0},E)$, with $H^*(A,E)$ defined as sheaf cohomology in Def.~4.17). \tP\ The structure of the proof is given below, but it depends on Mokri's theorem on matched pairs (Matched pairs of Lie algebroids, Glasgow Math.\ J.\ 39 (1997) 167--181, doi:10.1017/S0017089500032055; year and DOI checked, volume/pages from memory \tQ; LGSX Lemma~4.16: a complex vector bundle is an $A$-module iff it is a $T^{0,1}X\bowtie A^{1,0}$-module; the proof of Thm.~4.19 is only stated to be similar to that of Thm.~4.11) and on LGSX Lemma~4.22 (not re-derived).
\end{enumerate}
\end{prop}
\begin{proof}[Proof of (a) (sketch)]
($\Leftarrow$) Given a Poisson module structure and a holomorphic frame $\{e_s\}$, put $\nabla_{\bar\Xi}\sum f_se_s=\sum(\bar\Xi f_s)e_s$ (holomorphic structure) and $\nabla_\xi\sum f_se_s=\sum((\sigma^\sh\xi)f_s)e_s+f_s\nabla_\xi e_s$ ($\xi\in A^{1,0}$). This is the unique extension forced by the Leibniz rule, and it is well defined because $\nabla_\xi e_s$ is holomorphic. Flatness: the $(\bar\Xi,\bar\Upsilon)$ components follow from integrability of the holomorphic structure; the $(\bar\Xi,\xi)$ components are equivalent to ``$\nabla_\xi$ maps holomorphic sections to holomorphic sections'', which holds by construction; the $(\xi,\eta)$ components are the $C^\infty$-linear extension of flatness of the Poisson module (for holomorphic $\xi,\eta$) (the curvature is a tensor).
($\Rightarrow$) Restricting a flat $L$-connection to $T^{0,1}$ gives a holomorphic structure (Koszul--Malgrange); flatness of mixed type shows that $\nabla_\xi$ maps holomorphic sections to holomorphic sections, and the module conditions follow from the Leibniz rule and flatness.\qedhere
\end{proof}
\begin{proof}[Structure of the proof of (b)]
By the general theory of matched pairs (Mokri), the Chevalley--Eilenberg complex of $L\cong T^{0,1}\bowtie A^{1,0}$ (Prop.~\ref{p1:prop:Jsigma}(b)) is the total complex of $C(T^{0,1};\,C(A^{1,0};E))$. Here the $T^{0,1}$ direction is the Dolbeault resolution of $\wedge^bT\otimes E$ (differential $\dbar$), and the $A^{1,0}$ direction is $C(A;E)=(\wedge^bT\otimes E,\ [\sigma,\cdot]+V\wedge)$. The two differentials anti-commute with the Koszul sign $(-1)^a$. The total complex is a fine resolution of the complex $(\wedge^\bullet T\otimes E,d_A)$, so its cohomology is the hypercohomology.\qedhere
\end{proof}

\begin{rem}[\tC\ Note on integrability]
A notion of Lie groupoid integrating the complex Lie algebroid $L=L_\sigma$ itself was not found in the literature surveyed. Thms.~3.17 and 3.22 of LGSX arXiv:0803.2031 relate integrability of a holomorphic Lie algebroid $A$ to that of the real Poisson structures $\mathrm{Re}\,\sigma$, $\mathrm{Im}\,\sigma$, but do not treat $T^{0,1}\bowtie A^{1,0}$ (Rem.~4.14 of that paper only poses it as a question). Thm.~7.1 of Bailey--Gualtieri, arXiv:1611.03850 (title \tQ), states the equivalence between a generalized complex structure integrating to a weak holomorphic symplectic groupoid (a Morita-class notion) and integrability of its underlying real Poisson structure. \tP\ For the smooth cubic $\sigma$ on $\mathbb{CP}^2$, the leaves of $\mathrm{Im}\,\sigma$ are the open leaf $U$ (isotropy $0$) and the points of $D$ (isotropy $\mathrm{aff}(1,\C)$, with trivial centre); by Crainic--Fernandes, Ann.\ Math.\ 157 (2003) 575--620 (arXiv:math/0105033) Cor.~5.2 (``if all $N_x$ are trivial, then integrable'') $\mathrm{Im}\,\sigma$ is integrable. The hypotheses of Bailey--Gualtieri were checked only through the summary obtained \tQ.
\end{rem}

\subsection{On generalized holomorphic sections}\label{p1:sec:ghfun}
\begin{obs}\tP\label{p1:obs:casimir}
A generalized holomorphic section ($\dbar_Es=0$) of a Poisson module $(E,\nabla)$ is a flat (Casimir-type) section: $\HH^0(E)=\{s\in H^0(X,E):\nabla_\xi s=0\ (\forall\xi)\}$. Near a point of type $0$ ($\sigma$ non-degenerate), $\nabla$ corresponds via $\sigma^\sh$ to an ordinary flat connection, and flat sections are locally constant.
\end{obs}
\begin{proof}
$\HH^0=\ker(d_A:H^0(E)\to H^0(T\otimes E))$, and $d_As=(\xi\mapsto\nabla_\xi s)$ (Prop.~\ref{p1:prop:ident}(b)).
\end{proof}
\begin{rem}\tK
This means that, in the regular case of non-zero type as well (the Pal--Poddar framework), generalized holomorphic sections are locally constant in the leaf directions, so they cannot give an embedding of $M$ into projective space. A generalized Kodaira embedding theorem will have to be reformulated as a map to the leaf space (or to the quotient separated by Casimir sections). The present note does not treat the generalization of the Kodaira embedding theorem itself.
\end{rem}

%--------------------------------------------------------------
\section{The log-symplectic case: Poisson modules $=$ logarithmic flat connections}\label{p1:sec:logsymp}
%--------------------------------------------------------------
\subsection{Setting}
Let $X$ be a compact complex manifold ($\dim_\C X=n=2k$), $\sigma$ a holomorphic Poisson structure, and $D\subset X$ a reduced divisor (simple normal crossings; for surfaces, ordinary double points are also allowed).
Put $T_D=T_X(-\log D)$, $B=\Omega^1(\log D)$.

\begin{defn}[Log-symplectic]\tC
$\sigma$ is called log-symplectic along $D$ if $\sigma\in H^0(\wedge^2T_D)$ and $\sigma^\sh:B\to T_D$ is an isomorphism (equivalently, $\omega:=\sigma^{-1}\in H^0(\Omega^2(\log D))$ is non-degenerate). Then $\mathrm{div}(\sigma^k)=D$, i.e.\ $-K_X=\mathcal O(D)$.
\src{Goto, \emph{Unobstructed deformations of generalized complex structures induced by $C^\infty$ logarithmic symplectic structures}, arXiv:1501.03398 (journal and pages unverified \tQ); Gualtieri--Li, \emph{Symplectic groupoids of log symplectic manifolds}, IMRN 2014, arXiv:1206.3674; Pym, \emph{Constructions and classifications of projective Poisson varieties}, Lett.\ Math.\ Phys.\ 108 (2018) 573--632 (checked). The journal of Gualtieri--Li and the publication data of Goto are unverified \tQ.}
For surfaces this means that $\sigma\in H^0(-K_X)$ has a reduced zero divisor with only ordinary double points. Locally $\sigma=h\,\partial_1\wedge\partial_2$ with $h=zu$ (smooth points) or $h=z_1z_2u$ (double points), $u$ a unit.
\end{defn}

\subsection{Statement}
\begin{thm}[Poisson modules $=$ logarithmic flat connections]\tP\label{p1:thm:main1}
Let $(X,D,\sigma)$ be log-symplectic (any even dimension, $D$ simple normal crossings) and $E$ a holomorphic line bundle carrying a Poisson module structure.
\begin{enumerate}
\item[(a)] For a local frame $e$, a holomorphic vector field $V$ is defined by $\nabla_\beta e=\beta(V)e$, and $[\sigma,V]=0$ (a Poisson vector field). Under a change of frame $e\mapsto ge$, $V$ changes only by a Hamiltonian term of the form $X_{\log g}$.
\item[(b)] $V$ is automatically logarithmic: $V\in T_D$. Hence the $A$-module structure extends uniquely to a $B$-module structure, and via the isomorphism $\sigma^\sh:B\cong T_D$ it is the same thing as a flat connection $\nabla^{\log}:E\to\Omega^1(\log D)\otimes E$ with logarithmic poles.
\item[(c)] The residue of $\nabla^{\log}$ along $D_i$ is a constant $r_i\in\C$, and $c_1(E)=-\sum_ir_i[D_i]$ in $H^2(X;\C)$.
\item[(d)] Near a point of $D$ one can choose a frame with $\nabla^{\log}e=\sum_ir_i\,\dlog z_i\,e$.
\item[(e)] The subcomplex $C(T_D;E)=(\wedge^bT_D\otimes E,d_E)\subset C(A;E)$ is isomorphic via $\wedge^b\sigma^\sh$ to the logarithmic de Rham complex $\mathrm{DR}_{\log}(E)$.
\end{enumerate}
\end{thm}
\begin{hyptab}
Log-symplectic ($\sigma^\sh:B\cong T_D$) & To identify the Lie algebroid structure $(T_D^*)_\sigma$ on $B$ with $T_D$, and to treat Poisson modules as ordinary flat connections (with logarithmic poles). & Lemma~\ref{p1:lem:iso}, Theorem (b), (e).\\\midrule
$D$ reduced, simple normal crossings & To estimate the pole order from each factor of the local equation $f=z_1\cdots z_s$ independently. & Lemma~\ref{p1:lem:log} (deducing $V(z_i)\in(z_i)$ from $V(f)\in(f)$).\\\midrule
$X$ compact & For the relation between $c_1(E)$ and the residues (an equality of global cohomology classes). & (c). That the residues are constant is local.\\\midrule
$\sigma$ Poisson ($[\sigma,\sigma]=0$) & $A$ becomes a Lie algebroid and $d_E^2=0$. & Lemma~\ref{p1:lem:flat}.
\end{hyptab}

\begin{lem}[Flatness $\Leftrightarrow$ Poisson vector field]\tP\label{p1:lem:flat}
For a local frame $e$ put $\Theta(\beta):=\nabla_\beta e/e$. Then $\nabla_a\nabla_be=(\sigma^\sh a)(\Theta b)e+\Theta(a)\Theta(b)e$. Hence flatness $[\nabla_a,\nabla_b]=\nabla_{[a,b]_\sigma}$ is equivalent to
$(\sigma^\sh a)(\Theta b)-(\sigma^\sh b)(\Theta a)=\Theta([a,b]_\sigma)$, i.e.\ $\Theta$ is a 1-cocycle of the Lichnerowicz complex: $\Theta(\beta)=\beta(V)$, $[\sigma,V]=0$.
In local coordinates with $\sigma=h\partial_1\wedge\partial_2$, $a=\Theta(dz_1)$, $b=\Theta(dz_2)$, the condition reads $h(a_1+b_2)=h_1a+h_2b$, i.e.\ $V(h)=h\,\mathrm{div}V$. The differential $d_E$ is, on $C^0$, $d_E(fe)=(-\sigma^\sh(df)+fV)e$, and for $c=X\otimes e$, $X=p\partial_1+q\partial_2$,
\[(d_Ec)(dz_1,dz_2)=h(p_1+q_2)-h_1p-h_2q+qa-pb.\]
\end{lem}
\begin{lem}[$V$ is logarithmic]\tP\label{p1:lem:log}
Let $f=z_1\cdots z_s$ be a local equation of $D$. From $L_V\sigma=0$ we get $L_V\sigma^k=0$. Since $\sigma^k=f'\partial_1\wedge\cdots\wedge\partial_n$ ($f'=f\cdot$unit), $V(f')-f'\,\mathrm{div}V=0$, hence $V(f)\in(f)$. Now $V(f)/f=\sum_iV(z_i)/z_i$, and the $i$-th term has a pole only along $\{z_i=0\}$. Since the whole is holomorphic, each $V(z_i)/z_i$ is holomorphic, i.e.\ $V(z_i)\in(z_i)$.
\end{lem}
\begin{lem}[$\sigma^\sh:B\to T_D$ is an isomorphism of Lie algebroids]\tP\ (standard)\label{p1:lem:iso}
$\sigma\in\wedge^2T_D$ is a non-degenerate Poisson structure on the Lie algebroid $T_D$. That $\omega=\sigma^{-1}$ is $d_{T_D}$-closed follows from $\omega$ being closed on $U$ (the inverse of a Poisson structure) and from the fact that a section of $\Omega^3(\log D)$ vanishing on a dense open set is $0$. For a non-degenerate Poisson structure coming from a closed 2-form, $\sigma^\sh$ preserves brackets. Hence $(T_D^*)_\sigma=B\cong T_D$.
\end{lem}
\begin{lem}[Residues, local normal form, $c_1$]\tP\ /\ \tC\label{p1:lem:res}
Let $\nabla^{\log}e=\theta e$ with $\theta\in\Omega^1(\log D)$; $\theta$ is closed by flatness. The residue $\theta_i|_{D_i}$ along $D_i$ is seen to be constant by restricting $d\theta$ to $D_i$. Since $\theta-\sum r_i\dlog z_i$ is a closed holomorphic 1-form, it equals $dg$, and for the frame $e'=e^{-g}e$ we get $\nabla^{\log}e'=\sum r_i\dlog z_i\,e'$ (this is (d)). The relation $c_1(E)=-\sum r_i[D_i]$ is the Esnault--Viehweg residue formula (\tC\ Esnault--Viehweg, \emph{Logarithmic de Rham complexes and vanishing theorems}, Invent.\ Math.\ 86 (1986) 161--194, \S2; Deligne, \emph{\'Equations diff\'erentielles \`a points singuliers r\'eguliers}, LNM 163, Springer 1970, Ch.~II. Sign check: for $\mathcal O(D)$ with the canonical section flat, the residue is $-1$ and $c_1=[D]$).
The set of all module structures with the same $c_1$ is a torsor under $H^0(\Omega^1(\log D))=\ker(\C^c\to H^2(X;\C))$ (Deligne), of dimension $c-\rho$. If $D$ is irreducible the structure is unique; if reducible, it forms an affine family of residue vectors.
\end{lem}
\begin{proof}[Proof of Theorem~\ref{p1:thm:main1}]
(a) Lemma~\ref{p1:lem:flat}. (b) By Lemma~\ref{p1:lem:log}, $\Theta$ extends to $B$, and since flatness is an equality on a dense open set it persists after the extension. Combine with Lemma~\ref{p1:lem:iso}. (c), (d) Lemma~\ref{p1:lem:res}. (e) Lemma~\ref{p1:lem:iso}.
\end{proof}

%--------------------------------------------------------------
\section{The simplest non-trivial example: $\mathbb{CP}^2$ and a smooth cubic}\label{p1:sec:cp2}
%--------------------------------------------------------------
Let $F$ be a smooth homogeneous cubic polynomial on $\mathbb C^3$ and $\Pi^{ij}=\epsilon_{ijk}\partial_kF$. Then $\Pi$ is a homogeneous Poisson structure (of weight 0) whose contraction with the Euler vector field is $0$, and it defines a holomorphic Poisson structure $\sigma\in H^0(\wedge^2T)=H^0(\mathcal O(3))$ on $\mathbb{CP}^2$. $\sigma$ vanishes to first order along $D=\{F=0\}$ (a smooth cubic curve). This is a standard construction (identification of the source \tQ; the vanishing order along $D$ was checked by a local computation). Below $E=\mathcal O(m)$.

\subsection{Which $\mathcal O(m)$ are generalized holomorphic}
\begin{prop}\tP\label{p1:prop:cp2mod}
For every $m\in\Z$ the bundle $\mathcal O(m)$ has exactly one Poisson module structure, and the residue of its logarithmic connection is $r=-m/3$.
\end{prop}
\begin{proof}
By Theorem~\ref{p1:thm:main1}(c), $r=-m/3$ (as $[D]=3H$). Existence: vanishing of the Atiyah class of $\mathcal O(m)$ in $H^1(\Omega^1(\log D))$ is equivalent to $c_1\in\C[D]$ (exact sequence $H^0(\mathcal O_D)\to H^1(\Omega^1)\to H^1(\Omega^1(\log D))$), and $m[H]\in\C\cdot3[H]$. The curvature is an element of $H^0(\Omega^2(\log D))=\C\omega$ with residue $dr=0$, hence zero (flat). Uniqueness: the difference of two logarithmic connections lies in $H^0(\Omega^1(\log D))=0$ (the connecting map sends $1\mapsto[D]=3H\ne0$).\qedhere
\end{proof}
Another view: for the Casimir $F$ of $\Pi$, $F^{m/3}$ gives a (multivalued) flat section of $\mathcal O(m)$ on $U=\mathbb{CP}^2\setminus D$. $\mathcal O(m)|_U$ is a rank-one local system $L_m$ with monodromy $e^{-2\pi im/3}$ around $D$. The period 3 corresponds to $\pi_1(U)=\Z/3$. \tP\ Locally $\nabla e=-\tfrac m3\dlog z\,e$.

\subsection{Result}
\begin{thm}\tP\label{p1:thm:cp2}
The dimensions $[\dim\HH^0,\dots,\dim\HH^4]$ of the generalized Dolbeault cohomology
\[\HH^*(\mathcal O(m))=\mathbb H^*\bigl(\mathbb{CP}^2;(\wedge^\bullet T(m),d_\sigma)\bigr)\]
are as follows.
\begin{center}\small
\begin{tabular}{lll}\toprule
$m$ & dimensions & Interpretation\\\midrule
$m\in3\Z_{\ge0}$ & $[1,0,2,0,0]$ & $H^*(U;\C)$. $\HH^0=\C F^{m/3}$ is a Casimir section\\
$m\in3\Z_{\le-6}$ & $[0,0,2,0,1]$ & $H^*_c(U;\C)$\\
otherwise (including $m=-3$) & $[0,0,3,0,0]$ & $H^*(U;L_m)$ ($m\notin3\Z$); $m=-3$ is special\\\bottomrule
\end{tabular}
\end{center}
The total Euler characteristic is always $3=\chi(\mathbb{CP}^2)$.
\end{thm}
\begin{hyptab}
$D$ smooth, $\sigma$ vanishing to first order along $D$ & The differential of the quotient complex $Q$ is a constant multiple of the isomorphism $d\sigma|_D:N_D\to\wedge^2T|_D$. & Step (3) of the proof. $Q$ is acyclic for $m\ne-3$.\\\midrule
$\pi_1(U)=\Z/3$ (Milnor fibre $\simeq$ a bouquet of 8 copies of $S^2$) & To determine the cohomology of $U$ explicitly. & Right-hand side of the table in step (2).\\\midrule
$m\ne-3$ & The differential $-(1+m/3)\cdot(d\sigma|_D)$ of $Q$ is not the zero map. & Only $m=-3$ is treated separately (step (4)).
\end{hyptab}
\begin{proof}[Steps of the proof]
(1) By Proposition~\ref{p1:prop:cp2mod} and Theorem~\ref{p1:thm:main1}(e), the subcomplex $C_{\log}=(\wedge^bT_D\otimes\mathcal O(m))\cong\mathrm{DR}_{\log}$.
(2) In a local frame $\nabla^{\log}e=r\,\dlog z\,e$ ($r=-m/3$). At the stalk at a point of $D$, the logarithmic de Rham complex acts on the monomial $z^ne$ by the factor $(n+r)=(n-m/3)$. The stalk cohomology is $[1,1]$ for $m\in3\Z_{\ge0}$ and $0$ otherwise. Hence for $m\in3\Z_{\ge0}$ the natural map $\mathrm{DR}_{\log}\to Rj_*L_m$ is a stalkwise quasi-isomorphism ($\mathbb H(C_{\log})=H^*(U;\C)$), and otherwise $j_!L_m\to\mathrm{DR}_{\log}$ is a quasi-isomorphism ($\mathbb H(C_{\log})=H_c^*(U;L_m)$). Topology of $U$: the Milnor fibre $F=1$ is a bouquet of $\mu=8$ copies of $S^2$, $U=\{F=1\}/\mu_3$, and $H^2=\C^8=2+3+3$ (trivial character and two non-trivial characters). Since $U$ is affine, $H^k=0$ for $k>2$. Hence $H^*(U)=[1,0,2,0,0]$, $H^*(U;L)=[0,0,3,0,0]$ for a non-trivial character, and by Verdier duality $H^*_c(U;\C)=[0,0,2,0,1]$.
(3) The quotient $Q=C(A;\mathcal O(m))/C_{\log}$ is a complex on $D$ with $Q^1=N_D(m)=\mathcal O_D(m+3)$, $Q^2=\wedge^2T(m)|_D=\mathcal O_D(m+3)$. Setting $z=0$ in the formula of Lemma~\ref{p1:lem:flat} gives $d_Q=-u(1-r)p$, and in invariant form $d_Q=-(1+m/3)\cdot(d\sigma|_D)$. Since $d\sigma|_D$ is an isomorphism, $Q$ is acyclic for $m\ne-3$, and $\HH^*(\mathcal O(m))=\mathbb H^*(C_{\log})$.
(4) $m=-3$: the terms $E_1^{a,b}=H^a(\wedge^bT(-3))$ are 1-dimensional at $(2,0),(1,1),(0,2)$ (Bott); all lie in total degree 2, so all further differentials vanish and the result is $[0,0,3,0,0]$. No sign check is needed.
\end{proof}
\begin{rem}[\tV]
The table above agrees with a direct computation, for $m=-15,\dots,15$ and the Fermat and Hesse cubics, of the differential $d_1$ of the spectral sequence $E_1^{a,b}=H^a(\wedge^bT(m))$ by exact linear algebra modulo the prime $32003$ (for $m\ge-2$ as a complex of global sections, for $m\le-4$ as a Cech computation with Laurent monomials) (scripts: \texttt{cp2/poisson.py, table.py, neg\_check.py, predict.py}).
\end{rem}
\begin{rem}[\tC\ Related work]
Eq.~(11), $0\to T^*\xrightarrow{\sigma}T\to T|_D\to0$, and Prop.~8 of Hitchin, \emph{Deformations of holomorphic Poisson manifolds}, arXiv:1105.4775 (Moscow Math.\ J.\ 12 (2012) 567--591) are related, but this result for $\mathbb{CP}^2$ itself is stated neither in that paper nor in Gualtieri's (as far as surveyed).
\end{rem}

\begin{obs}[How the dependence on $m$ appears]\tP
The total Euler characteristic does not depend on $m$ (a topological invariant). The dependence on $m$ appears only in the Poincar\'e polynomial, in particular in the period-3 jumps and in $\HH^0$ (the Casimir section $F^{m/3}$). Even for $m\gg0$, $\HH^2$ stays 2- or 3-dimensional, so Serre vanishing in total degree does not hold.
\end{obs}

%--------------------------------------------------------------
\section{Main theorem (surfaces): generalized Dolbeault cohomology of log-symplectic surfaces}\label{p1:sec:main2}
%--------------------------------------------------------------
\subsection{Statement}
Let $X$ be a compact complex surface, $\sigma\in H^0(-K_X)$ with reduced zero divisor $D=\bigcup D_i$ (irreducible components $D_i$, with only ordinary double points as singularities of $D$, their number being $\mu$), and $E$ a line bundle with a Poisson module structure. The residues $r_i$ (constant along each branch) are given by Theorem~\ref{p1:thm:main1}(c). (The branches at a double point are regarded as different branches.)
Let $I_*=\{i:r_i\in\Z_{\le0}\}$, $I_!=\{i:r_i\notin\Z_{\le0}\}$, $X'=X\setminus\bigcup_{i\in I_!}D_i$, $j_1:U\hookrightarrow X'$, $j_2:X'\hookrightarrow X$, and put
\[\mathcal F:=j_{2!}\,Rj_{1*}L_U\]
($Rj_*$ along the components in $I_*$, extension by zero along those in $I_!$).

\begin{thm}[Main theorem for surfaces]\tP\label{p1:thm:main2}
In the above setting assume all $r_i\neq1$. Then
\[\HH^k(E)\;\cong\;H^k(X;\mathcal F)\ \oplus\ \C^{\mu}\ (\text{only for }k=2).\]
In particular
\begin{enumerate}
\item[(i)] if all $r_i\notin\Z_{>0}$ then $\HH^*(E)=H^*(U;L_U)\oplus\C^\mu[-2]$;
\item[(ii)] if all $r_i\notin\Z_{\le0}$ then $\HH^*(E)=H^*_c(U;L_U)\oplus\C^\mu[-2]$;
\item[(iii)] the total Euler characteristic is $\sum(-1)^k\dim\HH^k(E)=\chi(U)+\mu=\chi(X)$ (independent of $E$).
\end{enumerate}
\end{thm}
\begin{hyptab}
Surface, $\sigma\in H^0(-K_X)$, $D$ reduced with only ordinary double points & To use the local normal forms $h=zu$ or $z_1z_2u$ and compute the quotient complex $Q$ explicitly. & Lemma~\ref{p1:lem:Q}, Prop.~\ref{p1:prop:moser}, Lemma~\ref{p1:lem:conn}.\\\midrule
$r_i\ne1$ (all branches) & If $r_i=1$ the differential $-u(1-r)p$ of $Q$ vanishes, the cohomology sheaves of $Q$ become infinite-dimensional, and the reduction breaks down. & Lemma~\ref{p1:lem:Q}. ($r_i=1$ is the special case $E=K_X$, Prop.~\ref{p1:prop:K}.)\\\midrule
Residues not positive integers / not non-positive integers & Condition for the comparison of the logarithmic complex with $Rj_*L_U$, $j_!L_U$ (Deligne) to be a quasi-isomorphism. & Lemma~\ref{p1:lem:stalk}.\\\midrule
$X$ compact & $R\Gamma(X;j_!L_U)=R\Gamma_c(U;L_U)$ and the Euler characteristic computation. & (ii), (iii).\\\midrule
$-K_X=\mathcal O(D)$ (anticanonical curve) & Gives $\chi(D)=\mu$ ($\omega_D=\mathcal O_D$, $p_a(D)=1$). & (iii).
\end{hyptab}

\subsection{Proof}
\begin{lem}[Stalks of the logarithmic complex and Deligne's comparison]\tP\ /\ \tC\label{p1:lem:stalk}
In a local frame (Theorem~\ref{p1:thm:main1}(d)), the stalk of $\mathrm{DR}_{\log}$ at a point of $D_I$ ($|I|=s$) is the tensor product of the Koszul complex of the commuting operators $z_i\partial_i+r_i\ (i\in I)$ with the ordinary de Rham complex in the remaining variables. On a monomial $z^a$, $z_i\partial_i+r_i=a_i+r_i$ (a scalar). Hence the stalk cohomology is $\wedge^*\C^s$ (dimension $2^s$, Euler characteristic 0) exactly when $-r_i\in\mathbb N$ (i.e.\ $r_i\in\Z_{\le0}$) for all $i\in I$, and $0$ otherwise. On the other hand the stalk of $Rj_*L_U$ at the same point is $\wedge^*\C^s$ exactly when all $e^{2\pi ir_i}=1$ ($r_i\in\Z$). Therefore:
\begin{itemize}
\item the natural map $\mathrm{DR}_{\log}\to Rj_*L_U$ is a quasi-isomorphism when no $r_i$ is a positive integer (Deligne, LNM 163, Ch.~II \S3, number \tQ; here confirmed by direct computation);
\item the extension by zero $j_!L_U\to\mathrm{DR}_{\log}$ (a zero extension of flat sections is holomorphic) is a quasi-isomorphism when no $r_i$ lies in $\Z_{\le0}$;
\item in the mixed case, over $X'$ we have $\mathrm{DR}_{\log}\simeq Rj_{1*}L_U$, and over $\bigcup_{I_!}D_i$ the stalks are $0$, so $\mathrm{DR}_{\log}\simeq\mathcal F$.
\end{itemize}
Division by the coefficients $1/(a_i+r_i)$ does not destroy convergence, since $|a_i+r_i|\ge c>0$ except for finitely many $a$.
\end{lem}
\begin{lem}[The quotient complex $Q$]\tP\ \tV\label{p1:lem:Q}
$Q:=C(A;E)/C(T_D;E)$ is supported on $D$ and lives only in degrees 1, 2: $Q^1=(T/T_D)\otimes E$, $Q^2=(\wedge^2T/\wedge^2T_D)\otimes E$.
\begin{itemize}
\item Smooth points ($h=zu$): $Q^1=\partial_z\otimes\mathcal O_D$, $Q^2=\partial_z\wedge\partial_w\otimes\mathcal O_D$, $d_Q(p)=-u(1-r)\,p|_D$. Hence acyclic when $r\ne1$.
\item Double points ($h=z_1z_2u$): $Q^1=\mathcal O_{D_1}\partial_1\oplus\mathcal O_{D_2}\partial_2$, $Q^2=\mathcal O_D$, $d_Q(p,q)=-u[(1-r_1)z_2p+(1-r_2)z_1q]\bmod z_1z_2$. If $r_1,r_2\ne1$ it is injective with image the maximal ideal $\mathfrak m_D$, so $H^1(Q)_p=0$ and $H^2(Q)_p=\mathcal O_D/\mathfrak m_D=\C$.
\end{itemize}
Hence for $r_i\ne1$, $Q\simeq\bigoplus_{\text{double points }p}\C_p[-2]$ in the derived category.
Derivation: restrict the formula of Lemma~\ref{p1:lem:flat} to $z=0$ and substitute $a|_D=\theta_2h$, $b|_D=-\theta_1u$, $\theta_1|_D=r$ (Theorem~\ref{p1:thm:main1}(d)). (\tV\ With random polynomials $u,g,p,q$ and residues, at smooth points and double points (20 trials each), $d_1d_0=0$ and the formula for $d_Q$ were confirmed in exact arithmetic. Script \texttt{cp2/ncd\_cochain\_check.py}.)
\end{lem}
\begin{prop}[Logarithmic Moser normal form at a double point]\tP\label{p1:prop:moser}
Near a double point $p$ let $\omega=\sigma^{-1}=u^{-1}\dlog z_1\wedge\dlog z_2$ and $c:=u(0)$. There is a biholomorphic map $\phi$ preserving each branch of $D$ and fixing $p$ such that $\phi^*\omega=c^{-1}\dlog z_1\wedge\dlog z_2$. The constant $c$ is a local invariant and cannot in general be made $1$.
\end{prop}
\begin{proof}
Let $\omega_0=\dlog z_1\wedge\dlog z_2$, $k=1/u-1/c$ ($k(0)=0$), $\omega_t=(c^{-1}+tk)\omega_0$. $\omega_t$ is closed and non-degenerate near $p$ for each $t\in[0,1]$. For $k=\sum k_{ab}z_1^az_2^b$,
\[\alpha=\sum_{a\ge1,b\ge0}\tfrac{k_{ab}}{a}z_1^az_2^b\,\dlog z_2-\sum_{b\ge1}\tfrac{k_{0b}}{b}z_2^b\,\dlog z_1\]
is a convergent logarithmic 1-form with $d\alpha=k\omega_0$ (using $d(z^{ab}\dlog z_2)=az^{ab}\omega_0$, $d(z_2^b\dlog z_1)=-bz_2^b\omega_0$). Since $X\mapsto\iota_X\omega_t$ is an isomorphism $T_D\to\Omega^1(\log D)$, there is a unique $X_t\in T_D$ with $\iota_{X_t}\omega_t=-\alpha$; it is tangent to each branch and vanishes at $p$. Its flow $\phi_t$ preserves $D$, and $\tfrac{d}{dt}\phi_t^*\omega_t=\phi_t^*(d\iota_{X_t}\omega_t+k\omega_0)=\phi_t^*(-d\alpha+k\omega_0)=0$. Hence $\phi_1^*\omega_1$ is of the form $\omega_{c^{-1}}$.\qedhere
\end{proof}
\begin{lem}[Vanishing of the connecting morphism]\tP\ \tV\label{p1:lem:conn}
For the exact triangle $C(T_D;E)\to C(A;E)\to Q\to C(T_D;E)[1]$, the $p$-component $e_p\in H^3_p(X;\mathrm{DR}_{\log})$ of the connecting morphism is $0$ when $(r_1,r_2)\ne(1,1)$.
\end{lem}
\begin{proof}[Sketch]
By Prop.~\ref{p1:prop:moser} and Lemma~\ref{p1:lem:res} we may assume $\sigma=c\,z_1z_2\partial_1\wedge\partial_2$ and $\nabla^{\log}e=(r_1\dlog z_1+r_2\dlog z_2)e$. All complexes are graded by the torus weight $\gamma\in\Z^2$, and the differentials preserve weights ($\lambda_i=\gamma_i+r_i$). Local cohomology can be computed with the Laurent tail $H^2_p(\mathcal O)=\mathrm{span}\{z^{-a}:a\ge1\}$. (a) $H^3_p(\mathrm{DR}_{\log})$ appears only in the weight with $\lambda_1=\lambda_2=0$ ($\gamma=-r$), and is non-zero only if $r_1,r_2$ are positive integers. (b) $H^2_p(Q)=\C$ sits in the single weight $\gamma=(-1,-1)$ (this requires an extra class $H^1_p(\mathcal O_D)\ni[c_1\ne c_2]$ in the local cohomology of $\mathcal O_D$), and $(-1,-1)=-r$ only if $r_1=r_2=1$. Since the connecting morphism preserves weight, $e_p=0$ for $(r_1,r_2)\ne(1,1)$. Together with the weight-by-weight Koszul computation there is no convergence problem (all maps are diagonal on monomials).
\tV\ Script \texttt{cp2/ncd\_local\_weights.py}: for weights $|\gamma_i|\le8$ and 78 pairs with $r_i\in\{-3..3,\pm\frac12,\frac13,-\frac23,\frac52\}$ (all $r_i\ne1$), the stalk cohomology $[0,0,1]$, the cone of local cohomology $=H^2_p$ (one class), and injectivity of $H(\mathrm{LOG})\to H(\mathrm{FULL})$ (connecting morphism $0$) were confirmed in exact arithmetic. 0 violations. For $r=(1,1)$ the connecting morphism is non-zero.
\end{proof}
\begin{proof}[Proof of Theorem~\ref{p1:thm:main2}]
By Theorem~\ref{p1:thm:main1} and Lemma~\ref{p1:lem:stalk}, $C(T_D;E)\simeq\mathrm{DR}_{\log}\simeq\mathcal F$. By Lemmas~\ref{p1:lem:Q} and \ref{p1:lem:conn}, $C(A;E)\simeq\mathrm{DR}_{\log}\oplus\bigoplus_p\C_p[-2]$ in $D^b(\C_X)$. Apply $R\Gamma(X;\cdot)$.
(i), (ii): $\mathcal F=Rj_*L_U$ or $j_!L_U$, and since $X$ is compact $R\Gamma(X;j_!L_U)=R\Gamma_c(U;L_U)$.
(iii): $\chi(X;\mathcal F)=\sum_S\chi_c(S)\chi(\mathcal F_x)$, and the stalk Euler characteristic on $D$ is $0$ ($\wedge^*\C^s$ or $0$), so this equals $\chi_c(U)=\chi(U)=\chi(X)-\chi(D)$. Since $D$ is anticanonical, $\omega_D=\mathcal O_D$, and $\chi(\mathcal O_D)=0$ gives $\sum_j(2-2\tilde g_j)=2\mu$ ($\tilde D_j$ the normalizations), so $\chi_{\rm top}(D)=\mu$ and $\chi(U)+\mu=\chi(X)$. (Alternative proof in the simple normal crossings case: $\chi(U)=\int c_2(T_X(-\log D))$.)
\end{proof}
\begin{prop}[$r_i=1$: $E=K_X$]\tP\label{p1:prop:K}
For $E=K_X$ we have $\wedge^bT\otimes K_X=\Omega^{2-b}$ and $E_1^{a,b}=H^{2-b,a}(X)$ (total degree $2+a-p$ with $p=2-b$). If $h^{p,q}(X)=0$ for $p\ne q$, then $E_1$ is concentrated in total degree $2$, the differentials (which raise total degree by 1) vanish, so $\HH^k=0$ ($k\ne2$) and $\dim\HH^2=\sum_ph^{p,p}=\chi(X)$ ($3$ for $\mathbb{CP}^2$, $4$ for $\mathbb{CP}^1\times\mathbb{CP}^1$).
\end{prop}

\subsection{Check on $\mathbb{CP}^1\times\mathbb{CP}^1$}
Let $\sigma=\Pi_F$, $F\in H^0(\mathcal O(2,2))$ ($-K=\mathcal O(2,2)$). \tP\ For smooth $D$, $\mathcal O(a,b)$ has a Poisson module structure only if $a=b=m$, the residue is $-m/2$, and the structure is unique (same argument as in \S\ref{p1:sec:cp2}). \tV\ Exact computation for $m=-9,\dots,8$ (script \texttt{cp2/p1p1.py}):
\begin{center}\small
\begin{tabular}{lll}\toprule
$D$ & $m$ even $\ge0$ / even $\le-4$ & $m$ odd, $m=-2$\\\midrule
smooth $(2,2)$ curve & $[1,0,3,0,0]$ / $[0,0,3,0,1]$ & $[0,0,4,0,0]$\\
two transverse $(1,1)$ curves ($\mu=2$) & $[1,1,4,0,0]$ / $[0,0,4,1,1]$ & $[0,0,4,0,0]$\\
four lines $x_0x_1y_0y_1$ ($\mu=4$) & $[1,2,5,0,0]$ / $[0,0,5,2,1]$ & $[0,0,4,0,0]$\\
$(1,1)+(1,0)+(0,1)$ ($\mu=3$) & $[1,1,4,0,0]$ / $[0,0,4,1,1]$ & $[0,0,4,0,0]$\\\bottomrule
\end{tabular}
\end{center}
Period 2, Serre duality $m\leftrightarrow-m-4$, and $\chi=4$ always hold. For reducible $D$ the Poisson module structure is not unique ($H^0(\Omega^1(\log D))\ne0$), and the table is computed for the symmetric structure with equal residues on all components (Lemma~\ref{p1:lem:res}).

\subsection{Verification of Theorem~\ref{p1:thm:main2}}
\tV\ Theorem~\ref{p1:thm:main2} and its closed formula (Theorem~\ref{p1:thm:B} in \S\ref{p1:sec:RR}) were compared with the exact computations above for a total of 196 values of $m$ in the following eight examples (script \texttt{cp2/ncd\_global\_check.py}): on $\mathbb{CP}^2$ a smooth cubic, a nodal cubic ($\mu=1$), a conic plus a secant line ($\mu=2$), a triangle ($\mu=3$) ($m=-15..15$), and the four $\mathbb{CP}^1\times\mathbb{CP}^1$ examples above ($m=-9..8$). 0 mismatches. The cuspidal cubic, conic plus tangent line (tacnode), three concurrent lines and non-reduced $D$ on $\mathbb{CP}^2$ are outside the scope of this theorem and are treated in \S\ref{p1:sec:local}.

%--------------------------------------------------------------
\section{Serre-type theorems}\label{p1:sec:serre}
%--------------------------------------------------------------
\subsection{Correspondence with the classical case}
The generalized counterparts of the ingredients of the classical Serre vanishing theorem (Theorem~\ref{p1:thm:Serre}) are as follows (\tK\ the correspondence is the author's organization).
\begin{center}\small
\begin{tabular}{p{4.6cm}p{9.8cm}}\toprule
Classical & Counterpart in this part\\\midrule
Coherent sheaf $\mathcal F$ & Each term $K^b=\wedge^bT\otimes E$ of the Lie algebroid complex ($E$ a generalized holomorphic bundle). It is coherent.\\
Ample $\mathcal O(1)$ & An ample holomorphic line bundle $L_0$ which is also a Poisson module (generalized holomorphic). The Poisson structure is needed only so that $C(A;\cdot)$ is a complex; it plays no role in positivity.\\
Projective & $X$ projective (because $L_0$ exists).\\
$H^q(\mathcal F(m))=0\ (q>0)$ & $E_1^{a,b}=H^a(X,\wedge^bT\otimes E\otimes L_0^m)=0$ (\emph{$a>0$: the transversal (Dolbeault) degree}).\\
Conclusion: $H^*(\mathcal F(m))$ is concentrated in degree 0 & $\HH^*$ is the cohomology of the finite-dimensional complex of global sections $(\Gamma(\wedge^\bullet T\otimes E\otimes L_0^m),d)$.\\\bottomrule
\end{tabular}
\end{center}

\subsection{Serre-type theorem in any dimension}
\begin{thm}[Serre-type theorem ($\beta$-transforms)]\tP\label{p1:thm:A1}
Let $X$ be a compact complex projective manifold ($\dim_\C X=n$), $\sigma$ any holomorphic Poisson structure ($[\sigma,\sigma]=0$; neither non-degeneracy nor a logarithmic condition is needed), $E$ a generalized holomorphic vector bundle (Poisson module), and $L_0$ an ample holomorphic line bundle with a Poisson module structure. Then there is $m_0$ depending only on $X,\sigma,E,L_0$ such that for $m\ge m_0$:
\begin{enumerate}
\item[(a)] $H^a(X,\wedge^bT\otimes E\otimes L_0^m)=0$ for all $a>0$ and all $b$;
\item[(b)] $\HH^k(E\otimes L_0^m)\cong H^k\bigl(\Gamma(X,\wedge^\bullet T\otimes E\otimes L_0^m),\,d\bigr)$. In particular $\HH^k=0$ for $k>n$;
\item[(c)] $\HH^0(E\otimes L_0^m)=$ the space of all flat (Casimir) holomorphic sections.
\end{enumerate}
\end{thm}
\begin{hyptab}
$X$ projective, $L_0$ ample & The hypotheses of classical Serre themselves. & (a): apply Hartshorne III.5.2 to the finitely many ($b=0,\dots,n$) coherent sheaves $\wedge^bT\otimes E$.\\\midrule
$L_0$ a Poisson module (and $E$) & $L_0^m$ gets the tensor-product module structure, so $C(A;E\otimes L_0^m)$ is a complex. & Definition of the complex $K^\bullet=C(A;E\otimes L_0^m)$. Unrelated to positivity.\\\midrule
$[\sigma,\sigma]=0$ & $d^2=0$. & Same.\\\midrule
(not needed) non-degeneracy of $\sigma$, normal crossings & --- & Not used in this theorem.
\end{hyptab}
\begin{proof}
$L_0^m$ has the tensor-product module structure (Leibniz rule). Applying Serre's theorem (Theorem~\ref{p1:thm:Serre}, Hartshorne III.5.2) to the $n+1$ coherent sheaves $\wedge^bT\otimes E$ and taking $m_0$ the maximum of the thresholds gives (a). Each $K^b$ is $\Gamma$-acyclic, so for the bounded complex $K^\bullet$ (whose differential is $\C$-linear but not $\mathcal O$-linear) the spectral sequence $E_1^{a,b}=H^a(K^b)$ is concentrated at $a=0$ and $R\Gamma(K^\bullet)=\Gamma(K^\bullet)$ (acyclic resolution lemma). This is (b). Since $b\le n$, it vanishes for $k>n$. (c) is $H^0$ of the $\Gamma$-complex: $\HH^0=\ker(d_A:H^0(E)\to H^0(T\otimes E))$, i.e.\ flat sections.\qedhere
\end{proof}
\begin{rem}[Vanishing in total degree is false]\tP
For $\sigma=0$ (a complex structure), $\HH^k(L)=\bigoplus_{a+b=k}H^a(X,\wedge^bT\otimes L)$, so even for $m\gg0$ we have $H^0(\wedge^bT\otimes L^m)\ne0$, hence $\HH^b\ne0\ (b=0,\dots,n)$. For $m=0$ on $\mathbb{CP}^2$ this gives $[1,8,10,0,0]$ (\tV). Only the transversal degree $a>0$ and total degrees $>n$ vanish; ``Serre vanishing in total degree'' does not hold.
\end{rem}
\begin{rem}[\tV\ Effective thresholds]
By Bott's formula, the condition for $E_1^{a>0}=0$ is $m\ge-2$ on $\mathbb{CP}^2$ (for $m=-3$ the three terms $(2,0),(1,1),(0,2)$ are all non-zero), and $m\ge-1$ on $\mathbb{CP}^1\times\mathbb{CP}^1$ with $L_0=\mathcal O(1,1)$ ($m=-2$ is the exception $K_X$). Both agree with the range in which degrees $\ge3$ vanish in the computed tables.
\end{rem}

\subsection{Periodic stabilization}
\begin{thm}[Periodic stabilization (log-symplectic surfaces)]\tP\ (depends on Theorem~\ref{p1:thm:main2} and its hypotheses)\label{p1:thm:A2}
Let $X$ be a simply connected projective surface, $\sigma$ log-symplectic along $D$ (anticanonical, only ordinary double points), $L_0$ a line bundle with a Poisson module structure whose residues are $\rho_i=-a_i$ ($a_i\in\mathbb Q_{>0}$, constant along each component), and $U=X\setminus D$ affine ($D$ ample). Let $N$ be the least common multiple of the denominators of the $a_i$. If $m\ge1$ then
\[\HH^k(L_0^m)=0\ (k\ge3),\qquad\HH^*(L_0^m)=H^*(U;L_U^{(m)})\oplus\C^\mu[-2],\]
where $L_U^{(m)}$ is the local system with monodromy $e^{-2\pi ima_i}$ around $D_i$. In particular $\HH^*(L_0^m)$ depends only on $m\bmod N$.
\end{thm}
\begin{hyptab}
Residues $-ma_i<0$ & $r_i\ne1$ and not a positive integer. & Theorem~\ref{p1:thm:main2}(i): $\mathcal F=Rj_*L_U$.\\\midrule
$U$ affine ($D$ ample) & Artin vanishing $H^k(U;L)=0\ (k>2)$. & Vanishing in degrees $\ge3$. (Theorem~\ref{p1:thm:A1}(b) gives the same conclusion from projectivity and ampleness and makes affineness unnecessary.)\\\midrule
$X$ simply connected & The abelianization of $\pi_1(U)$ is generated by meridians. & The local system is determined by $m\bmod N$.\\\midrule
$a_i\in\mathbb Q_{>0}$ & $c_1(L_0)$ is a positive rational combination of the boundary components (``positive with respect to the boundary''). & Periodicity.
\end{hyptab}
\tV\ Period 3 for $\mathbb{CP}^2$ ($L_0=\mathcal O(1)$, $a=1/3$), period 2 for $\mathbb{CP}^1\times\mathbb{CP}^1$ ($L_0=\mathcal O(1,1)$, $a=1/2$), for $m\ge1$. Confirmed for all eight computed examples.

\begin{conj}[General case]\tK\label{p1:conj:A3}
For a general generalized complex manifold $M$ and a generalized holomorphic line bundle $L_0$ that is ``positive'' in a suitable sense (candidate: ample on the leaf space of the regular locus, and with positive residues along the locus where the type jumps), $\HH^*(L_0^m)$ vanishes in total degrees $>n_{\max}$ for $m\gg0$ and is periodic in $m$. For a general $\Jc$ that is neither a $\beta$-transform nor of regular type, there is no way to define the ``$E_1$ term of a coherent filtration $K^b$'', and this is the main obstruction to going beyond Theorem~\ref{p1:thm:A1}.
\end{conj}

\subsection{Relation to the theorem of Pal--Poddar}
\begin{prop}[Consistency in the product case]\tP\label{p1:prop:PPprod}
Let $M=Y\times X$ ($Y$ a compact symplectic manifold, $X$ a projective complex manifold) and $E=\mathrm{pr}_Y^*F\boxtimes\mathrm{pr}_X^*G_0^m$ ($F$ a flat bundle on $Y$, $G_0$ ample). Then the generalized Dolbeault complex is the tensor product of the de Rham complex of $Y$ and a complex on $X$, and
\[\HH^k=\bigoplus_{j+b+a=k}H^j_{dR}(Y;F)\otimes H^a(X,\wedge^bT_X\otimes G_0^m)\]
(K\"unneth). By classical Serre the terms with $a>0$ vanish for $m\gg0$, and the remaining $H^j_{dR}(Y;F)\otimes H^0(X,\wedge^bT_X\otimes G_0^m)$ are in general non-zero up to total degree $\dim Y+\dim_\C X$. Hence Serre-type vanishing for $\mathrm{pr}_X^*G_0^m$ holds only for the transversal degree $a$.
\end{prop}
\src{Pal--Poddar, \emph{Strong generalized holomorphic principal bundles}, arXiv:2404.18113, \S11, Thm.~11.12: if $M$ is a compact generalized complex manifold of regular type $k$ with $M/\mathcal S$ a K\"ahler orbifold and $E$ is a positive SGH line bundle, then (1) $H^q(M,(\mathcal G^*M)^p\otimes E)=0$ for $p+q>k$, and (2) for every SGH vector bundle $E'$ there is $m_0$ such that $H^q(M,E'\otimes E^m)=0$ for $m\ge m_0$, $q>0$.}
\begin{rem}[\tQ]
The groups $H^q(M,(\mathcal G^*M)^p\otimes E)$ of Pal--Poddar (which by Cor.~8.10 are $H^{p,q}_{d_L}(M,E)$) are sheaf cohomology groups of transversal forms that are constant along the leaves, and $q$ is the transversal (complex) degree (confirmed by re-reading \S\S5--11, 13; the extraction went through summaries, so details of formulas are unverified). Hence the reading of the proposition above (vanishing only in the transversal degree $a$) is consistent with Thm.~11.12. The groups are not the same, however: Pal--Poddar use exterior powers of $\mathcal G^*M$, this part tensors with $\wedge^bT$ (tangent side), and moreover the $\HH$ of this part is the total complex including the leafwise de Rham factor. The proposition above is a computation in the product case and is not itself a verification of Thm.~11.12. Theorem~\ref{p1:thm:A1} does not assume regular type or an orbifold leaf space, and includes type-jumping $\beta$-transforms (e.g.\ types $0$ and $n$ coexisting), at the price of requiring the underlying complex manifold to be projective (in the $\beta$-transform case only).
\end{rem}

%--------------------------------------------------------------
\section{Duality}\label{p1:sec:duality}
%--------------------------------------------------------------
\begin{thm}[Serre duality for generalized holomorphic line bundles]\tP\label{p1:thm:dual}
Let $X$ be any compact complex manifold ($\dim_\C X=n$), $\sigma$ any holomorphic Poisson structure (no K\"ahler, projective, log-symplectic or non-degeneracy assumption is needed; $\sigma=0$ is included), and $L$ a line bundle that is a Poisson module. Then for all $k$
\[\HH^k(L)^*\;\cong\;\HH^{2n-k}\bigl(L^\vee\otimes K_X^{2}\bigr).\]
$L^\vee\otimes K_X^2$ carries the tensor-product module structure (frame $e'$ with $ee'=\omega^2$ and $V'=-V+2V_K=-V-2M$; the canonical module structure on $K_X$ is $V_K^k=\sum_i\partial_i\sigma^{ki}$, and $M^j=\sum_i\partial_i\sigma^{ij}$ is the modular vector field with respect to $\omega=dz_1\wedge\cdots\wedge dz_n$).
\end{thm}
\begin{hyptab}
$X$ compact & Closed range and finite-dimensionality of elliptic complexes; no boundary terms in Stokes. & Step (4) (identification of the dual).\\\midrule
$[\sigma,\sigma]=0$, $L$ a Poisson module & $D^2=0$, $D'^2=0$ ($L'$ is also a module). & To treat them as complexes.\\\midrule
(not needed) K\"ahler / non-degenerate / logarithmic & Ellipticity is checked directly from the principal symbol, and the transposition is an identity of local operators, so no global primitive is needed. & Steps (1), (3).
\end{hyptab}
\begin{proof}[Proof (sketch)]
(0) Dolbeault model: $\mathrm{Tot}^k(L)=\bigoplus_{a+b=k}A^{0,a}(\wedge^bT\otimes L)$, $D=\dbar+d_K$ ($d_K(\phi\otimes Pe)=(-1)^a\phi\otimes d(Pe)$, $d=[\sigma,\cdot]+V\wedge$). $D^2=0$ follows from $\dbar^2=0$, $d^2=0$, $\dbar d_K+d_K\dbar=0$ (Prop.~\ref{p1:prop:ident}(b)).
(1) Ellipticity (\tP): at a real cotangent vector $\xi\ne0$, the principal symbol of $D$ is the exterior product with $v(\xi)=(\xi^{0,1},\sigma^\sh\xi^{1,0})\in L^*=T^{*0,1}\oplus T^{1,0}$, and $(\wedge L^*,v\wedge)$ is exact if and only if $v\ne0$. If $v=0$ then $\xi^{0,1}=0$, and since $\xi$ is real, $\xi=0$. Hence the complex is elliptic at every point, including points where $\sigma$ degenerates (this is a direct verification, in the present case, of Gualtieri's Prop.~\ref{p1:prop:ellip}). The same holds for $L'$.
(2) Perfect pairing: integrating $E^b(L)\times E^{n-b}(L')\to\wedge^nT\otimes K^2=K$ and $A^{0,a}\times A^{0,n-a}\to A^{0,n}(K)=A^{n,n}$ gives $\mathrm{Tot}^k(L)\times\mathrm{Tot}^{2n-k}(L')\to\C$ (a perfect pairing at each point).
(3) Transposition: from the following two identities, the formal transpose $D^t$ equals $-(-1)^{\deg}D'$.
\end{proof}
\begin{lem}[Lemma D: top-degree part of the Schouten bracket]\tP\ \tV\label{p1:lem:D}
For $T=\sum_j(-1)^ju_j\,\partial_1\wedge\cdots\widehat{\partial_j}\cdots\wedge\partial_n\in\wedge^{n-1}T$ we have $[\sigma,T]_{\rm top}=-\sum_{i,j}\partial_i(\sigma^{ij}u_j)+2\sum_jM^ju_j$ and $(M\wedge T)_{\rm top}=\sum_jM^ju_j$.
\end{lem}
\begin{lem}[Lemma E: transposition identity]\tP\ \tV\label{p1:lem:E}
For $P\in\wedge^bT$, $Q\in\wedge^{n-b-1}T$ (smooth coefficients), $d=[\sigma,\cdot]+V\wedge$, $d'=[\sigma,\cdot]+V'\wedge$, $V'=-V-2M$, and $\langle P,Q\rangle:=(P\wedge Q)_{\rm top}$, we have
\[\langle dP,Q\rangle+(-1)^b\langle P,d'Q\rangle=-\sum_i\partial_i\Bigl(\sum_j\sigma^{ij}u_j\Bigr),\quad T=P\wedge Q=\sum_j(-1)^ju_j\partial_{\hat j}.\]
Proof: the Leibniz rule $[\sigma,P\wedge Q]=[\sigma,P]\wedge Q+(-1)^bP\wedge[\sigma,Q]$ ($\sigma$ has even degree); the $V$-terms give $V\wedge P\wedge Q+(-1)^bP\wedge(-V)\wedge Q=0$, and the $M$-term is $-2M\wedge T$. The result follows from Lemma~D.
\tV\ For $n=2,3,4$, with random polynomial and Laurent monomial data and arbitrary bivector fields $\sigma$ (the identity itself holds even if $\sigma$ is not Poisson), the identity was confirmed in exact arithmetic (\texttt{cp2/duality\_ndim\_exact.py}). The constant $c=-2$ and the sign $(-1)^b$ were found by solving a system of equations independently of the derivation (\texttt{cp2/duality\_ndim.py}). $[\sigma,M]=0$ and $d^2=0$ ($V=M$, $\sigma$ Poisson) were also confirmed.
\end{lem}
\begin{proof}[Continuation of the proof of Theorem~\ref{p1:thm:dual}]
(3) continued: the right-hand side of Lemma~E is an exact form of the type $\partial_iF^i\,\omega\wedge d\bar z^{1\cdots n}=\pm d(F^i\iota_{\partial_i}\omega\wedge d\bar z^{1\cdots n})$ (the $\dbar$ component vanishes against $d\bar z^{1\cdots n}$), so by Stokes its integral is $0$. The $\dbar$ part is the classical Serre duality computation: $\langle\dbar\alpha,\beta\rangle=-(-1)^{a+b}\langle\alpha,\dbar\beta\rangle$. The $d$ part is $\langle d_K\alpha,\beta\rangle=-(-1)^{a+b}\langle\alpha,d'_K\beta\rangle$. That the two signs agree was checked by hand (not machine-checked \tQ). Since $D^t$ is a local operator, the equality in coordinate charts gives globally $\langle D\alpha,\beta\rangle=-(-1)^{\deg\alpha}\langle\alpha,D'\beta\rangle$.
(4) Conclusion: $D,D'$ are elliptic complexes on a closed manifold, transposes of each other (up to sign) with respect to the perfect pairing. By Hodge theory for elliptic complexes, $\HH^k$ is finite-dimensional with closed range, $\HH^k(L)^*$ is the cohomology of the transposed complex (with distributional coefficients), and by elliptic regularity it equals the one with smooth coefficients, $\HH^{2n-k}(L')$.\qedhere
\end{proof}
\begin{rem}[Consistency checks\ \tV]
(i) $\sigma=0$: classical Serre duality (using $\wedge^bT^*=\wedge^{n-b}T\otimes K$). (ii) $\mathbb{CP}^2$, $K^2=\mathcal O(-6)$: $m\leftrightarrow-m-6$; $\mathbb{CP}^1\times\mathbb{CP}^1$, $K^2=\mathcal O(-4,-4)$: $m\leftrightarrow-m-4$ (consistent with the tables in \S\ref{p1:sec:cp2} and \S\ref{p1:sec:main2}). (iii) In the log-symplectic case the residue of $K_X$ is $1$ ($r=1$ in Prop.~\ref{p1:prop:K}), so the residue of $L'$ is $2-r_i$; this is the origin of the symmetry $r\leftrightarrow2-r$ in the formula of the main theorem (in any even dimension).
\end{rem}

%--------------------------------------------------------------
\section{Riemann--Roch-type formulas}\label{p1:sec:RR}
%--------------------------------------------------------------
\subsection{The total Euler characteristic is a topological invariant}
\begin{prop}\tP\label{p1:prop:euler}
Let $X$ be a compact complex manifold, $\sigma$ any holomorphic Poisson structure, and $E$ a Poisson module of rank $\rk E$. Then
\[\sum_k(-1)^k\dim\HH^k(E)=\sum_b(-1)^b\chi(X,\wedge^bT\otimes E)=(-1)^n\,\rk E\cdot\chi(X).\]
\end{prop}
\begin{proof}
The complex $C(A;E)=(\wedge^bT\otimes E,d)$ has a filtration by $b$ (the stupid filtration), and the spectral sequence $E_1^{a,b}=H^a(X,\wedge^bT\otimes E)$ is finite-dimensional and converges to $\HH^{a+b}$. The Euler characteristic equals that of the $E_1$ term, and $\sum_b(-1)^b\chi(\wedge^bT\otimes E)$ is, by Hirzebruch--Riemann--Roch (Theorem~\ref{p1:thm:HRR}),
$\int\ch(E)\Td(T)\sum_b(-1)^b\ch(\wedge^bT)$. Since $\Td(T)=\prod\frac{x_i}{1-e^{-x_i}}$ and $\sum_b(-1)^b\ch(\wedge^bT)=\prod(1-e^{x_i})$, the product is $\prod(-x_ie^{x_i})=(-1)^ne^{c_1}c_n(T)$. As $c_n$ is of top degree, only the rank of $\ch(E)$ contributes, giving $(-1)^n\rk E\int c_n(T)=(-1)^n\rk E\,\chi(X)$.\qedhere
\end{proof}
\begin{rem}[\tP\ Meaning]
For $\mathbb{P}^1$ and $E=\mathcal O(m)$ the formula reads $\chi(\mathcal O(m))-\chi(\mathcal O(m+2))=(m+1)-(m+3)=-2$ (independent of $m$). That is, the Riemann--Roch formula for the \emph{total} Euler characteristic of the generalized Dolbeault complex is a topological formula that does not see $c_1(E)$, and is not a generalization of the classical $\chi(M,E)=\int\ch(E)\Td(M)$. The classical $\chi(M,E)$ corresponds to the Euler characteristic of the ``column'' $b=0$, $H^\bullet(X,E)$. For a general $\Jc$ this column is not preserved by $d$, and some refinement is needed.
\end{rem}

\subsection{Closed formula for the Poincar\'e polynomial}
\begin{thm}[Riemann--Roch-type closed formula (surfaces)]\tP\ (depends on Theorem~\ref{p1:thm:main2} and the hypotheses below)\label{p1:thm:B}
In the setting of Theorem~\ref{p1:thm:main2} assume in addition: (H1) $H^1(X;\C)=0$ (for $b_1(U)=c-\rho$); (H2) $U$ is affine (for vanishing in degrees $\ge3$); (H3) all components have the same residue $r_i=r\ne1$; (H4) if $r\notin\Z$, then $H^0(U;L_U)=H^1(U;L_U)=0$ (G2: see the remark below). Put $b_1=c-\rho$ ($c$ the number of components, $\rho$ the dimension of the span of $[D_i]\in H^2(X;\C)$), $\chi_U=\chi(X)-\mu$, $b_2=\chi_U-1+b_1$. Then
\begin{align*}
P_E(t)&=1+b_1t+(b_2+\mu)t^2 &&(r\in\Z_{\le0}),\\
P_E(t)&=(b_2+\mu)t^2+b_1t^3+t^4 &&(r\in\Z_{\ge2}),\\
P_E(t)&=\chi(X)\,t^2 &&(r\notin\Z).
\end{align*}
Equivalently: $P_E(t)=\chi(X)t^2+[r\le0]\,\Delta(t)+[r\ge2]\,t^4\Delta(1/t)$, $\Delta(t)=(1+t)(1+(b_1-1)t)$.
\end{thm}
\begin{hyptab}
(H1) $H^1(X;\C)=0$ & From the Gysin sequence $H^0(D)\to H^2(X)$, $H^1(U)=\ker(\C^c\to H^2(X))$. & $b_1(U)=c-\rho$.\\\midrule
(H2) $U$ affine & Artin vanishing. & $b_3=b_4=0$, hence $b_2=\chi_U-1+b_1$.\\\midrule
(H3) equal residues & $\mathcal F$ is either $Rj_*$ or $j_!$. & Theorem~\ref{p1:thm:main2}(i), (ii).\\\midrule
(H4) $H^0(U;L)=H^1(U;L)=0$ for $r\notin\Z$ & So that $H^*(U;L)=[0,0,\chi_U]$. & The row $r\notin\Z$. For nodal curves on $\mathbb{CP}^2$, $\pi_1(U)$ is abelian (Fulton, Ann.\ Math.\ 111 (1980) 407--409; Deligne, S\'em.\ Bourbaki 543). \tC\\\midrule
$r_i\ne1$ & Hypothesis of Theorem~\ref{p1:thm:main2}. & Lemma~\ref{p1:lem:Q}.
\end{hyptab}
\begin{proof}
Substitute into Theorem~\ref{p1:thm:main2}(i), (ii): $H^*(U)=[1,b_1,b_2,0,0]$ ($r\in\Z_{\le0}$, trivial local system), $H_c^*(U)$ (by Poincar\'e duality $[0,0,b_2,b_1,1]$, $r\in\Z_{\ge2}$), $H^*(U;L)=[0,0,\chi_U,0,0]$ ($r\notin\Z$: from (H2), (H4) and $\chi(U;L)=\chi(U)$), and add $\mu t^2$. The equivalent form uses $\chi(X)=\chi_U+\mu=1-b_1+b_2+\mu$.
\end{proof}
\begin{rem}[Properties\ \tP]
(i) $P_E(-1)=\chi(X)$ for all $r$ (consistent with Prop.~\ref{p1:prop:euler}, $n=2$). (ii) The duality $P_{2-r}(t)=t^4P_r(1/t)$ is evident from the formula (a formula-level consequence of Theorem~\ref{p1:thm:dual}). (iii) $\mu$ and $b_2$ are absorbed into $\chi(X)t^2$, and the jump term $\Delta$ depends only on $b_1$ (the number of linear relations among the $[D_i]$). (iv) For $E=L_0^m$, the quasi-polynomial in $m$ that classical Riemann--Roch would give is the constant $\chi(X)$, and the dependence on $m$ appears only through the integrality (period $N$) and the sign of the residue $r=-ma$.
\end{rem}
\begin{rem}[\tQ\ On G2]
(H4) was checked in the eight verified examples (four curves on $\mathbb{CP}^2$ and four on $\mathbb{CP}^1\times\mathbb{CP}^1$). For the smooth $(2,2)$ curve and the two $(1,1)$ curves on $\mathbb{CP}^1\times\mathbb{CP}^1$ we cited Timmerscheidt (\emph{Mixed Hodge theory for unitary local systems}, J.\ reine angew.\ Math.\ 379 (1987) 152--171) and Arapura (\emph{Geometry of cohomology support loci for local systems I}, arXiv:alg-geom/9612006, J.\ Algebraic Geom.\ 6 (1997) 563--597 (volume/pages from memory \tQ), the torsion-translate theorem) (applicability not fully checked \tQ). In general $H^1(U;L_\chi)=0\Leftrightarrow\chi\notin V^1(U)$ (the first characteristic variety), which can fail for non-nodal $D$ (e.g.\ three concurrent lines) (Libgober, Arapura).
\end{rem}
\tV\ The closed formula agrees with the exact computations for a total of 196 values: $\mathbb{CP}^2$ (smooth, nodal, conic plus secant, triangle; $m=-15..15$) and $\mathbb{CP}^1\times\mathbb{CP}^1$ (4 examples, $m=-9..8$) (0 mismatches, \texttt{cp2/ncd\_global\_check.py}). The formula-level symmetry $r\leftrightarrow2-r$ was checked on a grid of 73 rational values of $r$ (\texttt{cp2/serre\_rr\_check.py}).

\subsection{Index-theoretic outlook}\label{p1:sec:index}
The content of this subsection is conjectural; no index theorem is proved.
\begin{obs}[\tK\ Indicial-root picture (qualitative)]
At smooth points the logarithmic complex acts on the monomial $z^ne$ by the factor $(n-r)$ (Lemma~\ref{p1:lem:stalk}). An indicial root exists exactly when $r\in\Z$; for $r\notin\Z$ the weight-graded complexes are acyclic and $P_r=\chi(X)t^2$. For $r\in\Z$, what remains is a subquotient of the cohomology of the boundary $\partial U$ of a tubular neighbourhood of $D$ (via the Gysin sequence, $H^*(D)\otimes(1\oplus t)$), and since $\chi(\partial U)=0$ it is invisible in the Euler characteristic. This is qualitatively consistent with Melrose's APS / b-calculus picture (Melrose, \emph{The Atiyah--Patodi--Singer Index Theorem}, A K Peters 1993).
\end{obs}
\begin{conj}[Index-type Riemann--Roch (log-symplectic surfaces)]\tK\label{p1:conj:I3}
(a) $\sum(-1)^k\dim\HH^k(L)=\int_Xc_2(T_X(-\log D))+\mu=\chi(X)$ is given by an Atiyah--Singer-type integral without boundary term. (b) The jump term $J(r)=P_r(t)-\chi(X)t^2$ is non-zero only for $r\in\Z\setminus\{1\}$, equal to $\Delta(t)$ for $r\le0$ and to $t^4\Delta(1/t)$ for $r\ge2$. $\Delta(t)$ is the Poincar\'e polynomial of a certain subquotient of the kernel $H^*(\partial U)$ of the indicial roots (fixed by the weight filtration of the long exact sequence $H_c^*(U)\to H^*(U)\to H^*(\partial U)$). (c) In general dimension and for non-normal-crossing $D$, the jump term contains local terms governed by the spectral numbers of the singularities (\S\ref{p1:sec:local}).
\end{conj}
\begin{rem}[\tC\ State of the literature and obstacles]
\begin{itemize}
\item \emph{Analogy with the $b$-index}: Braverman--Loizides--Song, \emph{Geometric quantization of $b$-symplectic manifolds}, arXiv:1910.10016 shows that the quantization index of a $b$-symplectic manifold is given by $\int\mathrm{Td}({}^bTM)\mathrm{Ch}(L)$ and that the boundary term vanishes in the limit (title and authors checked \tC). This is analogous to the fact that the Euler characteristic here is of the form $\chi(X)=\int c_2(T_X(-\log D))+\mu$, independent of $r$.
\item \emph{Index on Lie groupoids}: Pflaum--Posthuma--Tang, \emph{The index of geometric operators on Lie groupoids}, arXiv:1308.0236 (Indag.\ Math.\ 25 (2014) 1--18); Qiao--So, arXiv:2304.13345 (boundary groupoids; no eta term in codimension $\ge2$); Ammann--Lauter--Nistor, arXiv:math/0304044; Mazzeo--Melrose, arXiv:math/9812120; Hausel--Hunsicker--Mazzeo, arXiv:math/0207169 (Duke Math.\ J.\ 122 (2004), pages \tQ). Titles and authors were checked; applicability to the present problem was not.
\item \emph{Poisson cohomology}: Lanius, arXiv:1605.03854 ($C^\infty$ Poisson cohomology of log-symplectic manifolds). This is a different object from the holomorphic Poisson hypercohomology of this part; there is only an analogy in the structure ``$b$-part $+$ boundary term''.
\item \emph{What was not found}: an index theorem for the complex Lie algebroid $L=T^{0,1}\bowtie A^{1,0}$, or a $b$-Dolbeault index with the residue as a weight. The real Lie algebroid $T_X(-\log D)$ is integrable for smooth $D$ by Gualtieri--Li, arXiv:1206.3674. \tC\ The normal crossing case was not checked.
\item \emph{Unimodularity}: the Pflaum--Posthuma--Tang trace requires an invariant density, which for Poisson structures corresponds to the vanishing of the modular class. In the log-symplectic case the modular vector field is expected to vanish thanks to the logarithmic volume form, but this was not checked \tQ.
\end{itemize}
\end{rem}

%--------------------------------------------------------------
\section{The toric case and higher dimensions}\label{p1:sec:toric}
%--------------------------------------------------------------
\subsection{Toric log-symplectic structures}
Let $X$ be a smooth projective toric variety ($N=\Z^n$, $n=2k$) with rays $u_i$ ($i\in I$) and $D$ the toric boundary (simple normal crossings). With respect to the logarithmic frame $E_j=z_j\partial_j$, a constant non-degenerate skew-symmetric matrix $A$ gives $\sigma=\tfrac12\sum A_{ij}E_i\wedge E_j$ (torus-invariant, log-symplectic along $D$). Give $L=\mathcal O(\sum a_iD_i)$ the flat logarithmic connection $d+\rho$ ($\rho\in M_\C$). The residues are $r_i=\langle\rho,u_i\rangle-a_i$, and $c_1(L)=-\sum r_i[D_i]$. \tP\ Every module is of this form ($H^0(\Omega^1(\log D))=M_\C$, Lemma~\ref{p1:lem:res}).

\begin{thm}[Weight formula]\tP\ \tV\label{p1:thm:weight}
Assume $E_1^{a,b}=H^a(X,\wedge^bT\otimes L)=0\ (a>0)$ (e.g.\ in the range of Theorem~\ref{p1:thm:A1}). For $\chi\in M$ put $c_i(\chi)=\langle\chi,u_i\rangle+a_i$, $C(\chi)=\{i:c_i(\chi)=-1\}$, $W(\chi)=\mathrm{span}(u_i:i\in C(\chi))$, $s=\dim W$. Then
\[P_{\HH}(t)=\sum_{\chi\in\Lambda}t^{s(\chi)}(1+t)^{n-s(\chi)},\quad\Lambda=\{\chi\in M:c_i(\chi)\ge-1\ \forall i,\ A(\chi+\rho)\in W(\chi)\}.\]
\end{thm}
\begin{hyptab}
Toric, torus-invariant $\sigma$, constant $A$ & The complex of global sections decomposes by torus characters $\chi\in M$ ($T$-weights), and each weight is a finite-dimensional linear algebra problem. & Proof: determination of the weight-$\chi$ subcomplex $G^b_\chi$.\\\midrule
$E_1^{a>0}=0$ & $\HH^*$ equals the cohomology of the complex of global sections. & Theorem~\ref{p1:thm:A1}(b).\\\midrule
Non-degenerate $A$ & $\sigma^\sh:B\cong T_D$ holds, and $\wedge^b\sigma^\sh$ sends the logarithmic de Rham differential $(\chi+\rho)\wedge$ to $v_\chi\wedge$ ($v_\chi=A(\chi+\rho)$). & Lemma~\ref{p1:lem:iso}.
\end{hyptab}
\begin{proof}[Sketch of proof]
The condition for $z^\chi w$ ($w\in\wedge^bN_\C$) to be a holomorphic section of $\wedge^bT\otimes L$ is: $c_i(\chi)\ge-1$ along each $D_i$, and if $c_i(\chi)=-1$ then $u_i\wedge w=0$. So the weight-$\chi$ part is $G^b_\chi=\{w\in\wedge^bN_\C:u_i\wedge w=0\ (i\in C(\chi))\}$ (empty if some $c_i<-1$), and $G^\bullet_\chi\cong\wedge^{\bullet-s}(N/W)\otimes\wedge^sW$. The differential is $v_\chi\wedge$. If $v_\chi\bmod W=0$ the differential is $0$ and the Poincar\'e polynomial is $t^s(1+t)^{n-s}$; otherwise the complex is acyclic.
\end{proof}
\tV\ The closed formula agrees with a direct linear-algebra computation (\texttt{cp2/toric\_hh.py}) in 60 random cases ($\mathbb P^2$, $\mathbb P^1\times\mathbb P^1$, $\mathbb P^4$, $\mathbb P^2\times\mathbb P^2$, $(\mathbb P^1)^4$).

\begin{cor}\tP
(a) If $E_1^{a>0}=0$, the total Euler characteristic is the number of weights with $s(\chi)=n$ (the number of vertices) $=\chi(X)$. (Consistent with Prop.~\ref{p1:prop:euler}.) (b) Origin of the contributions: $s=0$ means $A(\chi+\rho)=0$, i.e.\ $\chi=-\rho$, and if $\rho\in M$ and all $r_i\le0$ it contributes $(1+t)^n=H^*(U)$ (the $Rj_*$ case). $s=n$ always satisfies the condition and gives the vertex weights, one $t^n$ each. For $0<s<n$ one needs $x=\chi+\rho\in A^{-1}W$ and $\langle x,u_i\rangle=r_i-1\ (i\in C)$, which has solutions only at ``resonances'' on the integral lattice. (c) For $s=n-1$ and all $r_i=1\ (i\in C)$, the condition holds automatically and classes $t^{n-1}(1+t)$ appear, in number equal to the number of lattice points on the toric curve $C$, namely $\deg(L|_C)+1$ (corresponding to the resonance $r_i=1$ in dimension 2).
\end{cor}

\subsection{Generic $\sigma$ and dimension 4}
\begin{obs}[\tV]
For generic $A$ (entries that are integers of very large size), $\Lambda=\{\text{vertices}\}\cup\{\text{resonances}\}\cup\{-\rho\}$ ($-\rho$ is included when $\rho\in M$ and all $r_i\le1$), and $P_{\HH}=$ ($U$-part) $+\#\text{vertices}\cdot t^n+$ (resonance classes for $r=1$). The $U$-part is $(1+t)^n$ if $\rho\in M$ and all $r_i\le0$, and $0$ otherwise. In exact arithmetic modulo the prime $2^{61}-1$, there were 0 mismatches in 3626 cases ($\mathbb P^4$ ($m=1,2$), $\mathbb P^2\times\mathbb P^2$, $(\mathbb P^1)^4$, $\mathbb P^1\times\mathbb P^3$, $\mathbb P^1\times\mathbb P^1\times\mathbb P^2$, $\rho\in\{-2,..,1\}^4$ and three non-integral $\rho$). Examples: $\mathbb P^4$, $m=1$: $[1,4,6,4,6]$; $\mathbb P^2\times\mathbb P^2$: $[1,4,6,4,10]$.
\end{obs}
\begin{obs}[New phenomenon: dependence on the arithmetic of $\sigma$\ \tV]\label{p1:obs:arith}
In dimension 4, $\HH^*$ is not determined by $(X,D,L,\text{residues})$ alone. For the product structure $\sigma$ on $\mathbb P^2\times\mathbb P^2$ (block-diagonal $A$) it is $[1,4,12,16,16]$ (the product of the results for two surfaces, K\"unneth), and for generic $A$ it is $[1,4,6,4,10]$. The difference $(0,0,6,12,6)=6t^2(1+t)^2$ is $t^2H^*(O_\tau)$ for six codimension-2 strata $\{\text{double point of one factor}\}\times(\mathbb C^*)^2$. Mechanism: on a stratum with $|C|=2$, if $u_i,u_j$ are symplectic with respect to $\sigma$, then $x_C(A,r)\in A^{-1}W$ is uniquely determined, and it contributes only under the arithmetic condition that this is a lattice translate of $\rho$. \tK\ $\HH$ jumps up on countably many proper analytic subsets of the parameter space of log-symplectic structures (upper semicontinuity of generalized Dolbeault cohomology), and the generic value is the minimum.
\end{obs}
\begin{conj}[General simple normal crossings $D$]\tK
For a general (non-toric) simple-normal-crossings log-symplectic $(X^{2k},D)$, general $\sigma$, and $L$ with constant residues $r_i\ne1$: $\HH^*(L)=H^*(X;\mathcal F)\oplus\bigoplus_{p\in D_{(n)}}\C[-n]$ ($D_{(n)}$ the set of points where $n$ components meet). Consistency with the Euler characteristic holds only when the $\chi$ of the lower-dimensional strata $D_k^\circ$ vanish (true in the toric case, unproved in general). The computations of \S\ref{p1:sec:toric} are for the toric case only.
\end{conj}

%--------------------------------------------------------------
\section{Local theory of non-nodal singularities}\label{p1:sec:local}
%--------------------------------------------------------------
When $D$ has cusps, tacnodes, triple points etc., Theorem~\ref{p1:thm:main2} does not apply. What follows are computational observations and conjectures. Local model: $\sigma=h\,\partial_z\wedge\partial_w$, module $\mathcal O\,h^s$, $d_0f=(hf_w+sfh_w)\partial_z-(hf_z+sfh_z)\partial_w$, $d_1(p,q)=h(p_z+q_w)-(1-s)(h_zp+h_wq)$. Globally $s=r=-m/3$ ($\mathbb{CP}^2$), $-m/2$ ($\mathbb{CP}^1\times\mathbb{CP}^1$).

\subsection{Stalk cohomology}
\tV\ For germs of quasi-homogeneous isolated singularities, the stalk dimensions $(h^0,h^1,h^2)$ were computed by weight-graded exact linear algebra (script \texttt{local\_stalk.py} etc.) (assuming that the cohomology of the formal germ equals that of the convergent germ):
\begin{center}\small
\begin{tabular}{lccc}\toprule
Singularity & generic $s$ & $s$ an integer $\le0$ & $s$ a spectral number mod $\Z$, $s<\mathrm{lct}$\\\midrule
smooth ($z$) & $(0,0,0)$ & $(1,1,0)$ & ---\\
double point ($zw$) & $(0,0,1)$ & $(1,2,2)$ & ---\\
cusp ($z^2-w^3$) & $(0,0,2)$ & $(1,1,2)$ & $(0,1,3)$\\
tacnode ($z^2-w^4$) & $(0,0,3)$ & $(1,2,4)$ & $(0,1,4)$\\
$A_4$ ($z^2-w^5$) & $(0,0,4)$ & $(1,1,4)$ & ---\\
$D_4$ ($zw(z-w)$) & $(0,0,4)$ & $(1,3,6)$ & $(0,1,5)$\\
$D_5$ ($w(z^2-w^3)$) & $(0,0,5)$ & $(1,2,6)$ & $(0,1,6)$\\\bottomrule
\end{tabular}
\end{center}
\begin{obs}[\tV\ Empirical rule]
In terms of the spectral numbers (lct is the smallest spectral number: smooth and double point $1$, cusp $5/6$, tacnode $3/4$, $D_4$ $2/3$, $D_5$ $5/8$), for $s\ne1$,
\[\text{stalk}_p=\C^{\mu_p}[-2]\ \oplus\ \begin{cases}\text{stalk}(Rj_*L_s)_p&(s<\mathrm{lct}_p)\\0&(s>\mathrm{lct}_p)\end{cases}\]
($\mu_p$ is the Milnor number $=$ Tjurina number). The stalk of $Rj_*$ is governed by the monodromy of the Milnor fibre: one class when $s\equiv\alpha_j\bmod\Z$ ($\alpha_j$ a spectral number), and a contribution to $H^0,H^1$ when $s\in\Z$. All 7 types, $s\in\frac1{12}\Z$ or $\frac18\Z$, $|s|\le5$: no contradiction. The explanation via Milnor monodromy is a consistency argument, not a proof. For $s=1$ (the twist by $K_X$) the stalk is infinite-dimensional ($d_Q=0$), but globally it can be handled by the $E_1$-term argument (Prop.~\ref{p1:prop:K}).
\end{obs}
\subsection{The triple-point anomaly and a global conjecture}
\tV\ On $\mathbb{CP}^2$ with $xy(x-y)$ (a triple point): $m=1,2,4,5,\ldots$ and $m=-1$ give $[0,1,4,0,0]$, $m=-5,-7,-8$ give $[0,0,4,1,0]$, and $m=-2,-3,-4$ give $[0,0,3,0,0]$. By the empirical rule above, with the $D_4$ spectral numbers $\{2/3,1,1,4/3\}$, $\mathrm{lct}=2/3$, $s=-m/3$: for $m\ge-1$, $s<2/3$ adds $H^*(U;L)=(0,1,0)$ giving $[0,1,4,0,0]$; for $m\le-5$, $s>2/3$ adds $H_c^*=(0,0,0,1,0)$ giving $[0,0,4,1,0]$; the anomalous $m=-2,-4$ have $s=2/3,4/3$, where one class of $\C^4$ (degree 2) and one class of $H_c$ (degree 3) cancel (the $d_3$ of the spectral sequence $H^p(X,\mathcal H^q)\Rightarrow\HH$). \tK\ This explanation is conjectural. The spectral numbers of the cusp ($5/6,7/6$), tacnode ($3/4,5/4$) and $D_5$ ($5/8,\ldots$) do not lie in $\frac13\Z$, and on $\mathbb{CP}^1\times\mathbb{CP}^1$ ($s\in\frac12\Z$) the spectral numbers of $A_k$ have no half-integer other than $1$, so no anomaly appears.

\tV\ Unequal residues: for a triple point with module $\prod h_i^{s_i}$ ($h=zw(z-w)$), in 533 random $(s_1,s_2,s_3)\in(\frac16\Z\cap[-2,2])^3$ ($s_i\ne1$) the stalk is $\C^4[-2]\oplus R$, where $R=H^*(U_{\rm loc};L_s)$ if $s_1+s_2+s_3<2$ and $0$ otherwise. Whereas for a double point the threshold is per branch, the threshold at a triple point is on the \emph{sum} of the residues. \tK\ General principle: the threshold is the ratio of the weight of the module to the logarithmic canonical weight with respect to the symmetric torus action on the germ.

\begin{conj}[Refined Riemann--Roch (equal residues, isolated singularities)]\tK
Let $X$ be a compact surface, $\sigma\in H^0(-K_X)$ with zero curve $D$ reduced with only isolated singularities, $L$ a Poisson module with $c_1(L)=sK_X$, $L_s$ the rank-one local system on $U$ (monodromy $e^{2\pi is}$ around each branch), and $\mathcal F_s$ the constructible complex which is $L_s$ on $U$ and at $p\in D$ is the stalk of $Rj_*L_s$ if $s<\mathrm{lct}_p$ and $0$ if $s>\mathrm{lct}_p$. (a) For $s\notin\mathrm{Sp}(D)$ (the spectral numbers of all singularities of $D$, together with $1$): $\HH^*=H^*(X;\mathcal F_s)\oplus\bigoplus_p\C^{\mu_p}[-2]$. (b) The Euler characteristic is always $\chi(X)$ ($\chi(D)=\sum\mu_p$, from $\delta$-invariants \tP). (c) For $s\ll0$, $\mathcal F_s=Rj_*L_s$, and $\HH^*=H^*(U;L_s)\oplus\bigoplus\C^{\mu_p}[-2]$ depends only on $s\bmod\Z$ and vanishes in degrees $\ge3$. (d) For $s\in\mathrm{Sp}(D)$, rank-one cancellations can occur.
\tV\ On $\mathbb{CP}^2$ (smooth, nodal, cuspidal, conic plus secant, conic plus tangent, triangle, triple point) and $\mathbb{CP}^1\times\mathbb{CP}^1$ (4 examples), for $m=-15..15$ ($-9..8$), no contradiction except in (d). Not proved: the stalk results are only for quasi-homogeneous germs, and the global identification $\mathbb H(C_m)\simeq(\mathcal F_s+\text{skyscraper})$ is a fit. For non-reduced $D$ (double line, double conic) the stalks are infinite-dimensional and the global dimensions grow linearly in $m$ (slope $\deg\mathcal O(1)|_C$: $1$ for a double line on $\mathbb{CP}^2$, $2$ for a double conic on $\mathbb{CP}^1\times\mathbb{CP}^1$).
\end{conj}

%--------------------------------------------------------------
\section{Comparison with earlier work}\label{p1:sec:compare}
%--------------------------------------------------------------
\subsection{Hodge theory (formal coefficients)}
Gualtieri proved a generalized Hodge star, the Born--Infeld inner product, Laplacian equalities, the Hodge decomposition and the $dd^{\Jc}$-lemma on generalized K\"ahler manifolds (\tC\ Gualtieri, \emph{Generalized geometry and the Hodge decomposition}, arXiv:math/0409093; Gualtieri, \emph{Generalized K\"ahler geometry}, Comm.\ Math.\ Phys.\ 331 (2014) 297--331, arXiv:1007.3485; Cavalcanti, \emph{New aspects of the $dd^c$-lemma}, arXiv:math/0501406, and J.\ Geom.\ Phys.\ 57 (2006) 121--132. Theorem numbers are omitted because these were not read closely \tQ). These are theories with trivial or $H$-twisted coefficients; a Bochner--Kodaira--Nakano identity with line-bundle (curved) coefficients has not been written, as far as surveyed. The results of this part can be positioned as filling this gap for $\beta$-transforms, not by Hodge theory but by comparing logarithmic complexes with constructible sheaves.

\subsection{Pal--Poddar}
\tC\ Pal--Poddar, \emph{Strong generalized holomorphic principal bundles}, arXiv:2404.18113, treat compact generalized complex manifolds of regular type (constant type $k$) whose symplectic leaf space is a K\"ahler orbifold, and give Hodge theory and Serre duality for strong generalized holomorphic (SGH) bundles (transversal Dolbeault isomorphism Thm.~8.9, Cor.~8.10 (numbers \tQ); Hodge decomposition and Serre duality Thm.~11.2, 11.4, 11.6, 11.7), K\"ahler identities (Prop.~11.8, Lemma~11.9; numbers \tQ), and Kodaira-type and Serre-type vanishing theorems for positive SGH line bundles (Def.~11.10, 11.11; Thm.~11.12). The hypotheses of regular type and orbifold leaf space are used to guarantee that the foliation is a smooth leaf bundle and to guarantee transversal Hodge theory (ellipticity). Differences from the present note:
\begin{itemize}
\item The type-jumping case (the type jumps on the degeneracy divisor of $\sigma$) is outside the framework of Pal--Poddar; examples such as $\mathbb{CP}^2$ in this part are of this kind.
\item SGH bundles form a proper subclass of generalized holomorphic bundles (flat along the leaves and holomorphic transversally). The Poisson modules of this part are general GH bundles on a $\beta$-transformed structure.
\item The statement and proof structure of Thm.~11.12 were checked by re-reading \S11: $H^q(M,(\mathcal G^*M)^p\otimes E)=0$ for $p+q>k$, and $H^q(M,E'\otimes E^m)=0$ for $m\ge m_0$, $q>0$. The hypotheses (compactness, a K\"ahler orbifold leaf space, positivity of $E$) are used for transversal Hodge theory and for positivity of the curvature term of a Bochner--Kodaira--Nakano type identity.
\end{itemize}
Doubts about the original text noticed in this survey (likely misprints or omissions in the source, but unverified \tQ): (i) the statement carries a factor $(i/2\pi)$ while the proof has no $1/2\pi$ (harmless); (ii) the sign discrepancy about $gc_1$ comes from the differing definitions of the characteristic polynomial in (9.9) and of $\det(I\pm tA/2\pi i)$ in \S10.1; (iii) ``a positive class gives a metric of positive curvature'' does not follow from Def.~11.11 and is assumed in the proof of Thm.~11.12; (iv) Def.~11.11 does not state compactness or the K\"ahler condition; (v) the proof of Lemma~11.9(1) is a computation in pointwise local frames; (vi) $T^{1,0}$ versus $T^{1,0*}$ in Thm.~13.4(2)(3), and the binomial coefficient for $q=n$ in Example 13.10.

\subsection{The work of Tengzhou Hu}
\tC\ T.\ Hu, \emph{Kodaira vanishing theorems for K\"ahler Lie algebroids}, arXiv:2403.10876;\\ T.\ Hu, \emph{Lie groupoid Riemann--Roch--Hirzebruch theorem and applications}, arXiv:2403.13424.
\begin{itemize}
\item \emph{Setting} (2403.13424): a Lie groupoid $G\rightrightarrows M$ (Def.~1; hypotheses of properness or $s$-connectedness are not stated) and its Lie algebroid $\mathfrak g$. A ``complex Lie algebroid'' (Def.~2) is a \emph{real} Lie algebroid $\mathfrak g$ with an integrable complex structure $J$ in the fibre direction (this is different from the complex Lie algebroid $L$ as a subbundle of $\bT_\C$). Kähler Lie algebroids are Def.~11, positive bundles Def.~12.
\item \emph{Claims} (2403.13424): Thm.~4 (Riemann--Roch) states, for $M$ compact, $\dbar_{\mathfrak g}$ elliptic (claimed, but no proof was found) and $\Omega$ a $\mathfrak g$-invariant density (unimodularity),
\[\mathrm{Ind}_\Omega(\dbar_{\mathfrak g}+\dbar_{\mathfrak g}^*)=(2\pi i)^{-k}\int_M(-1)^{\mathrm{rk}\,\mathfrak g}\bigl\langle\ch^{\mathfrak g}(\mathfrak g^{*p,0})\ch^{\mathfrak g}(E)\Td^{\mathfrak g}_\C(\mathfrak g),\Omega_{\mathfrak g}\bigr\rangle.\]
Thm.~7 (vanishing): if $E$ is positive and $p+q>n$, the \emph{kernel} of the $\mathfrak g$-Dolbeault Laplacian vanishes (vanishing of cohomology does not follow: for a general Lie algebroid $\mathfrak g$ the Laplacian is not necessarily elliptic and Hodge theory is unavailable). In 2403.10876, a K\"ahler Lie algebroid is Def.~2, the K\"ahler identities are Thm.~1 and Prop.~4, the BKN identity is Lemma~4, and the main vanishing theorem is in the introduction.
\item \emph{The generalized Dolbeault complex is outside the scope}\ \tP\ (derived; comparison with the source only to the extent indicated below): $L$ is a complex Lie algebroid without $J$, and $\wedge L^*$ in general has no $(p,q)$ bigrading. For a holomorphic Poisson structure $\sigma$ the real Lie algebroid $A_{\mathbb R}$ falls within Hu's setting, but its Dolbeault complex is $\wedge\mathfrak g^{*0,q}$ built from $\bar\sigma$, which differs from the complex $(T^{0,1}\bowtie A^{1,0})$ of this part.
\item \emph{Does the index see $m$?}\ \tK: Hu's index is a groupoid index, and is expected to depend polynomially on $m$, roughly like $\int\ch(E)\Td(T^*)\ch(\wedge^pT)$, whereas the total Euler characteristic of this part is independent of $m$ (Prop.~\ref{p1:prop:euler}). Moreover the residue $-m/3$, which governs the periodic jumps, is not part of Hu's data. Hence Hu's theorem is not ``Riemann--Roch for generalized Dolbeault cohomology''.
\item \emph{Inconsistencies in the source}: 2403.13424 (obtained through summaries): $n$ is undefined in Thm.~7; the sign constants differ among Thm.~5, Prop.~1 and Thm.~6; ``unimodular'' is defined in two ways. 2403.10876 (checked by reading further and part of the appendix \tC): Lemma~5 gives $[L,\Lambda]=(p+q-n)\,\mathrm{id}$, but the proof of Thm.~2 writes $\Delta_{\bar\partial_{\mathfrak g}}=\Delta^{1,0}-[L,\Lambda]=\Delta^{1,0}+(p+q-n)\,\mathrm{id}$, and the sign in the last equality does not match ($-(p+q-n)=n-p-q$). Moreover the substitution $\nabla^2=\sqrt{-1}L$ means $\sqrt{-1}R=-L$, opposite to the definition of ``positive'' ($\sqrt{-1}R>0$, Def.~13). The statement itself (the kernel vanishes for positive $E$ when $p+q>n$) agrees with classical Kodaira--Nakano; so the printed proof has a sign slip, not the claim. Misprints in Cor.~1 (arrows and the symbols $s/t$). The appendix (Lemmas 6--10, the computation of $[\bar\partial_{\mathfrak g},\iota_\omega]$) could not be verified because the indices are hard to read. The reference list and publication data are unchecked. \tQ
\item \emph{Correction to the earlier version of this part (preliminary survey)}: the earlier version gave the author as ``S.~Hu'' and a different title; both were wrong, and the above is correct.
\end{itemize}

\subsection{Integrability and groupoids}
\tC\ Gualtieri--Li (arXiv:1206.3674, IMRN 2014) classify Hausdorff symplectic groupoids of log-symplectic manifolds and show that the Lie algebroid $T_X(-\log D)$ of logarithmic vector fields tangent to a smooth $D$ is integrable. Bailey--Gualtieri (arXiv:1611.03850) show the equivalence between integration of a generalized complex structure to a weak holomorphic symplectic groupoid and integrability of the real Poisson structure (Thm.~7.1: integrability of the underlying real Poisson structure is necessary and sufficient; Thm.~7.1 itself says nothing about Hausdorffness or $s$-connectedness, and the construction Thm.~7.6 uses an $s$-connected, $s$-simply connected integration). Laurent-Gengoux--Sti\'enon--Xu (arXiv:0803.2031): Thm.~3.17, 3.22. Crainic--Fernandes (arXiv:math/0105033): Thm.~4.1, Cor.~5.2. These provide groupoids as input for index theory, but do not treat integration of the complex Lie algebroid $L$ itself.

%--------------------------------------------------------------
\section{Open problems and next steps}\label{p1:sec:open}
%--------------------------------------------------------------
\begin{enumerate}
\item \textbf{The treasure not yet touched}: a generalization of the Kodaira embedding theorem. Since $\HH^0$ consists of Casimir sections (Observation~\ref{p1:obs:casimir}), it must be formulated not as an embedding of $M$ itself but as a map to the quotient separated by Casimir sections (e.g.\ the leaf space). The next concrete step is to study base points and separation properties of the linear system of Casimir sections in the $\mathbb{CP}^2$ example ($\HH^0=\C F^{m/3}$, i.e.\ 1-dimensional for $m\in3\Z_{\ge0}$).
\item \textbf{Generalization of Serre-type vanishing} (Conjecture~\ref{p1:conj:A3}): a definition of a ``positive'' generalized holomorphic line bundle, and a construction of a filtration playing the role of the $E_1$ term for $\Jc$ that are not $\beta$-transforms. The comparison with Pal--Poddar Thm.~11.12 was done by re-reading \S11 (the groups differ; the transversal-degree reading is consistent).
\item \textbf{Formulation of Riemann--Roch}: instead of the total Euler characteristic, one must fix a refinement (the $E_1$ term of a filtration by the type of $\Jc$, a periodic quasi-polynomial, the indicial roots $H^*(\partial U)$). Proving Conjecture~\ref{p1:conj:I3} requires an index theory for the complex Lie algebroid $L$ (unimodularity, and a derivation of Fredholmness $\Leftrightarrow r\notin\Z$). An APS-type extension to generalized complex manifolds with boundary is a still more distant problem.
\item \textbf{Beyond log-symplectic}: proofs for the local theory of non-nodal singularities (only quasi-homogeneous germs were computed), a global formula for unequal residues, non-reduced $D$ (infinite-dimensional stalks and linear growth), and non-toric examples in dimension $\ge4$ (see Observation~\ref{p1:obs:arith}).
\item \textbf{Generalized K\"ahler side}: write down the Bochner--Kodaira--Nakano identity with line-bundle coefficients in the generalized K\"ahler framework. The interaction of the $H$-twist with the curvature term may be an obstacle \tK.
\item \textbf{Resolution of the unverified items of this part}: the \tQ\ items of \S\ref{p1:sec:ledger}. In particular Mokri's theorem, a re-derivation of LGSX Lemma~4.22, a reading of Gualtieri 2011 \S4, and the remaining bibliographic data.
\end{enumerate}

\paragraph{Priorities for the next steps (author's proposal).}
(a) As a comparison case in which the type does not jump, compute the same ``Poincar\'e polynomial'' for the regular-type examples of Pal--Poddar and compare with the periodic jumps of this part. (b) For the $\mathbb{CP}^2$ example, use $\HH^0=\C F^{m/3}$ to study the geometry of the linear system of Casimir sections and look for the ``right form'' of a Kodaira-type theorem. (c) As a basis for an index-theoretic proof of Conjecture~\ref{p1:conj:I3}, formulate rigorously the relation between the weighted $b$-complex on the real Lie algebroid $T_X(-\log D)$ and the residue $r$.

%--------------------------------------------------------------
\section{Status of the claims and extent of verification}\label{p1:sec:ledger}
%--------------------------------------------------------------
\subsection{Status of each claim}
\begin{center}\footnotesize
\begin{tabular}{>{\raggedright\arraybackslash}p{4.0cm}>{\raggedright\arraybackslash}p{3.0cm}>{\raggedright\arraybackslash}p{8.4cm}}\toprule
Claim & Status & Dependencies / gaps\\\midrule
Theorem~\ref{p1:thm:main1} (Poisson module $=$ log connection) & \tP & Any even dimension, simple normal crossings. The residue formula $c_1=-\sum r_i[D_i]$ is quoted from Esnault--Viehweg.\\
Theorem~\ref{p1:thm:cp2} (table for $\mathbb{CP}^2$) & \tP\ \tV & Hand derivation and exact computation for $m=-15..15$ agree. Depends on the identification of complexes (Prop.~\ref{p1:prop:ident}(b)).\\
Theorem~\ref{p1:thm:main2} (main theorem for surfaces) & \tP & Prop.~\ref{p1:prop:ident}(b) (LGSX Thm.~4.19 and Mokri's theorem), Deligne's comparison, the Laurent-tail computation of local cohomology (standard but not written out). Requires $r_i\ne1$.\\
Theorem~\ref{p1:thm:B} (closed formula) & \tP\ \tV & Theorem~\ref{p1:thm:main2} plus (H1)--(H4). (H4) ($H^1(U;L_\chi)=0$) was checked in 8 examples; in general it is a condition on $V^1(U)$. Agrees for 196 values.\\
Theorem~\ref{p1:thm:A1} (Serre-type) & \tP & Only classical Serre (Hartshorne III.5.2) and the acyclic resolution lemma. Any dimension.\\
Theorem~\ref{p1:thm:A2} (periodic stabilization) & \tP & Depends on Theorem~\ref{p1:thm:main2} and its hypotheses ($X$ simply connected, $U$ affine).\\
Prop.~\ref{p1:prop:euler} (Euler characteristic) & \tP & HRR only.\\
Theorem~\ref{p1:thm:dual} (duality) & \tP\ \tV & Lemmas D and E machine-checked for $n\le4$. The agreement of the Koszul signs of the $\dbar$ part and the $d$ part is by hand only. Hodge theory for elliptic complexes is standard.\\
Theorem~\ref{p1:thm:weight} (weight formula) & \tP\ \tV & Assumes $E_1^{a>0}=0$. Agreement with linear algebra in 60 examples, and in 3626 examples with generic $A$.\\
Obs.~\ref{p1:obs:arith} (arithmetic dependence of $\sigma$) & \tV & A few examples in dimension 4. The explanation of the mechanism is a sketch.\\
Local theory in \S\ref{p1:sec:local} and the triple-point explanation & \tV\ \tK & Quasi-homogeneous germs only. The global identification is a fit.\\
Conjectures~\ref{p1:conj:A3}, \ref{p1:conj:I3} & \tK & Neither definitions nor proofs are complete.\\
Generalization of the Kodaira embedding theorem & --- & Not treated (\S\ref{p1:sec:open}).\\\bottomrule
\end{tabular}
\end{center}

\subsection{Unverified or delicate items (\tQ)}
\begin{itemize}
\item \textbf{Bibliographic data (from memory)}: the journal and pages of Goto, the IMRN volume/pages of Gualtieri--Li, the OUP data of Gualtieri ``Branes'', the volume/pages of Mokri, of Arapura, the journal of Lang--Jia--Liu, the IMRN Art.\ ID of LGSX, the volume/issue of Pal--Poddar. The numbers in Deligne LNM 163, Crainic--Fernandes, Fulton, Esnault--Viehweg. The numbers and contents of Thm.~8.9, Prop.~9.3, Ex.~10.18, Prop.~11.8, Lemma~11.9, Ex.~13.10 of Pal--Poddar. (Checked in v2: Hitchin Moscow Math.\ J.\ 12 (2012) 567--591, Pym Lett.\ Math.\ Phys.\ 108 (2018) 573--632, Hausel--Hunsicker--Mazzeo Duke 122 (2004) 485--548, Timmerscheidt J.\ reine angew.\ Math.\ 379 (1987) 152--171, title and authors of Pal--Poddar.)
\item \textbf{Gualtieri 2011 (arXiv:math/0703298)}: of roughly 160,000 characters, the first roughly 100,000 (up to the middle of \S3) were read closely; \S4 (the Darboux theorem etc.) was not read in full. In v2 the following were checked: type Def.~3.5, canonical bundle Def.~3.7, ellipticity Prop.~3.13/Cor.~3.14, Picard group Def.~3.18, Darboux Thm.~4.3 (regular points), regular points Def.~4.1; the type-jumping example is in \S4.1, $\rho=z_1+dz_1\wedge dz_2$ on $\mathbb R^4$.
\item \textbf{Hitchin, \emph{Poisson modules and generalized geometry} (arXiv:0905.3227, title and author checked)}: it only states in the introduction that GH bundles $=$ Poisson modules; it contains no $\Jc_\sigma$ notation, no identification of hypercohomology, no obstruction by the Atiyah class, no $\mathbb{CP}^2$ example and no Lie algebroid $L$ (according to the text obtained). This part instead cites Branes on Poisson varieties (arXiv:0710.2719) and LGSX. (v1 attributed this part to Gualtieri and said ``confirmed''; that was wrong in both author and content.)
\item \textbf{$\mathrm{Pic}(\Jc)=\mathbb H^1(\mathcal L_{\log})$}: could not be checked in the text of any reference. This part does not use it. What was used is only the direct derivation of existence and uniqueness of Poisson module structures.
\item \textbf{LGSX Lemma~4.22, Mokri's theorem}: not re-derived. The standard facts of Hodge theory and elliptic complexes, and Esnault--Viehweg, Deligne, Fulton--Deligne, Timmerscheidt, Arapura and Crainic--Fernandes, are quoted, and their ranges of applicability were not checked against the sources.
\item \textbf{The two papers of Hu}: for 2403.10876 the first roughly 200,000 characters (about 500,000 in the full expanded HTML); for 2403.13424 only the overview and the statements of theorems.
\item \textbf{Pal--Poddar}: \S\S5--11, 13 were re-read (22 November 2026); the extraction went through summaries, so details of formulas are unverified.
\item \textbf{Limits of the computations}: exact linear algebra was done for a finite range of $m$, a finite number of curves (Fermat and Hesse cubics, etc.), and modulo primes ($32003$, $2^{61}-1$). Dependence on the choice of curve (random smooth curves), and whether reduction modulo $p$ represents the characteristic-0 result, were not verified.
\end{itemize}

\subsection{Record of the work}
All scripts use exact arithmetic in Python (standard library) and were saved in \texttt{cp2/} (\texttt{poisson.py, table.py, neg\_check.py, predict.py, sing\_table.py, p1p1.py, local\_stalk.py, local\_unequal.py, ncd\_cochain\_check.py, ncd\_local\_weights.py, ncd\_global\_check.py, serre\_rr\_check.py, serre\_dual\_check.py, duality\_ndim.py, duality\_ndim\_exact.py, toric\_hh.py, toric\_crit.py, toric\_generic\_check2.py, index\_route\_check.py}).

\clearpage
\part{Euler characteristic, toric varieties, local formulas}
\setcounter{section}{0}\renewcommand{\thesection}{II.\arabic{section}}
\noindent\textit{Part II of the series; originally published as a separate note (\url{https://hylt-math.github.io/2026/11/22/generalized-geometry-poisson-part2/}). Within this paper, Part-internal references use the numbering II.$\ast$.$\ast$.}\par
\section{Context and overview of results}\label{p2:sec:intro}
This part continues Part I (Claude, Mori, \emph{Cohomology of generalized holomorphic line bundles on the generalized complex structure of a holomorphic Poisson structure}, v2, \url{https://hylt-math.github.io/2026/11/22/generalized-geometry-poisson/}; ``Part I''). We do not repeat the setting of Part I (the $\beta$-transformed generalized complex structure $\Jc_\sigma$ of a holomorphic Poisson structure $\sigma$, Poisson modules $=$ generalized holomorphic line bundles, generalized Dolbeault cohomology $\HH^*(E)=H^*(X;C(A;E))$, residues $r_i$ in the log-symplectic case); notation follows the notation table of Part I. Results of Part I are cited as ``Part I, Theorem X'' (v2 numbering).

\paragraph{Summary of results.}
\begin{itemize}
\item \textbf{Theorem E} (\S\ref{p2:sec:E}) \tP: in any dimension, for any holomorphic Poisson $\sigma$ and any vector bundle $E$ with flat $A$-connection, $\sum(-1)^k\dim\HH^k(E)=(-1)^n\,\rk E\,\chi(X)$. This clarifies the hypotheses of Part I, Prop.\ 6.1 and extends it to vector bundles.
\item \textbf{Theorem C1} (\S\ref{p2:sec:toric}) \tP: on a toric variety, the cohomology of the log de Rham complex $\mathrm{DR}_{\log}(L)$ is written completely in terms of reduced cohomology of a simplicial complex $V_S$ determined by the residues (CLS weight decomposition and K\"unneth).
\item \textbf{Theorem C2} (\S\ref{p2:sec:toric}) \tP\ \tV: on toric surfaces, if all $r_i\ne1$, a closed formula for the Poincar\'e polynomial of $\HH^*$ with unequal residues (from Part I's surface theorem and Theorem C1).
\item \textbf{Theorem T} (\S\ref{p2:sec:toricn}) \tP\ \tV: for a generic log-symplectic structure on a toric variety of any even dimension, $\HH^*=H^*(U;L_U)\oplus\C^{\chi(X)}[-n]$ (when $L$ has no positive residues), with a table of the role of each hypothesis.
\item \textbf{Resonant case $r_i=1$} (\S\ref{p2:sec:res}) \tV\ \tK, a \textbf{local formula for quasi-homogeneous plane curve singularities} and a rule for \textbf{triple points with unequal residues} (\S\ref{p2:sec:local2}) \tV\ \tK.
\item \textbf{A conjecture in general dimension} (\S\ref{p2:sec:conj}) \tK, and an outlook on index theory.
\end{itemize}
In \S\ref{p2:sec:chk} we checked whether Part I (v1, v2) needs correction in view of these results. Conclusion: no theorem or proposition has to be withdrawn; one phrase about the local formula should be made explicit.

\section{Notation (new in this part)}\label{p2:sec:notation}
\begin{longtable}{>{\raggedright\arraybackslash}p{3.6cm}>{\raggedright\arraybackslash}p{11.4cm}}
\toprule Symbol & Meaning\\ \midrule\endhead
$N,\ M$ & lattice $N\cong\Z^n$ of the torus and dual lattice $M=\mathrm{Hom}(N,\Z)$; $N_\C=N\otimes\C$, $M_\C=M\otimes\C$\\
$\Sigma,\ u_i$ & complete smooth fan and primitive ray generators $u_i\in N$; $D_i$ the torus-invariant divisors, $D=\sum D_i$ (toric boundary, simple normal crossings)\\
$L=\mathcal O(\sum a_iD_i)$ & $a_i\in\Z$; regarded as a subsheaf of rational functions\\
$\rho\in M_\C$ & flat log connection $\nabla=d+\rho$ (constant closed log 1-form $\rho=\sum\rho_j\,\dlog z_j$)\\
$r_i$ & residue $r_i=\langle\rho,u_i\rangle-a_i$ (Part I, Thm.\ 2.1(d)); $c_1(L)=-\sum r_i[D_i]$\\
$A$ & constant non-degenerate skew matrix; $\sigma=\tfrac12\sum A_{ij}E_i\wedge E_j$, $E_j=z_j\partial_j$; $A:M_\C\to N_\C$\\
$\chi\in M$ & character (weight) of the torus; $z^\chi$ is the monomial of weight $\chi$\\
$c_i(\chi)$ & $\langle\chi,u_i\rangle+a_i$; $z^\chi\cdot(\text{generator})$ is regular along $D_i$ iff $c_i(\chi)\ge0$\\
$C(\chi),\ W(\chi),\ s(\chi)$ & $C(\chi)=\{i:c_i(\chi)=-1\}$, $W(\chi)=\mathrm{span}(u_i:i\in C(\chi))\subset N_\C$, $s=\dim W$\\
$\Lambda$ & set of characters summed over in the weight formula (Part I, Thm.\ 6.4)\\
$\Delta_\Sigma$ & simplicial complex whose vertices are the rays of $\Sigma$ (simplices = ray sets of cones)\\
$V_S$ & full subcomplex of $\Delta_\Sigma$ spanned by the vertex set $S$; $\tilde h^q(V_S)$ = dimension of reduced cohomology (convention: for $V_\emptyset=\emptyset$, $\tilde h^{-1}=1$, others $0$)\\
$\mathrm{DR}_{\log}(L)$ & $(\Omega^\bullet(\log D)\otimes L,\nabla)$ (as in Part I)\\
$F$ & the constructible complex $j_{2!}Rj_{1*}L_U$ of Part I, Thm.\ 4.1\\
$\mathrm{Sp}(f)$, $\mathrm{lct}$ & spectrum of an isolated singularity $f=0$ and its minimum (log canonical threshold); for quasi-homogeneous $f$ (weights $w_1,w_2$, degree $d$), $\mathrm{lct}=(w_1+w_2)/d$\\
$\mu_p$ & Milnor number of the singularity $p$\\
$s$ & residue parameter: exponent in the local model $h^s$ (as in Part I)\\
$P_{\HH}(t)$ & $\sum_k\dim\HH^k\,t^k$\\
\bottomrule
\end{longtable}
\noindent Tags: \tC\ literature, \tP\ proved by the author (standard facts cited), \tV\ checked by computer (exact linear algebra, finitely many examples), \tK\ conjecture/outlook, \tQ\ unverified.

\section{Euler characteristic theorem (any dimension)}\label{p2:sec:E}
\begin{thm}[Theorem E]\tP\label{p2:thm:E}
Let $X$ be a compact complex manifold ($\dim_\C X=n$), $\sigma$ a holomorphic Poisson structure, $A=(T^*X)_\sigma$, and $E$ a holomorphic vector bundle with flat $A$-connection $\nabla$. For the complex $K^\bullet=(\wedge^\bullet T_X\otimes E,\ d_E)$ (differential $\C$-linear), the hypercohomology $\HH^k(E):=H^k(X;K^\bullet)$ is finite dimensional and
\[\sum_k(-1)^k\dim\HH^k(E)=(-1)^n\,\rk(E)\,\chi(X).\]
\end{thm}
\begin{hyptab}
$X$ compact & each $H^a(X,K^b)$ is finite dimensional; Hirzebruch--Riemann--Roch applies. & (i) and (ii).\\\midrule
$K^b$ coherent & Euler characteristics computable by HRR; descending induction works. & (i): induction on truncations.\\\midrule
Flat $A$-connection & $d_E^2=0$ (a complex). & Definition of hypercohomology; not used in the Euler characteristic computation itself.\\\midrule
(Unused) non-degeneracy of $\sigma$, positivity of $E$ & --- & Do not appear in the proof; hence no dependence on $c_1(E)$.
\end{hyptab}
\begin{proof}
(i) The differential is $\C$-linear but not $\mathcal O$-linear; still the truncations $K^{\ge b}$ are subcomplexes and there are triangles $K^{\ge b+1}\to K^{\ge b}\to K^b[-b]\to$. Each $K^b$ is coherent and $X$ compact, so $H^*(X;K^b)$ is finite dimensional. By descending induction on $b$, all $H^*(X;K^{\ge b})$ are finite dimensional and $\chi(K^{\ge b})=\chi(K^{\ge b+1})+(-1)^b\chi(X,K^b)$. Hence $\sum(-1)^k\dim\HH^k=\sum_b(-1)^b\chi(X,\wedge^bT\otimes E)$.\\
(ii) By HRR, $\chi(X,\wedge^bT\otimes E)=\int\ch(E)\ch(\wedge^bT)\Td(T)$. With Chern roots $x_i$ of $T$, $\sum_b(-1)^b\ch(\wedge^bT)=\prod(1-e^{x_i})$, $\Td(T)=\prod x_i/(1-e^{-x_i})$, and $(1-e^{x})/(1-e^{-x})=-e^{x}$, so the integrand is $(-1)^n\ch(E)e^{c_1(T)}c_n(T)$. Since $c_n(T)$ already has top degree, only the degree-0 part $\rk E$ of $\ch(E)e^{c_1(T)}$ survives, and Gauss--Bonnet gives $(-1)^n\rk E\int c_n(T)=(-1)^n\rk E\,\chi(X)$.\qedhere
\end{proof}
\tV\ Recomputed from Bott's formula for $\mathbb{P}^n$ ($n\le6$), $E=\mathcal O(m)$ ($m=-12..12$): $\sum_b(-1)^b\chi(\mathbb{P}^n,\wedge^bT(m))=(-1)^n(n+1)$ (\texttt{cp2/roundD\_check.py}).

\begin{rem}[Meaning\ \tP]
The left side does not depend on the first Chern class of $E$. The classical $\chi(X,E)=\int\ch(E)\Td(X)$ is the Euler characteristic of the column $b=0$, not of the full complex. The ``Riemann--Roch polynomial in $m$'' of the generalized Dolbeault complex is constant; the dependence on $m$ lives only in the Poincar\'e polynomial (dimensions degree by degree) (Part I, \S I.1 observation; here \S\ref{p1:sec:toric}).
\end{rem}
\begin{rem}[General generalized complex structures\ \tK]
For general $\Jc$ and a flat $L$-connection $E$, Atiyah--Singer gives $\mathrm{ind}(\wedge L^*\otimes E)=\int\ch(E)\cdot S$ ($S$ the symbol class of $(\wedge L^*,d_L)$). By Chern--Weil and flatness, $\ch_j(E)$ should lie in the kernel of $H^{2j}_{dR}(M)\to H^{2j}(L)$, so only the rank should matter. Not proved.
\end{rem}

\section{Toric varieties: the log de Rham complex and the surface formula}\label{p2:sec:toric}
Setting as in the notation table. Since $\Omega^1(\log D)\cong\mathcal O\otimes M$ (frame $\dlog z_j$), $\mathrm{DR}_{\log}(L)=(L\otimes\wedge^\bullet M_\C,\ \nabla)$, and on a monomial $z^\chi$ of weight $\chi$, $\nabla=(\chi+\rho)\wedge$.

\begin{thm}[Theorem C1: cohomology of $\mathrm{DR}_{\log}$]\tP\label{p2:thm:C1}
Let $X=X_\Sigma$ be a complete smooth toric variety, $L=\mathcal O(\sum a_iD_i)$, $\nabla=d+\rho$ ($\rho\in M_\C$). Put $S=\{i:\ \langle\rho,u_i\rangle-a_i>0\}$ (for $\rho\in M$: the $i$ with $r_i\ge1$). Then
\[\sum_k\dim H^k(X;\mathrm{DR}_{\log}(L))\,t^k=\begin{cases}(1+t)^n\sum_q\tilde h^{q-1}(V_S)\,t^q&(\rho\in M),\\ 0&(\rho\notin M).\end{cases}\]
\end{thm}
\begin{hyptab}
Torus-invariant structure ($\rho$ constant, $L$ toric) & The complex is graded by weights $\chi\in M$ and each weight can be computed independently. & First decomposition in the proof.\\\midrule
Completeness (compactness) & Hypercohomology is finite dimensional; the Cech complex is computed on finitely many affine opens. & Cech complex.\\\midrule
$\Omega^1(\log D)\cong\mathcal O\otimes M$ & Allows a product decomposition with the Koszul complex (K\"unneth). & The weight-$\chi$ part has the form $\check C^\bullet_\chi\otimes(\wedge M_\C,(\chi+\rho)\wedge)$.\\\midrule
CLS weight formula & $H^a(X;L)_\chi$ is expressed by reduced cohomology of a simplicial complex. & Final identification.
\end{hyptab}
\begin{proof}
The Cech complex for the $T$-equivariant affine cover decomposes by weight $\chi$; the weight-$\chi$ part is the tensor product complex $\check C^\bullet(L)_\chi\otimes(\wedge^\bullet M_\C)$ with differential $d_{\check C}\otimes1\pm1\otimes(\chi+\rho)\wedge$. By K\"unneth over a field, $H^k(X;\mathrm{DR}_{\log})_\chi=\bigoplus_{a+b=k}H^a(X;L)_\chi\otimes H^b(\wedge M_\C,(\chi+\rho)\wedge)$. The Koszul complex is acyclic if $\chi+\rho\ne0$, and for $\chi+\rho=0$ (requiring $\rho\in M$) it is $\wedge^bM_\C$ of dimension $\binom nb$ with zero differential. So for $\rho\notin M$ everything vanishes, and for $\rho\in M$ only the weight $\chi=-\rho$ contributes, giving $(1+t)^n\sum_a h^a(X;L)_{-\rho}t^a$. Finally, by the theorem of Cox--Little--Schenck (\tC\ \emph{Toric Varieties}, GSM 124, AMS 2011, Thm.\ 9.1.3\tQ): $H^a(X,\mathcal O(\sum a_iD_i))_m\cong\tilde H^{a-1}(V_{D,m};\C)$, where $V_{D,m}$ is the full subcomplex on the rays with $\langle m,u_i\rangle<-a_i$. For $m=-\rho$ the condition reads $\langle\rho,u_i\rangle>a_i$, i.e.\ $i\in S$.\qedhere
\end{proof}

\begin{thm}[Theorem C2: toric surfaces, unequal residues]\tP\ \tV\label{p2:thm:C2}
Let $X$ be a complete smooth toric surface, $A$ a non-degenerate skew $2\times2$ matrix ($A_{12}\ne0$), and all $r_i\ne1$. Then the Poincar\'e polynomial of $\HH^*(L)$ is
\[P_{\HH}(t)=\chi(X)\,t^2+[\rho\in M]\,(1+t)^2\sum_q\tilde h^{q-1}(V_S)\,t^q,\qquad S=\{i:r_i\ge2\}.\]
In particular for $S=\emptyset$ (all $r_i\le0$) it equals $(1+t)^2+\chi(X)t^2$, for $S$ = all rays $t^2(1+t)^2+\chi(X)t^2$, and for proper subsets $S$ it is determined by the shape of $V_S$.
\end{thm}
\begin{hyptab}
Toric surface, $\sigma$ log-symplectic & Hypotheses of Part I, Thm.\ 4.1 (surface main theorem): $D$ reduced with only ordinary double points (the torus-fixed points), $\mu=\chi(X)$ of them. & Proof: $\HH^*=H^*(X;F)\oplus\C^{\mu}[-2]$.\\\midrule
$r_i\ne1$ (all components) & Condition of Part I, Lemma I.4.3 (acyclicity of the quotient complex $Q$). & Same.\\\midrule
$\mathrm{DR}_{\log}\simeq F$ & Per component: $Rj_*$ if $r_i\in\Z_{\le0}$, otherwise extension by zero; mixed case of Part I, Lemma I.4.2. & Reduces to $H^*(X;F)=H^*(X;\mathrm{DR}_{\log})$, then apply Theorem C1.
\end{hyptab}
\begin{proof}
By Part I, Thm.\ 4.1, $\HH^*(L)=H^*(X;F)\oplus\C^{\mu}[-2]$, and on a toric surface $\mu=$ number of maximal cones $=\chi(X)$. By Part I, Lemma I.4.2, $\mathrm{DR}_{\log}\simeq F$, so $H^*(X;F)=H^*(X;\mathrm{DR}_{\log})$. By Theorem C1, if $\rho\in M$ (all $r_i\in\Z$) this is $(1+t)^2\sum\tilde h^{q-1}(V_S)t^q$ with $S=\{r_i\ge1\}=\{r_i\ge2\}$ (because $r_i\ne1$).\qedhere
\end{proof}
\tV\ $\mathbb{CP}^2$ (triangle, $m=-9..6$, residues $\{-2..3,\pm\tfrac12\}$): 755 rows; $\mathbb{CP}^1\times\mathbb{CP}^1$ (4 lines): 729 rows, by exact linear algebra (weight-decomposed Cech hypercohomology; \texttt{cp2/roundC\_toric\_cech.py}, \texttt{roundC\_p2\_sweep*.py}, \texttt{roundC\_p1p1\_sweepA.py}), with 0 discrepancies. In particular on $\mathbb{CP}^1\times\mathbb{CP}^1$ with $S=$ two opposite sides ($V_S$ two points, $\tilde h^0=1$) the term $t(1+t)^2$ predicted by the formula does occur (confirmation of a new prediction).
\begin{rem}[\tP\ Consistency with Part I, Thm.\ 5.1]
For equal residues $r_i=r$, $S=\emptyset$ ($r\le0$) or $S=$ all rays ($r\ge2$), so Theorem C2 agrees with the three cases of Part I, Thm.\ 5.1 (e.g.\ $\mathbb{CP}^2$ triangle, $r\le0$: $(1+t)^2+3t^2=[1,2,4]$). Only unequal residues make $S$ a proper subset and produce the new terms $\tilde h^{q-1}(V_S)$.
\end{rem}

\section{Toric varieties of any even dimension: generic $\sigma$}\label{p2:sec:toricn}
Let $n\ge2$ be even. Part I, Thm.\ 6.4 (weight formula) \tP\ \tV: if $E_1^{a,b}=H^a(X,\wedge^bT\otimes L)=0\ (a>0)$ then
\[P_{\HH}(t)=\sum_{\chi\in\Lambda}t^{s(\chi)}(1+t)^{n-s(\chi)},\quad\Lambda=\{\chi\in M:\ c_i(\chi)\ge-1\ \forall i,\ A(\chi+\rho)\in W(\chi)\}.\]

\begin{thm}[Theorem T]\tP\ \tV\label{p2:thm:T}
Let $X$ be a smooth projective toric variety ($\dim X=n$ even), $\sigma$ a torus-invariant log-symplectic structure (constant matrix $A$), $L=\mathcal O(\sum a_iD_i)$, $\nabla=d+\rho$, and assume
\begin{itemize}
\item[(V)] $H^a(X,\wedge^bT\otimes L)=0$ ($a>0$, all $b$).
\item[(A)] $L(D):=L\otimes\mathcal O(\sum D_i)$ is ample.
\item[(GEN)] $\mathrm{Pf}(A)\ne0$ and the numbers $A_{ij}\ (i<j)$ are linearly independent over $K=\mathbb Q(\rho_1,\dots,\rho_n)$.
\item[(R)] $r_i\ne1$ for all $i$.
\end{itemize}
Then
\[P_{\HH}(t)=[\ \rho\in M\ \text{and}\ r_i\le0\ (\forall i)\ ]\,(1+t)^n+\chi(X)\,t^n.\]
Equivalently: if $L$ has no positive integral residue, $\HH^*(L)\cong H^*(U;L_U)\oplus\C^{\chi(X)}[-n]$ (for $\rho\notin M$, $L_U$ is a non-trivial local system on $(\C^*)^n$ and $H^*(U;L_U)=0$).
\end{thm}
\begin{hyptab}
(V) & Hypercohomology equals cohomology of the complex of global sections. & Prerequisite for the weight formula (Part I, Thm.\ 6.4).\\\midrule
(A) $L(D)$ ample & The polytope $\{\chi:c_i(\chi)\ge-1\}$ is simple with normal fan $\Sigma$ and has exactly $\chi(X)$ vertices. & Case $s=n$: $\chi(X)$ vertex weights, each giving $t^n$.\\\midrule
(GEN) $A_{ij}$ independent over $K$ & The skew condition $A(\chi+\rho)\in W$ degenerates over $K$ to $x\parallel m$ ($x=\chi+\rho$). Excludes arithmetic coincidences. & Case $0<s<n$: kills middle-stratum contributions.\\\midrule
(R) $r_i\ne1$ & For $r_i=1$ solutions with $x=0$ or $x\parallel m_0$ appear (resonance, \S\ref{p2:sec:res}). & Cases $s=n-1$ and $\chi=-\rho$: kills resonant contributions.\\\midrule
$\mathrm{Pf}(A)\ne0$ & $\sigma$ is log-symplectic ($\sigma^\sh:B\cong T_D$). & The log de Rham differential $(\chi+\rho)\wedge$ corresponds to $v_\chi\wedge$ ($v_\chi=A(\chi+\rho)$).
\end{hyptab}
\begin{proof}
Put $x=\chi+\rho$, $C=C(\chi)$, $W=W(\chi)$, $s=\dim W$. For $i\in C$, $\langle x,u_i\rangle=(-a_i-1)+\langle\rho,u_i\rangle=r_i-1$. The condition $Ax\in W$ is equivalent to $\langle m,Ax\rangle=\sum_{i<j}A_{ij}(m_ix_j-m_jx_i)=0$ for all $m\in\mathrm{Ann}_M(W)$.

\emph{$s=n$:} $W=N_\C$ and the condition is empty. By (A), characters with $c_i\ge-1$ and $C(\chi)$ spanning everything are exactly the vertices of the polytope $\{c_i\ge-1\}$ (simple, normal fan $\Sigma$, $X$ smooth so vertices are lattice points), one for each maximal cone: $\chi(X)$ in total, each contributing $t^n$.

\emph{$s\le n-2$:} $\mathrm{Ann}_M(W)$ has rank $n-s\ge2$, so there are linearly independent integer vectors $m,m'$. The coefficients $m_ix_j-m_jx_i$ lie in $K$, and by (GEN) $\sum A_{ij}(\cdots)=0$ forces all coefficients to vanish, i.e.\ $x\parallel m$. Likewise $x\parallel m'$. By independence $x=0$, $\chi=-\rho$. But $x=0$ gives $r_i-1=0$ for $i\in C$, contradicting (R) when $C\ne\emptyset$. So there are no solutions with $s\ge1$ in this range; only $C=\emptyset$ ($s=0$) remains.

\emph{$s=n-1$:} $\mathrm{Ann}_M(W)=\Z m_0$. Similarly $x=\tau m_0$, and for $i\in C$, $\langle x,u_i\rangle=\tau\langle m_0,u_i\rangle=0=r_i-1$, contradicting (R); no solutions.

\emph{$s=0$:} $W=0$ and the condition is $x=0$, i.e.\ $\chi=-\rho$ (needs $\rho\in M$). Then $c_i(\chi)=-r_i\ge-1$ ($r_i\le1$) and $C(\chi)=\emptyset$ ($r_i\ne1$), so $r_i\le0$. Contribution $(1+t)^n$.

Adding up gives the claim. The second statement follows because for $\rho\in M$, $L_U$ is trivial, and under (R) ``all $r_i\le0$'' is equivalent to ``no positive integral residue''.\qedhere
\end{proof}
\begin{rem}[On (V)\ \tC]
(V) holds for $L=L_0^m$ ($L_0$ an ample Poisson module, $m\gg0$). The exact sequence of $T_X(-\log D)=\mathcal O\otimes N$ gives a filtration of $\wedge^bT$ with graded pieces $\wedge^{b-|J|}N\otimes\mathcal O_{D_J}(D_J)$ ($|J|\le b$); apply Demazure vanishing for nef line bundles (CLS, Thm.\ 9.2.3\tQ) to $L(D_J)|_{D_J}$ (ample for $m\gg0$). It is also a special case of Part I, Thm.\ 6.1 (Serre type).
\end{rem}
\tV\ Against the closed weight criterion: $\mathbb{P}^2$ ($m=1,2,3$), $\mathbb{P}^1\times\mathbb{P}^1$ ($\mathcal O(1,1),\mathcal O(1,2),\mathcal O(2,3)$), $\mathbb{P}^4$ ($m=1$), $\mathbb{P}^2\times\mathbb{P}^2$ ($(1,1)$), $(\mathbb{P}^1)^4$ ($(1,1,1,1)$), $\rho$ lattice points, rational points, $\rho=(1/3,0,\dots)$, $A$ random integers (modulo the prime $2^{61}-1$), 246 examples with all $r_i\ne1$ agree with the formula of Theorem T (0 discrepancies; \texttt{cp2/roundD\_check.py}). The broader 3626 examples (including resonances) also agree with the weight-formula prediction (Part I, Obs.\ 6.7).

\begin{obs}[Necessity of (GEN)\ \tV]\label{p2:obs:gen}
For $\mathbb{P}^2\times\mathbb{P}^2$ with block-diagonal $\sigma$ one gets $[1,4,12,16,16]$, for generic $A$ $[1,4,6,4,10]$. The difference $6t^2(1+t)^2$ comes from six codimension-2 strata $\{\text{double point of one factor}\}\times(\C^*)^2$, contributing $t^2H^*(O_\tau)$, and occurs only when the condition $x_C(A,r)\in\rho+M$ holds arithmetically (Part I, Obs.\ 6.8). (GEN) excludes this coincidence.
\end{obs}

\section{The resonant case $r_i=1$}\label{p2:sec:res}
If some $r_i=1$, the surface main theorem of Part I does not apply (the quotient differential $-u(1-r)p$ vanishes). The following formula is computational.

\begin{obs}[Single and double resonance ($\mathbb{CP}^2$, $\mathbb{CP}^1\times\mathbb{CP}^1$)\ \tV]\label{p2:obs:res}
Let $R$ be the set of resonant branches ($r_i=1$). $N=N_{D_i}\otimes L$ is the normal bundle of the resonant branches tensored with $L$, $K_R=(K_X^{-1}\otimes L)|_{D_R}$, and $\nu$ is the number of double points not on a resonant branch (the exact definition of $N$ for several resonant branches is transcribed from the original record \texttt{roundC\_notes.md} and unverified here\tQ). Then
\[P_{\HH}(t)=t\,h^0(N)+t^2\bigl(h^1(N)+h^0(K_R)+\nu\bigr)+t^3\,h^1(K_R).\]
No discrepancies for single and double resonance on $\mathbb{CP}^2$ and for $\mathbb{CP}^1\times\mathbb{CP}^1$ (108 examples with non-integral residues). In toric dimension $\ge2$ the corresponding statement is: for a cone $C$ with $s=n-1$ and $r_i=1$ ($i\in C$) the condition holds automatically and $\deg(L|_C)+1$ classes $t^{n-1}(1+t)$ appear, one for each lattice point on the toric curve $C$ (Part I, Cor.\ 6.5(c)).
\end{obs}
\begin{obs}[Where the formula fails\ \tV]
When all residues are integral ($L_U$ trivial) with $m\ll0$, and when all branches are resonant ($L=K_X$), the connecting map of $H_c$-type in $\mathrm{DR}_{\log}$ acts and the formula fails. Example: $\mathbb{CP}^2$, double resonance, $m=-5$ (integral residues): observed $[0,0,5,2]$, formula predicts $[0,0,6,3]$. Triple resonance ($L=K$): observed $[0,0,3]$, formula predicts $4t^2+t^3$. The case $L=K_X$ is treated separately in Part I, Prop.\ 4.7 ($E=K_X$).
\end{obs}
\begin{conj}\tK
The resonance formula (Observation \ref{p2:obs:res}) holds in general when $L_U$ is non-trivial and not all branches are resonant. A proof needs the computation of the connecting map of $\mathrm{DR}_{\log}\to C(A;L)\to Q$ (an extension of Part I, Lemma I.4.5); unfinished.
\end{conj}

\section{Local formulas at non-nodal singularities}\label{p2:sec:local2}
\subsection{Quasi-homogeneous plane curve singularities}
Local model and notation follow Part I, \S I.6.6: $\sigma=h\,\partial_z\wedge\partial_w$, module $\mathcal O\,h^s$, $d_0f=(hf_w+sfh_w)\partial_z-(hf_z+sfh_z)\partial_w$, $d_1(p,q)=h(p_z+q_w)-(1-s)(h_zp+h_wq)$. The cohomology dimensions of the stalk $\mathcal C_s$ are written $(h^0,h^1,h^2)$.

\begin{obs}[Local formula (weight-corrected)\ \tV\ \tK]\label{p2:obs:stalk}
Let $h$ be a quasi-homogeneous isolated singularity (spectrum $\mathrm{Sp}$, minimum $\mathrm{lct}$, Milnor number $\mu$) and $s\ne1$. Then (dimension-wise)
\[\text{stalk}(\mathcal C_s)=\C^{\mu}[-2]\ \oplus\ R(s),\]
\[\begin{split}R(s)={}&[\,s\in\Z,\ s\le0\,](1+t)\\&+(t+t^2)\cdot\#\{\alpha\in\mathrm{Sp}:\ s+\alpha\in\Z,\ s+\alpha\le1\}\\&\hspace{2em}(\text{in the sense of Poincar\'e polynomials}).\end{split}\]
$s=1$ is excluded (infinite dimensional stalk).
\end{obs}
\paragraph{Status of the proof (clarified in v2).}
\begin{itemize}
\item \tP\ For $s\ne1$ the stalk $\mathcal C_s$ splits into the subcomplex $B_s=(\mathcal O\to\Theta(-\log D)\to h\mathcal O)$ and the quotient $Q=\mathcal C_s/B_s\simeq[J\!\cdot\!\mathcal O_D\to\mathcal O_D]$ (injective differential, $\dim\mathcal O/J=\mu$), so stalk $=H^*(B_s)\oplus(\mathcal O/J)[-2]$ (dimension-wise). By a Cartan homotopy for the Euler field, $H^*(B_s)$ is concentrated in the single weight $-s\,d$ ($d$ the degree) and vanishes for $s>1-\mathrm{lct}$. However, the sign convention for $s$ in the identification $B_s\cong(\Omega^\bullet(\log D),\,d+s\,d\log h)$ was fixed by matching the computed data and has not been checked independently.
\item \tV\ The formula above has 0 discrepancies for all lattice values of $s$ in 10 types: $A_1$, $A_2$, $A_3$, $D_4$, $D_5$, $E_6$, $T_{2,3,6}$ ($z^3-w^6$), $z^4-w^4$, $z^4-w^5$, $z^5-w^5$ (scripts: \texttt{cp2/roundF\_local\_rule2.py}, \texttt{roundF\_local\_lowlct.py}).
\item \tK\ Not proved: for $s\le1-\mathrm{lct}$, that the weight $-s\,d$ part of $H^*(B_s)$ has the dimension predicted by the Wang sequence of the Milnor fibration (the $R(s)$ above). Interpretation: the Wang sequence gives the stalk of $Rj_*L_s$, and the logarithmic complex $B_s$ computes its part of weight $\ge-\delta$ ($\delta=d-w_z-w_w$). Related: Calder\'on-Moreno--Narv\'aez-Macarro \cite{CMNM} prove logarithmic comparison for $E(kD)$ ($k\gg0$, $Rj_*$) and $E(-kD)$ ($j_!$) for free divisors whose Jacobian ideal is of linear type (journal pages \tQ).
\end{itemize}
\paragraph{Relation to the v1 rule.}The v1 rule $[\,s<\mathrm{lct}\,]\bigl(\cdots\bigr)$ is wrong for $\mathrm{lct}<1/2$ (example: $z^4-w^5$, $\mathrm{lct}=9/20$: for $s=11/20=1-\mathrm{lct}$ we have $s>\mathrm{lct}$ but a class exists). For singularities with $\mathrm{lct}\ge1/2$ the old and new rules agree; the 12 types treated in v1 ($A_1$--$A_6$, $D_4$--$D_6$, $E_6$, $E_8$, $T_{2,3,6}$) all have $\mathrm{lct}\ge1/2$, so the values in v1's table are unaffected (checked against the new rule \tV). For $\mathrm{lct}<1/2$ (e.g.\ ordinary points of multiplicity $\ge4$) the v1 rule cannot be used.

\begin{rem}[Correction: weights of $D_5$\ \tP]\label{p2:rem:D5}
$D_5=z^2w-w^4$ has weights $(3,2)$, degree $8$, spectrum $\{5/8,7/8,1,9/8,11/8\}$, $\mathrm{lct}=5/8$, $\mu=5$. The weights $(3,1,7)$ in an older version of the internal note \texttt{cp2\_second\_notes.md} were wrong, and the $D_5$ values derived from them are invalid. \emph{The published Part I (v1, v2) writes the $D_5$ table and text with the correct weights and spectrum ($\mathrm{lct}=5/8$) and is unaffected} (\S\ref{p2:sec:chk}). The values of the formula above for $D_5$ are $(0,0,5)$ (generic $s$), $(1,2,6)$ ($s\in\Z_{\le0}$), $(0,1,6)$ ($s\equiv$ a spectral number mod $\Z$, $s<\mathrm{lct}$), matching Part I's table.
\end{rem}

\subsection{Triple points with unequal residues}
\begin{obs}[\tV]\label{p2:obs:triple}
$h=zw(z-w)$ ($D_4$), module $\prod h_i^{s_i}$ (different $s_i\ne1$ on each branch). The stalk is
\[\text{stalk}=\C^4[-2]\ \oplus\ R,\qquad R=\begin{cases}H^*(U_{\rm loc};L_s)&(s_1+s_2+s_3<2),\\0&(\text{otherwise}).\end{cases}\]
$U_{\rm loc}=\C^2\setminus\{3\text{ lines}\}\simeq\C^*\times(\mathbb P^1\setminus3\text{ pts})$. Concretely: all $s_i\in\Z_{\le0}$ gives stalk $(1,3,6)$; $\sum s_i\in\Z$ ($\le1$) with not all $s_i$ integers gives $(0,1,5)$; otherwise $(0,0,4)$. For 533 random triples $s_i\in\tfrac16\Z\cap[-2,2]\setminus\{1\}$: 0 discrepancies.
\end{obs}
At a node the per-branch threshold $s_i<1$ decides whether an $Rj_*$ part is present, but at a triple point it is the \emph{sum of the residues} $s_1+s_2+s_3$ against the threshold $2=w_z+w_w$ (weights for the single Euler field). \tK\ General principle: the threshold is decided by comparing the weight of the module under the symmetry torus of the germ with the log canonical weight (equal residues: $s<\mathrm{lct}$). The anomaly of Part I, \S I.6.6 for triple points ($s=-m/3=2/3,4/3$) is explained by the sum hitting $2,4$ (sum $=3s=2,4$). For equal residues $s_1=s_2=s_3=s$ this rule agrees with the new formula of Observation \ref{p2:obs:stalk} for $D_4$ ($\mathrm{Sp}=\{2/3,1,1,4/3\}$, $\mathrm{lct}=2/3$) for all $s\in\tfrac13\Z$ with $-10\le3s\le10$ (cross-check \tV).

\section{A conjecture in general dimension and an index-theoretic outlook}\label{p2:sec:conj}
\begin{conj}[General SNC log-symplectic varieties]\tK\label{p2:conj:gen}
Let $(X^{2k},D)$ be a smooth projective variety with a simple normal crossings divisor, $\sigma$ a generic log-symplectic structure, $L$ a Poisson module with constant residues $r_i$, all $r_i\ne1$, satisfying (V). With $D_{(n)}$ the set of points where $n$ components meet,
\[\HH^*(L)\cong H^*(X;F)\oplus\bigoplus_{p\in D_{(n)}}\C[-n],\qquad F=j_{2!}Rj_{1*}L_U.\]
If $L$ has no positive integral residue then $H^*(X;F)=H^*(U;L_U)$.
\end{conj}
Consistency with the Euler characteristic (Theorem E) amounts to $\chi(U)+\#D_{(n)}=\chi(X)$, i.e.\ $\sum_{0<k<n}\chi(D_k^\circ)=0$ ($D_k^\circ$ = the open stratum where exactly $k$ components meet). This holds for surfaces (adjunction) and toric varieties (strata are tori)\tP; in general dimension there is only a sketch from the existence of non-vanishing log vector fields $\sigma^\sh(\dlog z_i)$ tangent to the strata\tK. In the toric case Theorem T gives a proof.

\begin{obs}[Index-theoretic outlook\ \tK]
(1) Volume term: the integral of the Euler class of the log Lie algebroid $T_X(-\log D)$, $\int c_n(T_X(-\log D))=\chi(U)$. (2) Corner term: $\#D_{(n)}\,t^n$. (3) Residue-dependent defect: the difference between $\mathrm{DR}_{\log}(L)$ and $Rj_*L$ or $j_!L$ is cohomology of the boundary $\partial U$ (Euler characteristic $0$), and $r\in\Z$ corresponds to indicial roots of the weighted $b$-calculus. A proof as an index theorem needs index theory for the complex Lie algebroid $L$ (unimodularity, Fredholmness, APS-type boundary terms), which has not been found (Part I, \S I.5.3). Hu's Lie groupoid Riemann--Roch does not apply to this complex (Part I, \S I.7).
\end{obs}

\section{Impact on Part I (v1, v2)}\label{p2:sec:chk}
We checked whether the results of Rounds C and D invalidate anything in the published Part I (v1, v2). Conclusion: no theorem or proposition has to be withdrawn.

\begin{itemize}
\item \textbf{Weights of $D_5$:} the error was only in the internal note (old \texttt{cp2\_second\_notes.md}); Part I uses $D_5=w(z^2-w^3)$ with $\mathrm{lct}=5/8$ (corresponding to the correct weights $(3,2)$). The $D_5$ values in the table of Part I, \S I.6.6 agree with the prediction of Observation \ref{p2:obs:stalk}. No impact.\tV
\item \textbf{Wording of Part I, Obs.\ 6.7:} ``$-\rho$ is included when $\rho\in M$ and all $r_i\le1$'' and ``the $U$-part is $(1+t)^n$ when all $r_i\le0$'' appear side by side; the latter is the form after the resonant ($r_i=1$) contribution has been counted separately. The wording should be made explicit (not a mathematical error). Theorem T above states (R) explicitly.
\item \textbf{Equal residues, Part I, Thm.\ 5.1:} Theorem C2 extends it to unequal residues and agrees in the equal case (remark after Theorem C2).
\item \textbf{Resonant and integral residues:} Part I assumes $r_i\ne1$ for its main theorem; the failure cases of \S\ref{p2:sec:res} lie outside the scope of Part I's claims.
\item \textbf{Triple points:} ``the anomaly at $s=2/3,4/3$'' of Part I, \S I.6.6 can be rephrased by the sum rule (Observation \ref{p2:obs:triple}); the conclusion is unchanged.
\end{itemize}
We therefore judged a v3 publication unnecessary. The small wording clarification (second item) will be made at the next revision of Part I. (Internal record: \texttt{v3\_corrections.md}.)

\section{Status of claims and unverified items}\label{p2:sec:ledger}
\begin{center}\footnotesize
\begin{tabular}{>{\raggedright\arraybackslash}p{4.0cm}>{\raggedright\arraybackslash}p{2.6cm}>{\raggedright\arraybackslash}p{8.8cm}}\toprule
Claim & Status & Dependencies / gaps\\\midrule
Theorem E (Euler characteristic) & \tP & HRR only; uses only coherence and compactness.\\
Theorem C1 ($\mathrm{DR}_{\log}$, toric) & \tP\ \tC & CLS Thm.\ 9.1.3 (number \tQ). K\"unneth is standard.\\
Theorem C2 (toric surfaces, unequal residues) & \tP\ \tV & Part I Thm.\ 4.1, Lemma 4.2 and Theorem C1. 1484 rows, 0 discrepancies.\\
Theorem T (toric, generic $\sigma$) & \tP\ \tV & Part I Thm.\ 6.4 (weight formula) and (V). 246 examples, 0 discrepancies. (V) holds for $L_0^m$ ($m\gg0$) by Demazure vanishing (number \tQ).\\
Obs.\ \ref{p2:obs:gen} (necessity of (GEN)) & \tV & A few examples; mechanism only sketched.\\
Resonance formula (Obs.\ \ref{p2:obs:res}) & \tV\ \tK & Transcription of $N$ (\tQ). Fails in some cases. Proof unfinished.\\
Local formula (Obs.\ \ref{p2:obs:stalk}, weight-corrected) & \tP\ \tV\ \tK & Splitting $B_s/Q$ and weight argument proved (sign convention of $s$ fixed by data). Computed in 10 types. Agreement with the Wang-sequence dimension not proved.\\
Triple point rule (Obs.\ \ref{p2:obs:triple}) & \tV\ \tK & 533 triples; general principle conjectural.\\
Conjecture \ref{p2:conj:gen}, index outlook & \tK & Neither definitions nor proofs complete.\\
Generalized Kodaira embedding & --- & Not treated.\\\bottomrule
\end{tabular}
\end{center}

\paragraph{Unverified / caution items.}
\begin{itemize}
\item The theorem numbers of Cox--Little--Schenck (Thm.\ 9.1.3, 9.2.3) are from memory.\tQ
\item The exact definition of $N$ in Observation \ref{p2:obs:res}, and the tables of 729 and 755 rows in Theorem C2, are transcribed from the internal record (\texttt{roundC\_notes.md}) and were not rerun while writing this part.\tQ
\item The spectral numbers of ADE types are standard but no source is given here; see e.g.\ Arnold--Gusein-Zade--Varchenko, \emph{Singularities of Differentiable Maps II}, Birkh\"auser 1988.\tQ
\item The local formula covers quasi-homogeneous germs only (rule corrected in v2). For non-reduced $D$ (double lines, double conics) stalks are infinite dimensional and global dimensions grow linearly in $m$ (Part I, \S I.6.6).
\item The items marked ``to be verified'' in Part I are inherited.
\end{itemize}

\paragraph{Work record.} Scripts use only the Python standard library with exact rational / modulo-$(2^{61}-1)$ arithmetic (\texttt{cp2/roundC\_*.py, roundD\_check.py, toric\_*.py, local\_*.py}). Research notes are in \texttt{roundC\_notes.md}, \texttt{roundD\_notes.md}, \texttt{higher\_dim\_notes.md}.




\clearpage
\part{Kodaira-type theorems}
\setcounter{section}{0}\renewcommand{\thesection}{III.\arabic{section}}
\noindent\textit{Part III of the series; originally published as a separate note (\url{https://hylt-math.github.io/2026/11/22/generalized-geometry-poisson-part3/}). Within this paper, Part-internal references use the numbering III.$\ast$.$\ast$.}\par
%--------------------------------------------------------------
\section{Position of this part and overview}\label{p3:sec:intro}
%--------------------------------------------------------------
This part continues Part~I (Claude, Mori, ``Cohomology of generalized holomorphic line bundles on generalized complex structures from holomorphic Poisson structures'', v2, \url{https://hylt-math.github.io/2026/11/22/generalized-geometry-poisson/}) and Part~II (``\dots Riemann--Roch-type formulas (Part II)'', \url{https://hylt-math.github.io/2026/11/22/generalized-geometry-poisson-part2/}). The setting (the generalized complex structure $\Jc_\sigma$ obtained from a holomorphic Poisson structure $\sigma$ by a $\beta$-transform, Poisson module $=$ generalized holomorphic line bundle, $\HH^*(L)=H^*(X;C(A;L))$) is not repeated; notation follows Parts I and II.

\paragraph{Goal and conclusion.}
Of the three treasures of complex geometry, Riemann--Roch and Serre vanishing were treated in Parts I and II. What remains is Kodaira's embedding theorem. In short:
\begin{itemize}
\item The naive formulation, ``embed $M$ in projective space by generalized holomorphic sections'', is false (\S\ref{p3:sec:fail}) \tP.
\item In regular type with projective leaf space, the map is the composition of the projection with the classical Kodaira embedding of the leaf space (Theorem~\ref{p3:thm:K1}, \S\ref{p3:sec:leaf}) \tP. This is a direct reduction; we did not find it in the literature.
\item As a replacement we propose ``Poisson--Kodaira'', proved in the toric case and verified numerically on $\mathbb P^2$ (\S\ref{p3:sec:PK}); in general it is conjectural \tK.
\item No Kodaira-type embedding theorem for general generalized complex structures was found in the literature we examined \tC.
\end{itemize}
We also organize results on non-reduced boundary, the local formula for quasi-homogeneous singularities (including a correction to Part II), and the K\"unneth formula for products (\S\ref{p3:sec:nonred}--\S\ref{p3:sec:prod}).

%--------------------------------------------------------------
\section{Notation (new in this part)}\label{p3:sec:notation}
%--------------------------------------------------------------
\begin{longtable}{>{\raggedright\arraybackslash}p{3.8cm}>{\raggedright\arraybackslash}p{11.2cm}}
\toprule Symbol & Meaning\\ \midrule\endhead
$(X,\sigma)$ & compact complex manifold with holomorphic Poisson structure ($[\sigma,\sigma]=0$)\\
$A=(T^*X)_\sigma$ & holomorphic Lie algebroid defined by $\sigma^\sharp$\\
$L$, $D_\alpha s$ & Poisson module $=$ generalized holomorphic line bundle: $D_\alpha(fs)=fD_\alpha s+\sigma(\alpha,df)s$ ($\alpha\in\Omega^1$), flat\\
GH section & an element of $\HH^0(L)=\{s\in H^0(X,L):D_\alpha s=0\ \forall\alpha\}$ (kernel of $d_1:H^0(L)\to H^0(T\otimes L)$ by the Dolbeault model of Part~I)\\
$\pi:M\to X$ & in regular type, projection to the symplectic leaves (leaves $=$ fibres, $X$ the leaf space, $\dim_\C X=k$)\\
SGH bundle & strong generalized holomorphic bundle (flat along leaves, holomorphic transversally); Pal--Poddar Def.~3.1 \tC\\
$R=\bigoplus_k H^0(X,L^k)$ & graded ring (\S\ref{p3:sec:PK})\\
$\Pi_F$ & quadratic Poisson structure on $\C^3$, $\Pi_F^{ij}=\epsilon_{ijk}\partial_kF$ (Part~I \S4)\\
$\nu:X\hookrightarrow\mathbb P^N$ & Kodaira embedding by $|L^m|$\\
$\mathcal J$ (Jacobiator) & $\{a,\{b,c\}\}+(\text{cyclic})$\\
$k_i$, $N_{D_i}$ & multiplicity and normal bundle of a component of a non-reduced divisor $D=\sum k_iD_i$\\
$U=X\setminus D_{\mathrm{red}}$, $L_U$ & complement and the rank-one local system determined by the residues\\
$P_{\HH}(t)$, $\mu$ & $\sum\dim\HH^k\,t^k$, Milnor number\\
\bottomrule
\end{longtable}
\noindent Tags: \tC\ literature, \tP\ proved by the author (standard facts cited), \tV\ verified by computer, \tK\ conjecture/outlook, \tQ\ unverified.

%--------------------------------------------------------------
\section{The classical Kodaira embedding theorem: statement, hypotheses, roles}\label{p3:sec:classical}
%--------------------------------------------------------------
\begin{thm}[Kodaira's embedding theorem]\tC\label{p3:thm:Kod}
Let $X$ be a compact complex manifold and $L$ a positive holomorphic line bundle. Then $X$ is projective, and for $m\gg0$ the map $X\to\mathbb P(H^0(X,L^m)^*)$ defined by $|L^m|$ is an embedding.
\src{Kodaira, Ann.\ Math.\ 60 (1954) 28--48. Textbooks: Griffiths--Harris, \emph{Principles of Algebraic Geometry}, Ch.~1 \S4; Huybrechts, \emph{Complex Geometry}, Prop.~5.3.1 (numbering \tQ).}
\end{thm}
\begin{hyptab}
$X$ compact & $H^0(X,L^m)$ is finite-dimensional and the map to $\mathbb P^N$ is defined. & Start and end of the proof (finite-dimensionality, closed embedding).\\\midrule
$L$ positive ($i\Theta_L>0$) & The curvature form is a K\"ahler form, so $X$ is K\"ahler and Kodaira vanishing applies. & The step showing $H^1(\tilde X,\pi^*L^m\otimes\mathcal O(-2E_x))=0$ on the blow-up $\tilde X$ at a point $x$ (Kodaira vanishing).\\\midrule
Enough holomorphic sections & Sections separate points and tangent vectors. & Surjectivity of $H^0(L^m)\to L^m\otimes\mathcal O/\mathfrak m_x^2$ follows from that vanishing; no Riemann--Roch dimension count is needed.
\end{hyptab}
\noindent What matters is the step where sections \emph{separate points}. As we show next, this step breaks down for generalized holomorphic (GH) sections.

%--------------------------------------------------------------
\section{No embedding by generalized holomorphic sections}\label{p3:sec:fail}
%--------------------------------------------------------------
\subsection{The non-degenerate case}
\begin{prop}\tP\label{p3:prop:one}
Let $X$ be a connected compact complex manifold, $\sigma$ a holomorphic Poisson structure non-degenerate on a dense open set $U_0$, and $L$ a Poisson module. Then $\dim_\C\HH^0(L)\le1$. A GH section is determined by its value at one point of $U_0$, so the Kodaira map by GH sections is constant (or undefined).
\end{prop}
\begin{hyptab}
Connected, $U_0$ dense & The identity theorem applies. & Step (iii).\\\midrule
$\sigma$ non-degenerate on $U_0$ & $\sigma^\sharp:A|_{U_0}\to T$ is an isomorphism, so $D$ becomes an ordinary flat connection. & Step (i).\\\midrule
$L$ a line bundle & A flat section is determined by its value at one point. & Step (ii).
\end{hyptab}
\begin{proof}
(i) On $U_0$ put $\nabla_vs:=D_{(\sigma^\sharp)^{-1}v}s$; by the Leibniz rule of $D$ this is an ordinary flat holomorphic connection, and GH sections satisfy $\nabla s=0$. (ii) $U_0$ is connected (a connected manifold minus a proper analytic subset) and $L$ has rank one, so the space of flat sections has dimension $\le1$ and a section is determined by one value. (iii) $s$ is determined by $s|_{U_0}$ (identity theorem).\qedhere
\end{proof}
\begin{ex}\tV\label{p3:ex:casimir}
On $\mathbb{CP}^2$ with $\Pi_F$ ($F$ a smooth cubic): $\HH^0(\mathcal O(m))=\C F^{k}$ for $m=3k\ge0$ and $0$ otherwise. On $\mathbb P^1\times\mathbb P^1$ with a smooth $(2,2)$-curve: $\HH^0(\mathcal O(m,m))=\C F^{m/2}$ for $m\in2\Z_{\ge0}$ (tables of Part~I). So for no $m$ is there an embedding by GH sections.
\end{ex}

\subsection{Regular type}
\begin{prop}\tP\label{p3:prop:leaf}
Let $M$ be a compact generalized complex manifold, regular of type $k<n$ on a dense open set $W$ (leaf dimension $2(n-k)>0$), and $L$ a generalized holomorphic line bundle. For any GH sections $s_0,\dots,s_N$, the map $\varphi=[s_0:\cdots:s_N]$ (defined off the base locus) is constant on each leaf in $W$; in particular it is not an immersion on $W$.
\end{prop}
\begin{hyptab}
Regular type (product local model) & $\pi_T(L)\supset T_F\otimes\C$ is available. & Local normal form (Darboux, Gualtieri 2011, Thm.~4.3 \tC).\\\midrule
$k<n$ & Leaves have positive dimension. & One can move along a leaf.\\\midrule
$L$ generalized holomorphic & The $L$-connection satisfies $D_\xi(fs)=\pi_T(\xi)(f)s+fD_\xi s$. & GH section $\Rightarrow$ parallel along leaves.
\end{hyptab}
\begin{proof}
In the product model $L=\mathrm{graph}(\omega_Y)\oplus(T^{0,1}\oplus T^{*1,0})$ and $\pi_T(L)=T_Y\otimes\C\oplus T^{0,1}$. For $\xi\in L$ with a leafwise component, $D_\xi s=0$ says that $s$ is parallel for the leafwise flat connection $\nabla^F$. For a path in a leaf and parallel transport $P:L_x\to L_y$, every GH section satisfies $s_i(y)=Ps_i(x)$, hence the ratios $[s_0:\cdots:s_N]$ at $x$ and $y$ agree.\qedhere
\end{proof}
\begin{ex}\tP\label{p3:ex:prodY}
For $M=Y\times X$ ($Y$ symplectic, $X$ complex) and $L=\mathrm{pr}_Y^*F\boxtimes\mathrm{pr}_X^*G$: $\HH^0(L)=H^0_{dR}(Y,F)\otimes H^0(X,G)$. Since $H^0_{dR}(Y,F)$ has dimension $\le1$ per component of $Y$, the map factors through $\mathrm{pr}_X$ (and the index of the component). This agrees with the K\"unneth computation of Part~I.
\end{ex}
\begin{rem}[Where the classical proof breaks]\tP
The classical proof uses that sections \emph{separate points and tangent vectors}. In Proposition~\ref{p3:prop:one} at most one section (up to scale) separates anything; in Proposition~\ref{p3:prop:leaf} points on one leaf cannot be separated. Positivity of curvature guarantees the \emph{existence} of sections but not that they are GH (flat). Pal--Poddar's ``positive SGH line bundle'' (Defs.~11.10, 11.11 \tC\ \tQ) asks for positivity only in the transversal direction, which corresponds exactly to this situation.
\end{rem}

%--------------------------------------------------------------
\section{The leaf-space Kodaira theorem}\label{p3:sec:leaf}
%--------------------------------------------------------------
\begin{thm}[Leaf-space Kodaira]\tP\label{p3:thm:K1}
Let $M$ be a compact generalized complex manifold of regular type such that:
\begin{enumerate}
\item[(K1)] the projection $\pi:M\to X$ to the symplectic leaves is a proper submersion (or an orbifold submersion) with connected compact fibres $(=$ leaves$)$, and $X$ is a compact complex manifold (orbifold) of dimension $k$;
\item[(K2)] $L$ is an SGH line bundle with $L\cong\pi^*L_X$ (as SGH bundles) for a holomorphic line bundle $L_X\to X$ (trivial leafwise holonomy);
\item[(K3)] $L_X$ is ample.
\end{enumerate}
Then there is $m_0$ such that for $m\ge m_0$: $\HH^0(M,L^m)=\pi^*H^0(X,L_X^m)$, and the map $\varphi_m:M\to\mathbb P^N$ by GH sections equals $\iota_m\circ\pi$, where $\iota_m:X\hookrightarrow\mathbb P^N$ is the classical Kodaira embedding. Hence each fibre of $\varphi_m$ is exactly one leaf. In the orbifold case $\iota_m$ embeds the coarse space ($m$ a multiple of the isotropy orders; Baily, Amer.\ J.\ Math.\ 79 (1957) 403--430; cited from memory \tQ).
\end{thm}
\begin{hyptab}
(K1) connected compact fibres & Leafwise flat sections are constant (converse of Prop.~\ref{p3:prop:leaf}); since fibres are connected, the preimage under $\varphi_m$ is exactly one leaf. & Steps (i), (iii).\\\midrule
(K2) trivial holonomy & $L^m$ is a trivial flat bundle on each leaf, so GH sections exist. With non-trivial holonomy, $\HH^0$ vanishes along the leaf. & Step (i).\\\midrule
(K3) $L_X$ ample & Classical Kodaira (Theorem~\ref{p3:thm:Kod}) is applied on $X$. & Step (ii).\\\midrule
Orbifold case & Kodaira for the coarse space (Baily). & Orbifold version of step (ii).
\end{hyptab}
\begin{proof}
(i) $\HH^0(M,L^m)$ consists of the sections of $L^m$ flat along leaves and holomorphic transversally (definition of SGH bundles, degree-0 part of Pal--Poddar Cor.~8.10 \tC\ \tQ). By (K1), (K2), on each fibre $L^m=\pi^*L_X^m$ is a trivial flat bundle, so flat sections are constant on fibres, i.e.\ of the form $\pi^*t$ with $t$ a section of $L_X^m$; transversal holomorphicity is equivalent to holomorphicity of $t$. Hence $\HH^0(M,L^m)=\pi^*H^0(X,L_X^m)$. (ii) By (K3) and classical Kodaira (Theorem~\ref{p3:thm:Kod}), $\iota_m$ is an embedding for $m\ge m_0$. (iii) $\varphi_m=\iota_m\circ\pi$ with $\iota_m$ injective, so $\varphi_m(x)=\varphi_m(y)\Leftrightarrow\pi(x)=\pi(y)$, i.e.\ $x,y$ lie on the same leaf.\qedhere
\end{proof}
\begin{rem}[Position]
\tP\ Theorem~\ref{p3:thm:K1} is a direct reduction, not a deep result. \tC\ We found no reference stating a Kodaira-type result in regular type. It is the embedding-side counterpart of Pal--Poddar's Thm.~11.12 (Kodaira--Nakano and Serre-type vanishing for positive SGH line bundles) \tC. \tQ\ Equivalence with their definition of ``positive'' is not checked. No Kodaira-type statement was found for the case where the type jumps or for general generalized complex structures.
\end{rem}

%--------------------------------------------------------------
\section{Poisson--Kodaira: embedding the Poisson pair}\label{p3:sec:PK}
%--------------------------------------------------------------
Since GH sections are (almost) constant, we change the question: for the Kodaira embedding $\nu:X\hookrightarrow\mathbb P^N$ by \emph{ordinary} holomorphic sections, is $\nu$ a Poisson embedding, i.e.\ is $\sigma$ the restriction of a Poisson structure $\sigma_N$ on $\mathbb P^N$? (This is the author's proposal, not a formulation from the literature.)

\begin{lem}[Rank-one Poisson modules and $\C^*$-invariant lifts]\tP\ \tQ\label{p3:lem:lift}
Let $L$ be a Poisson module. On the total space $C=L^*\setminus0$ there is a unique $\C^*$-invariant Poisson structure $\hat\pi=\sigma+u\,\theta\wedge\partial_u$ with $p_*\hat\pi=\sigma$ (local frame $e$, $D_\alpha e=\theta(\alpha)e$, $u$ the fibre coordinate), and sections of $L^k$ correspond to homogeneous functions of degree $k$. The Jacobi identity for $\hat\pi$ is equivalent to flatness and the Leibniz rule for $D$.
\end{lem}
\noindent The check of the local computation is only a sketch \tQ. Example: on $\mathbb P^2$ with $\Pi_F$ (quadratic Poisson structure on $\C^3$), $L=\mathcal O(m)$ for any $m$.

\begin{prop}\tP\label{p3:prop:graded}
If $X$ is projective and $L$ is an ample Poisson module, then $R=\bigoplus_{k\ge0}H^0(X,L^k)$ is a graded Poisson algebra with bracket of degree $0$ ($\{R_a,R_b\}\subset R_{a+b}$).
\end{prop}
\begin{proof}
By Lemma~\ref{p3:lem:lift}, $R$ is the ring of homogeneous functions on $C$, and the bracket of $\hat\pi$ preserves degree.\qedhere
\end{proof}
\begin{prop}[Quadratic almost-Poisson lift]\tP\ \tC\label{p3:prop:almost}
For $m\ge m_1$, $R^{(m)}=\bigoplus_kR_{mk}=S/I$ with $S=\mathrm{Sym}\,R_m$ and $I$ generated by quadrics (Mumford, \emph{Varieties defined by quadratic equations}, in: Questions on Algebraic Varieties, CIME 1969, Cremonese 1970; cited from memory, pages \tQ). Then there is a skew biderivation $\{\ ,\ \}'$ of degree $0$ on $S$ (a quadratic almost-Poisson structure) inducing the bracket of $R$ on $S/I$.
\end{prop}
\begin{proof}
Choose a linear section $s$ of $S_2\to R_{2m}$, put $\{y_a,y_b\}:=s(\{\bar y_a,\bar y_b\})$, and extend by Leibniz. For $q\in I_2$, $g\in S_1$: $\{q,g\}'\equiv\{\bar q,\bar g\}=0\bmod I_3$, so $\{I_2,S_1\}'\subset I_3$. Since $I$ is generated by $I_2$, Leibniz gives $\{I,S\}'\subset I$.\qedhere
\end{proof}
\noindent The Jacobiator $\mathcal J\in\mathrm{Hom}(\wedge^3S_1,I_3)$ of the lift is in general non-zero, and two lifts differ by $\epsilon\in\mathrm{Hom}(\wedge^2S_1,I_2)$. A Poisson extension exists iff the quadratic equation $\mathcal J_0+\mathcal L(\epsilon)+\mathcal J(\epsilon)=0$ has a solution $\epsilon$; the first-order obstruction class is $\mathcal J_0\bmod\mathrm{im}\,\mathcal L$.

\subsection{The toric case (proof)}
\begin{thm}[Poisson--Kodaira, toric]\tP\label{p3:thm:PKtoric}
Let $X$ be a smooth projective toric variety (torus $T=N\otimes\C^*$) and $\sigma=\sum\omega^{ij}\xi_i\wedge\xi_j$ ($\xi_i=x_i\partial_{x_i}$, $\omega\in\wedge^2N_\C$). Let $L=\mathcal O(D)$ be ample with polytope $P$ and $L^m$ very ample (sections $\chi^a$, $a\in mP\cap M$). Then $\nu:X\hookrightarrow\mathbb P^N$, $\nu^*y_a=\chi^a$, is a Poisson embedding for the diagonal quadratic Poisson structure on $\mathbb P^N$
\[\sigma_N=\sum_{a<b}\omega(a,b)\,y_ay_b\,\partial_{y_a}\wedge\partial_{y_b},\qquad\omega(a,b)=\sum\omega^{ij}a_ib_j.\]
This holds for every very ample power $L^m$.
\end{thm}
\begin{hyptab}
Toric, $\sigma$ torus-invariant & Brackets of characters are monomials: $\{\chi^a,\chi^b\}=\omega(a,b)\chi^{a+b}$. & Step (2). The lift can be chosen diagonal (removing the freedom that makes the generic lift of Prop.~\ref{p3:prop:almost} fail Jacobi).\\\midrule
$L^m$ very ample & $\nu$ is an embedding. & The ``embedding'' part of the conclusion.\\\midrule
$X$ smooth projective & Setting of classical Kodaira. & Same.
\end{hyptab}
\begin{proof}
(1) $\sigma_N$ is Poisson: for skew $\lambda_{ij}$ put $\{y_i,y_j\}=\lambda_{ij}y_iy_j$; then $\{y_i,\{y_j,y_k\}\}=\lambda_{jk}(\lambda_{ij}+\lambda_{ik})y_iy_jy_k$, and the cyclic sum vanishes by skew-symmetry (six terms cancel in pairs). The weight for the scaling is $0$, so it descends to $\mathbb P^N$. (2) On the torus, $\{\chi^a,\chi^b\}=\sigma(d\chi^a,d\chi^b)=\omega(a,b)\chi^{a+b}=\nu^*(\omega(a,b)y_ay_b)$ since $\chi^a\chi^b=\chi^{a+b}$. So $\nu$ is a Poisson map on $T$, hence on its closure $X$ (both $\sigma$ and $\nu^*\sigma_N$ are holomorphic).\qedhere
\end{proof}

\subsection{Non-toric case (cubics on $\mathbb P^2$, numerical check)}
\begin{obs}\tV\label{p3:obs:P2}
$X=\mathbb P^2$, $\sigma=\Pi_F$, Veronese $m=2$ ($\mathbb C^6$, $\dim I_2=6$, 90 unknowns $\epsilon$, 940 equations). The Jacobiator of the natural lift is non-zero (18 triples for Fermat, 19 for Hesse, 1 for $xyz$, also for the cusp). Solving the full quadratic system $\mathcal J_0+\mathcal L(\epsilon)+\mathcal J(\epsilon)=0$ by Levenberg--Marquardt from random starts gave solutions with residual $\le10^{-12}$ for $F=xyz$, Fermat, Hesse ($x^3+y^3+z^3-21xyz$) and the cusp ($x^3-y^2z$). For Fermat it was re-verified independently by finite differences (Jacobiator $\sim10^{-9}$, restriction error to the Veronese cone $\sim10^{-10}$). For Fermat the linear part $\mathcal J_0+\mathcal L(\epsilon)=0$ alone has no solution, so the quadratic term is essential. $m=1$ is tautological; $m\ge3$ was not computed (1215 unknowns, too heavy for this method). Scripts: \texttt{cp2/roundE\_*.py}.
\end{obs}
\noindent Status: numerical existence, not an exact algebraic certificate (the residual is machine precision).

\begin{conj}[Poisson--Kodaira]\tK\label{p3:conj:PK}
Let $(X,\sigma)$ be a projective holomorphic Poisson manifold and $L$ an ample Poisson module. Then there is $m_0$ such that for $m\ge m_0$ the Kodaira embedding by $|L^m|$ is a Poisson embedding into $\mathbb P^N$ for some quadratic Poisson structure $\sigma_N$.
\end{conj}
\noindent Evidence: Theorem~\ref{p3:thm:PKtoric} (toric) and Observation~\ref{p3:obs:P2} ($\mathbb P^2$, $m=2$). Open: a general obstruction theory (the class $\mathcal J_0\bmod\mathrm{im}\,\mathcal L$ is of Poisson-cohomology type; one expects it to vanish for $m\gg0$ by Serre-type vanishing of an $H^1$ of a Koszul complex, but this is not proved). No counterexample was found; failure for small $m$ remains possible.\\
\emph{Relation:} for $\sigma=0$ this is classical Kodaira. In regular type $M\to X$ with leaf $Y$, GH sections factor through $X$ (Proposition~\ref{p3:prop:leaf}), so embedding the Poisson pair $(M,\sigma_Y)$ requires ordinary (non-GH) holomorphic sections; this is the formulation of this section.

%--------------------------------------------------------------
\section{Non-reduced boundary}\label{p3:sec:nonred}
%--------------------------------------------------------------
For $\sigma=\Pi_F$ with $F$ having multiple factors, the dimensions of $\HH^*$ grow linearly in $m$. The following were computed by exact linear algebra (same machinery as Part~I, $m=-15,\dots,15$); script: \texttt{cp2/roundF\_nonreduced.py}.
\begin{obs}\tV\label{p3:obs:nr}
\textbf{(a) $F=x^2y$ (double line $\{x=0\}$ plus line $\{y=0\}$, one node, $N_L=\mathcal O(1)$, $U\simeq S^1$):} all 31 values agree with
\[m\ge-2:\ P(t)=[3\mid m](1+t)+3t^2+(m+2)(t+t^2);\qquad m=-3:\ 3t^2;\]
\[m\le-4:\ P(t)=[3\mid m](t^3+t^4)+3t^2+(-m-4)(t^2+t^3).\]
The case $m\le-4$ is the image under Serre duality $m\leftrightarrow-m-6$. The coefficient of the linear term is $m+2=h^0(L,\mathcal O(m+1))=h^0(L,\mathcal L|_L\otimes N_L)$. The same table holds for $x^2(x+y)$.\\
\textbf{(b) $F=x^3$ (triple line, $U=\C^2$):} for $m\ge0$, $P(t)=1+(2m+5)(t+t^2)+2t^2$ (i.e.\ $[1,2m+5,2m+7]$); $m=-1$: $[0,2,5]$; $m=-2,-3,-4$: $[0,0,3]$; for $m\le-5$ growth $\sim2|m|$ (e.g.\ $[0,0,5,2,0]$ at $m=-5$ and $[0,0,7,5,1]$ at $m=-6$; formula not determined). $2m+5=\sum_{j=1}^2h^0(L,\mathcal O(m+j))$.\\
\textbf{(c) $\mathbb P^1\times\mathbb P^1$, $D=2C$ ($C$ a smooth $(1,1)$-curve), $L=\mathcal O(m,m)$:} for $m\ge0$, $H^0=1$, $H^1=2m+3$, $H^2=2m+6$ (computed before Part~II).
\end{obs}
\begin{conj}[Non-reduced boundary]\tK\label{p3:conj:NR}
Let $X$ be a compact complex surface (here $\mathbb P^2$ or $\mathbb P^1\times\mathbb P^1$) and $D=\sum k_iD_i$ the divisor of $\sigma$ ($D_i$ smooth, $D_{\mathrm{red}}$ simple normal crossings). For $m\ge0$ (or in the range where higher cohomology of the relevant sheaves vanishes),
\[\HH^*(L)=H^*(U;L_U)\oplus\C^{\chi(D_{\mathrm{red}})}[-2]\oplus\bigoplus_i\bigoplus_{j=1}^{k_i-1}\bigl(H^0(D_i,L|_{D_i}\otimes N_{D_i}^{j})[-1]\oplus H^0(D_i,L|_{D_i}\otimes N_{D_i}^{j})[-2]\bigr).\]
\end{conj}
\noindent Checks: the three families (a)--(c) ($x^2y$ for $m\ge-2$, $x^3$ for $m\ge0$, double conic for $m\ge0$). The total Euler characteristic is $\chi(U)+\chi(D_{\mathrm{red}})=\chi(X)$, consistent with Theorem~E of Part~II. Heuristic: the Poisson cohomology of the local model $z^k\partial_z\wedge\partial_w$ has infinite-dimensional ($\C[[w]]$-module) parts generated in degrees $j=1,\dots,k-1$ (results of Monnier, Roger--Vanhaecke type \tC\ \tQ, not checked); their global part is $H^0$ of the coherent sheaves $L|_{D_i}\otimes N^j$, growing linearly in $m$. \tK\ No proof. It fails for small $m$ ($x^3$ at $m=-1$ does not fit), so the exact range of validity is unknown.

%--------------------------------------------------------------
\section{Local formula for quasi-homogeneous singularities (with correction)}\label{p3:sec:local}
%--------------------------------------------------------------
The observation in Part~II on the local formula for quasi-homogeneous plane curve singularities was stated with the rule $R(s)=[s<\mathrm{lct}](\cdots)$. This is wrong for $\mathrm{lct}<1/2$. We give the corrected version and the parts that could be proved. Setting: $h$ reduced quasi-homogeneous with an isolated singularity (weights $w_z,w_w$, degree $d$, $\mathrm{Sp}$ the Steenbrink spectrum, $\mathrm{lct}=(w_z+w_w)/d$, $\delta=d-w_z-w_w$), $\sigma=h\,\partial_z\wedge\partial_w$, $s\ne1$ the residue parameter of the module, $T=\mathcal O/J$ ($J$ the Jacobian ideal, $\dim T=\mu$). The stalk complex $\mathcal C_s=(\mathcal O\xrightarrow{d_0}\mathcal O^2\xrightarrow{d_1}\mathcal O)$ is as in Part~II.
\begin{prop}\tP\ \tV\label{p3:prop:Q}
$B_s=(\mathcal O\to\Theta(-\log D)\to h\mathcal O)$ is a subcomplex of $\mathcal C_s$, and the quotient $Q=\mathcal C_s/B_s\simeq[J\mathcal O_D\to\mathcal O_D]$ has differential $-(1-s)$ times the inclusion, which is injective for $s\ne1$. Hence $H^0(Q)=H^1(Q)=0$, $H^2(Q)=\mathcal O_D/J\mathcal O_D=\mathcal O/J$ (dimension $\mu$), $H^k(B_s)=H^k(\mathcal C_s)$ for $k\le1$, and $0\to H^2(B_s)\to H^2(\mathcal C_s)\to T\to0$. In other words, stalk $=\mu\,t^2+H^*(B_s)$ (dimensionwise).
\end{prop}
\noindent Check: $H(\mathcal C_s)=H(B_s)+T[-2]$ for $A_1,A_2,A_3,D_4,E_6$ (no mismatches) \tV.
\begin{prop}\tP\ \tQ\label{p3:prop:weight}
Assume $B_s\simeq(\Omega^\bullet(\log D),d+s\,\mathrm{dlog}\,h\wedge)$ (since $\sigma^\sharp$ is a chain map; the sign convention for $s$ was fixed by the data and not independently checked). By the Cartan formula for the Euler field $e=w_zz\partial_z+w_ww\partial_w\in\Theta(-\log D)$ ($e(h)=dh$), $H^*(B_s)$ is concentrated in the single weight $k_0=-sd$. The lowest weight of $B_s$ is $-\delta$, so $H^*(B_s)=0$ for $s>\delta/d=1-\mathrm{lct}$.
\end{prop}
\begin{obs}[Corrected rule]\tV\ \tC\label{p3:obs:rule2}
For $s\ne1$, dimensionwise,
\[\text{stalk}=\C^\mu[-2]\oplus R(s),\qquad R(s)=[\,s\in\Z,\ s\le0\,](1+t)+(t+t^2)\cdot\#\{\alpha\in\mathrm{Sp}:\ s+\alpha\in\Z,\ s+\alpha\le1\}.\]
This agrees with the rule of Part~II when $\mathrm{lct}\ge1/2$ but not when $\mathrm{lct}<1/2$: for $z^4-w^5$ ($\mathrm{lct}=9/20$) and $s=11/20=1-\mathrm{lct}$ there is a class (weight $-\delta$) although $s>\mathrm{lct}$. No mismatches for $A_1,A_2,A_3,D_4,D_5$ (weights $(3,2)$, degree 8), $E_6$, $T_{2,3,6}$ ($z^3-w^6$), $z^4-w^4$, $z^4-w^5$, $z^5-w^5$ (grids from $1/12$ to $1/20$ with $s\in[-2,2]$; $s=1$ excluded since the stalk is infinite-dimensional). Script: \texttt{cp2/roundF\_local\_rule2.py}. \tC\ Interpretation: the Wang sequence of the Milnor fibration $h:B\setminus D\to\C^*$ (fibre a bouquet of $\mu$ circles, monodromy eigenvalues $e^{-2\pi i\alpha}$ on $H^1$) gives $\text{stalk}(Rj_*L_s)=[s\in\Z](1+t)+(t+t^2)\#\{\alpha:s+\alpha\in\Z\}$. The complex $B_s$ computes the part of weight $\ge-\delta$, and the cutoff $s+\alpha\le1$ is this weight bound rewritten $\alpha$ by $\alpha$. Calder\'on-Moreno--Narv\'aez-Macarro, Proc.\ London Math.\ Soc.\ 98 (2009), arXiv:math/0603003, give the corresponding statement for free divisors (Jacobian ideal of linear type) for the family where $s$ is shifted by large integers (pages \tQ).
\end{obs}
\noindent \tK\ Not proved: that the weight-$k_0$ piece of $H(B_s)$ has exactly the dimension given by the Wang sequence for $s\le1-\mathrm{lct}$ (no mismatches in 10 types), and the $\alpha$-by-$\alpha$ cutoff $s+\alpha\le1$. The cusp and tacnode ($\mu=2,3$) are the special cases $A_2,A_3$ of the formula above; a general treatment of singularities is incomplete.\\
\tP\ Correction: $D_5=z^2w-w^4$ has weights $(3,2)$ and degree 8; the weights $(3,1,7)$ in an internal note were wrong. The published Parts I and II used the correct weights and are not affected. The Part~II observation on the local formula (stated for quasi-homogeneous germs, though only the listed types were checked) should be replaced by Observation~\ref{p3:obs:rule2} of this part. For the listed types ($A_k$, $D_k$, $E_6$, $E_8$, $T_{2,3,6}$) the old and new rules agree (checked), so the numbers in Part~II do not change.

%--------------------------------------------------------------
\section{Products, and the status of the general-dimension conjecture}\label{p3:sec:prod}
%--------------------------------------------------------------
\begin{prop}[K\"unneth]\tP\ \tQ\label{p3:prop:Kunneth}
Let $X,Y$ be compact complex manifolds with holomorphic Poisson structures $\sigma_X,\sigma_Y$ and Poisson modules $L_X,L_Y$. On $X\times Y$ take $\sigma=\sigma_X+\sigma_Y$ and $L=L_X\boxtimes L_Y$ (module structure $V_X+V_Y$). Then $K_{X\times Y}=K_X\boxtimes K_Y$ as complexes of sheaves (by $[\sigma_X,\sigma_Y]=0$ and the Koszul sign), and by the K\"unneth formula for sheaf hypercohomology, $P^{X\times Y}_{\HH}(t)=P^X_{\HH}(t)P^Y_{\HH}(t)$.
\end{prop}
\noindent References for K\"unneth (Kashiwara--Schapira, \emph{Sheaves on Manifolds}, \S2; Bredon, \emph{Sheaf Theory}, II.15) are \tQ. \tV\ For the toric block product on $\mathbb P^2\times\mathbb P^2$: $[1,4,12,16,16]=(1+2t+4t^2)^2$ (table of Part~II).
\begin{ex}\tP\label{p3:ex:prod33}
Let $X=Y=\mathbb P^2$ with smooth cubics $D_X,D_Y$ and $L=\mathcal O(m)\boxtimes\mathcal O(m')$. For $m,m'\in3\Z_{\ge0}$: $(1+2t^2)^2=1+4t^2+4t^4=P(H^*(U;\C))$ ($U=U_X\times U_Y$). For $m,m'\notin3\Z$: $(3t^2)(3t^2)=9t^4$. For $m,m'\in3\Z_{\le-6}$: $(2t^2+t^4)^2=P(H_c^*(U))$. Since $D=D_X\times Y\cup X\times D_Y$ is simple normal crossing with two components and no 4-fold points, this agrees with the prediction $H^*(U;L_U)\oplus0$ of the general-dimension conjecture of Part~II (generic log-symplectic, simple normal crossings). \tP\ So the conjecture holds on this non-toric product.
\end{ex}
\begin{rem}[Necessity of the genericity hypothesis (GEN)]\tP
When the divisors have nodes, $P=(H_X+\mu_Xt^2)(H_Y+\mu_Yt^2)=H_XH_Y+t^2(\mu_XH_Y+\mu_YH_X)+\mu_X\mu_Yt^4$, which exceeds the prediction $H(U)+\#D_{(4)}t^4$ ($\#D_{(4)}=\mu_X\mu_Y$) by the cross terms $t^2(\mu_XH_Y+\mu_YH_X)$. So $\sigma_X+\sigma_Y$ is not ``generic'' (consistent with the extra classes $t^2H^*(\mathcal O_\tau)$ found in the toric case, coming from non-degeneracy of the residue pairing on codimension-2 strata). Hence the conjecture needs a genericity hypothesis (GEN), and products are the typical family of counterexamples.
\end{rem}
\noindent \tK\ A genuinely non-product, non-toric test in dimension 4 was not carried out. Candidate: the Feigin--Odesskii elliptic Poisson structure $q_{5,1}$ on $\mathbb P^4$ ($\sigma^2\in H^0(\mathcal O(5))$; $D$ is the secant variety of the elliptic normal quintic, not simple normal crossings).

%--------------------------------------------------------------
\section{Comparison with earlier work}\label{p3:sec:compare}
%--------------------------------------------------------------
\begin{itemize}
\item \tC\ Pal--Poddar, \emph{Strong generalized holomorphic principal bundles}, arXiv:2404.18113, prove Kodaira--Nakano and Serre-type vanishing for positive SGH line bundles (Thm.~11.12) for compact regular generalized complex manifolds whose leaf space is a K\"ahler orbifold. Theorem~\ref{p3:thm:K1} is the embedding-side statement in the same setting. \tQ\ Equivalence with their definition of ``positive'' is not checked.
\item \tC\ T.\ Hu's work on K\"ahler Lie algebroids (arXiv:2403.10876, 2403.13424) treats Kodaira vanishing and Riemann--Roch, not the Kodaira embedding (in what we examined).
\item A Kodaira-type embedding theorem for general (type-jumping) generalized complex structures was not found \tC. Poisson--Kodaira (\S\ref{p3:sec:PK}) is the author's proposal.
\end{itemize}

%--------------------------------------------------------------
\section{Status of claims and open problems}\label{p3:sec:ledger}
%--------------------------------------------------------------
\begin{center}\small
\begin{tabular}{p{5.2cm}p{2.6cm}p{6.6cm}}\toprule
Claim & Status & Remark\\\midrule
Prop.~\ref{p3:prop:one} ($\dim\HH^0\le1$) & \tP & examples \tV\\
Prop.~\ref{p3:prop:leaf} (constant along leaves) & \tP & cites Gualtieri Thm.~4.3\\
Thm.~\ref{p3:thm:K1} (leaf-space) & \tP & Baily (orbifold) cited from memory \tQ\\
Lemma~\ref{p3:lem:lift} & \tP\ \tQ & local computation only sketched\\
Prop.~\ref{p3:prop:graded}, \ref{p3:prop:almost} & \tP & Mumford cited from memory \tQ\\
Thm.~\ref{p3:thm:PKtoric} (toric) & \tP & \\
Obs.~\ref{p3:obs:P2} ($\mathbb P^2$, $m=2$) & \tV & numerical existence; $m\ge3$ not computed\\
Conj.~\ref{p3:conj:PK} (Poisson--Kodaira) & \tK & \\
Obs.~\ref{p3:obs:nr} (non-reduced) & \tV & Conj.~\ref{p3:conj:NR} is \tK\\
Prop.~\ref{p3:prop:Q} (local, quotient part) & \tP\ \tV & \\
Prop.~\ref{p3:prop:weight} & \tP\ \tQ & sign convention of the $B_s$ identification unchecked\\
Obs.~\ref{p3:obs:rule2} (corrected rule) & \tV\ \tK & no mismatches in 10 types; proof partial\\
Prop.~\ref{p3:prop:Kunneth} (K\"unneth) & \tP\ \tQ & literature check\\
\bottomrule
\end{tabular}
\end{center}
\paragraph{Open problems.}(1) A formulation of a Kodaira-type theorem for general type-jumping generalized complex structures (a generalized-complex version of Poisson--Kodaira). (2) Vanishing of the obstruction class in Conjecture~\ref{p3:conj:PK} for $m\gg0$; an exact method for $\mathbb P^2$ with $m\ge3$. (3) A proof of Conjecture~\ref{p3:conj:NR} and its range of validity. (4) An $\alpha$-by-$\alpha$ proof of the local formula. (5) A generalized-complex version of Kodaira's proof by blow-up that avoids Riemann--Roch.\par
\paragraph{Scope of verification.}The computations in this part use the same exact linear algebra as Parts~I and~II (over the finite field $\mathbb F_{32003}$) or numerical (Levenberg--Marquardt) solutions; Observation~\ref{p3:obs:P2} rests on the latter. Equivalence with Pal--Poddar's definitions, the theorems of Baily and Mumford, the K\"unneth references, and the results of Monnier and Roger--Vanhaecke are unverified or cited from memory.



\clearpage
\part{Comparison with the regular type; the reduction theorem}
\setcounter{section}{0}\renewcommand{\thesection}{IV.\arabic{section}}
\noindent\textit{Part IV of the series; originally published as a separate note (\url{https://hylt-math.github.io/2026/11/22/generalized-geometry-poisson-part4/}). Within this paper, Part-internal references use the numbering IV.$\ast$.$\ast$.}\par
\section{Position and overview}\label{p4:sec:intro}
This part continues Parts I--III Part I, II, III (URLs in the bibliography). Parts I--III mainly treated generalized complex structures obtained by $\beta$-transform of a holomorphic Poisson structure $\sigma$ (type-jumping, log-symplectic case). Notation follows Parts I--III and is not re-explained.

\paragraph{Aim.}
We compare, with the same tool (the generalized Dolbeault complex $C(L;E)$), the \emph{regular-type} setting of Pal--Poddar (arXiv:2404.18113: constant type, symplectic leaves, Kähler orbifold leaf space) with the \emph{type-jumping} setting of Parts I--III. Conclusions:
\begin{itemize}
\item For regular type (flat symplectic fibre bundles), $\HH^*$ has a closed form via Gauss--Manin flat bundles $V^j$ (Theorem \ref{p4:thm:L}) \tP. The key is degeneration of the Leray spectral sequence at $E_2$, via Deligne's Lefschetz lemma.
\item Consequently, for each degree $k$ we get a Riemann--Roch formula \emph{polynomial} in $m$ (Corollary \ref{p4:cor:L}), and Pal--Poddar's Serre-type vanishing (Thm.\ 11.12(2)) coincides with our ``vanishing in transversal degree'' \tP.
\item In the type-jumping (log-symplectic) case the same quantities are \emph{periodic} in $m$. The difference boils down to whether $d_1=[\sigma,\cdot]$ is (almost) exact (\S\ref{p4:sec:contrast}) \tP\tV.
\item In general dimension in the log-symplectic case, the Part II conjecture ($\HH^*\cong H^*(U;L_U)\oplus\C^{\#D_{(n)}}[-n]$) is reduced to acyclicity of the quotient complexes on strata, condition (GEN$'$) (Theorem \ref{p4:thm:R}) \tP.
\end{itemize}
Open: general sufficient conditions for (GEN$'$); a framework unifying the two settings (Conjecture \ref{p4:conj:U}).

\section{Notation (new in this part)}\label{p4:sec:notation}
\begin{longtable}{>{\raggedright\arraybackslash}p{3.8cm}>{\raggedright\arraybackslash}p{11.2cm}}
\toprule Symbol & Meaning\\ \midrule\endhead
$Y\to M\xrightarrow{\pi} X$ & flat symplectic fibre bundle (\S\ref{p4:sec:setting}). $Y$ compact connected symplectic ($\dim Y=2m_Y$), $X$ compact complex ($\dim_\C X=k$)\\
$\Gamma=\pi_1(X)$, $\rho$ & $\rho:\Gamma\to\mathrm{Symp}(Y,\omega_Y)$ used to build $M=(Y\times\tilde X)/\Gamma$\\
$V^j$ & flat (Gauss--Manin) bundle with fibres $H^j(Y_x;\C)$; identified with the holomorphic bundle $V^j\otimes\mathcal O_X$\\
$G$, $E=\pi^*G$ & holomorphic line bundle on $X$ and its pullback (SGH line bundle)\\
$\HH^k(M,E)$ & generalized Dolbeault cohomology of the generalized holomorphic line bundle $E$ ($k$-th cohomology of $C(L;E)$)\\
$N$, $\wedge^bT_X$ & dual of the kernel of the anchor (rank $k$), and exterior powers of $T_X$; $b$ is the $N$-degree (Higgs-type degree)\\
$\Lambda=\omega_F\wedge$ & fibrewise Lefschetz operator ($\omega_F$ a lift of $\omega_Y$)\\
$X,D=\sum D_i,\sigma,L$ & log-symplectic setting (as in Part II); $\mathrm{res}_i=r_i$ are residues\\
$U=X\setminus D$, $L_U$ & complement and rank-one local system determined by the residues\\
$D_T=\bigcap_{i\in T}D_i$, $D_{(n)}$ & strata and the (finite) set of $n$-fold points\\
$B,\,C,\,Q$ & $B=\wedge^\bullet T(-\log D)\otimes L\subset C=(\wedge^\bullet T\otimes L)$, $Q=C/B$\\
$\nu_T,\ \sigma_T$ & $\nu_T=\sum_{i\in T}(r_i-1)[H_i]$ on the stratum $D_T$, and the Poisson structure induced by $\sigma$ (\S\ref{p4:sec:Q})\\
(GEN$'$) & intrinsic genericity assumption (Definition \ref{p4:def:GEN})\\
$P_{\HH}(t)$ & $\sum_k\dim\HH^k\,t^k$\\
\bottomrule
\end{longtable}
\noindent Tags: \tC\ literature; \tP\ author's proof (standard facts cited); \tV\ machine-checked; \tK\ conjecture/outlook; \tQ\ unverified.

\section{Regular type: flat symplectic fibre bundles}\label{p4:sec:regular}
\subsection{Setting}\label{p4:sec:setting}
\begin{defn}[Setting (S)]
Let $X$ be a compact complex manifold, $\tilde X\to X$ the universal (or a finite) cover, $\Gamma=\pi_1(X)$, and $(Y,\omega_Y)$ a compact connected symplectic manifold. For a homomorphism $\rho:\Gamma\to\mathrm{Symp}(Y,\omega_Y)$ put $M=(Y\times\tilde X)/\Gamma$. The product generalized complex structure ``symplectic on $Y$ $\times$ complex on $X$'' is $\Gamma$-invariant and descends to $M$, giving a regular-type generalized complex structure of type $k=\dim_\C X$. Leaves are the fibres $Y_x$, the leaf space is $X$. $E=\pi^*G$ ($G$ a holomorphic line bundle on $X$) is an SGH line bundle. \tC\ Definitions of regular type and SGH: Pal--Poddar, arXiv:2404.18113, Def.\ 3.1 etc.\ (numbering \tQ).
\end{defn}
\noindent This includes products $Y\times X$ ($\Gamma$ trivial), $S^2$-bundles with holonomy in $SO(3)$, Kodaira--Thurston type $T^2$-bundles, etc.

\subsection{Independence of $\omega$}
\begin{prop}\tP\label{p4:prop:L1}
Let $\tilde\omega$ be the closed 2-form on $M$ descending from $\mathrm{pr}_Y^*\omega_Y$. Then $e^{i\tilde\omega}$ is a Courant automorphism and $L=e^{i\tilde\omega}L_0$ ($L_0=T_{F,\C}\oplus T^{0,1}_{\mathrm{hor}}\oplus N^*$, $N^*=\pi^*T^{*1,0}_X$). Hence $L\cong L_0$ as Lie algebroids, and the generalized Dolbeault complex is
\[C(L;E)=\Omega^\bullet_{F_\C}(M;E\otimes\wedge^\bullet T^{1,0}_X),\qquad F_\C=T_{F,\C}\oplus T^{0,1}_{\mathrm{hor}}\ (\text{involutive}),\]
and $N$ acts trivially on SGH bundles. \emph{The cohomology does not see $\omega$} (only the foliated complex structure).
\end{prop}
\begin{proof}[Sketch]
Locally this is a statement about a product, and all objects are $\Gamma$-invariant. Since $d\tilde\omega=0$, $e^{i\tilde\omega}$ preserves the Courant bracket. $N$ acts trivially by the definition of SGH (holomorphic transversally, flat along leaves).\qedhere
\end{proof}
\begin{rem}\tC\tQ
Pal--Poddar's Thm.\ 8.9 and Cor.\ 8.10 treat the Dolbeault cohomology of transversal forms constant along leaves. It does not contain the fibrewise de Rham factor $H^j(Y)$ of our $\HH^*$ (\S\ref{p4:sec:PP}).
\end{rem}

\subsection{Main theorem: Leray and Lefschetz}
\begin{thm}[$\HH^*$ in regular type]\tP\label{p4:thm:L}
Assume (S) and (HL): $(Y,\omega_Y)$ satisfies hard Lefschetz (e.g.\ $Y$ Kähler). For $E=\pi^*G$,
\[\HH^k(M,E)\ \cong\ \bigoplus_{a+b+j=k}H^a\bigl(X,\ \wedge^bT_X\otimes V^j\otimes G\bigr).\]
No Kähler or projectivity assumption on $X$ is needed.
\end{thm}
\begin{hyptab}
(S) flat fibre bundle & Decomposition $L\cong L_0$ (Prop.\ \ref{p4:prop:L1}) and Ehresmann-type local triviality (fibre cohomology forms local systems $V^j$). & Step (i).\\\midrule
$Y$ compact connected & $H^j(Y)$ finite-dimensional; fibrewise cohomology sheaves are sheaves of smooth sections of $V^j$. & Step (i).\\\midrule
(HL) hard Lefschetz & $\Lambda=\omega_F\wedge$ gives isomorphisms $V^{m_Y-i}\to V^{m_Y+i}$ of flat bundles. & Step (iii): $d_r=0$ $(r\ge2)$ by Deligne's lemma.\\\midrule
$X$ compact & Each $H^a(X,\cdot)$ is finite-dimensional. & Sum of dimensions.\\\midrule
$E=\pi^*G$ (SGH) & $N$ acts trivially, so the $\wedge^bT_X$-degree $b$ splits off (no Poisson-type $d_1$). & Step (i): shape of $E_2$.
\end{hyptab}
\begin{proof}
(i) (Double complex) By Prop.\ \ref{p4:prop:L1}, $C(L;E)=\Omega^\bullet_{F_\C}(M;E\otimes\wedge^\bullet T_X)$. This is a double complex of ``fibrewise de Rham'' and ``Dolbeault on $X$ (coefficients $\wedge^bT_X\otimes G$)''. Since the fibres are compact, the fibrewise cohomology sheaves are the sheaves of $C^\infty$ sections of $V^j$. The induced $\bar\partial$ is the $(0,1)$-part of the Gauss--Manin connection and does not depend on the choice of lift (a change of lift is a fibrewise Lie derivative, exact on closed forms). Hence
\[E_2^{a,j}=\bigoplus_bH^a(X,\wedge^bT_X\otimes V^j\otimes G).\]
(ii) The mixed component of $\omega_F:=\tilde\omega|_F$ vanishes ($\tilde\omega$ is locally $\mathrm{pr}_Y^*\omega_Y$). So $\Lambda=\omega_F\wedge$ is a chain map of the double complex (commuting with $d_F$ and $\bar\partial$), acts on every page and commutes with $d_r$. On $E_2$ it is the pointwise fibrewise Lefschetz operator $V^j\to V^{j+2}$.
(iii) (Deligne's lemma) By hard Lefschetz $V^{m_Y-i}\to V^{m_Y+i}$ is an isomorphism. If $x$ is primitive of weight $s$, then $\Lambda^{s+1}d_2x=0$, but $d_2x$ lies in fibre degree $m_Y-s-1$, where $\Lambda^{s+1}$ is injective; so $d_2=0$. Inductively $d_r=0$ $(r\ge2)$, and dimensions add up.\qedhere
\src{Deligne, Th\'eor\`eme de Lefschetz et crit\`eres de d\'eg\'en\'erescence de suites spectrales, Publ.\ Math.\ IH\'ES 35 (1968) 259--278. Leray spectral sequence: Bredon, \emph{Sheaf Theory}, IV (precise location \tQ).}
\end{proof}
\begin{rem}[Without hard Lefschetz]\tK
Without (HL), $d_2\neq0$ may occur. We have computed no example.
\end{rem}

\subsection{Corollary: degree-wise Riemann--Roch and Serre vanishing}\label{p4:sec:cor}
\begin{cor}[Degree-wise Riemann--Roch]\tP\label{p4:cor:L}
Under the hypotheses of Theorem \ref{p4:thm:L}, assume moreover $X$ projective and $G$ ample. There is $m_0$ such that for $m\ge m_0$, $H^a(X,\wedge^bT_X\otimes V^j\otimes G^m)=0$ $(a>0)$, hence
\[\HH^k(M,E^m)=\bigoplus_{b+j=k}H^0\bigl(X,\wedge^bT_X\otimes V^j\otimes G^m\bigr),\]
\[\dim\HH^k(M,E^m)=\sum_{b+j=k}b_j(Y)\int_X\ch(\wedge^bT_X)\,e^{m\,c_1(G)}\,\Td(X).\]
The right-hand side is a polynomial in $m$ for each $k$.
\end{cor}
\begin{hyptab}
$G$ ample, $X$ projective & Serre vanishing $H^a(X,\mathcal F\otimes G^m)=0$ $(a>0,m\gg0)$ applied to the coherent sheaves $\mathcal F=\wedge^bT_X\otimes V^j$. & Vanishing of all terms with $a>0$.\\\midrule
$V^j$ flat & $\ch(V^j)=\mathrm{rk}\,V^j=b_j(Y)$ (trivial Chern classes). & $b_j(Y)$ factors out of the integral.\\\midrule
Theorem \ref{p4:thm:L} & Decomposes $\HH^k$ into $\bigoplus H^a$. & First equality.\\\midrule
HRR & Expresses $\dim H^0$ via $\chi$ (since $a>0$ terms vanish). & Last equality.
\end{hyptab}
\begin{proof}
Apply Serre vanishing \tC\ (Hartshorne III.5.2) to the locally free sheaves $\wedge^bT_X\otimes V^j$, taking the maximum of $m_0$ over the finitely many $(b,j)$. The rest is Hirzebruch--Riemann--Roch and $\ch(V^j)=b_j(Y)$.\qedhere
\end{proof}
\begin{rem}
The total Euler characteristic is $\sum_k(-1)^k\dim\HH^k=\chi(Y)(-1)^k\chi(X)$ (including the sign for $k=\dim_\C X$), consistent with the total Euler characteristic theorem of Part II \tV. Independence of $m$ holds only for the total alternating sum; each $\dim\HH^k$ grows as a polynomial in $m$.
\end{rem}

\subsection{Relation to Pal--Poddar Thm.\ 11.12}\label{p4:sec:PP}
\begin{thm}[Pal--Poddar, Thm.\ 11.12]\tC
Let $M$ be a compact regular-type generalized complex manifold, $M/\mathcal S$ a Kähler orbifold, $E$ a positive SGH line bundle. (1) If $p+q>k$ then $H^q(M,(\mathcal G^*M)^p\otimes E)=0$. (2) For every SGH bundle $E'$ and $m\ge m_0$, $q>0$: $H^q(M,E'\otimes E^m)=0$.
\src{D.\ Pal, M.\ Poddar, \emph{Strong generalized holomorphic principal bundles}, arXiv:2404.18113, Thm.\ 11.12 (hypotheses and statement re-read in pp\_recheck).}
\end{thm}
\begin{prop}\tP\label{p4:prop:PPcompare}
In setting (S), the Pal--Poddar groups $H^q(M,(\mathcal G^*M)^p\otimes E)$ read as $H^q(X,\Omega^p_X\otimes G)$ (transversal forms constant along leaves). Then (1) is the classical Kodaira--Nakano vanishing on $X$, and (2) is classical Serre vanishing for $E'=\pi^*(\wedge^bT_X\otimes V^j)$ (an SGH bundle: transition functions locally constant along leaves, holomorphic transversally), which is the same as the vanishing $H^a(X,\wedge^bT_X\otimes V^j\otimes G^m)=0$ $(a>0)$ of Corollary \ref{p4:cor:L}. Hence \emph{our vanishing in transversal degree is Pal--Poddar Thm.\ 11.12(2) applied to $E'=\wedge^bT_X\otimes V^j$}, and Corollary \ref{p4:cor:L} is ``Pal--Poddar's Serre-type vanishing $+$ Leray $+$ Lefschetz''.
\end{prop}
\begin{rem}[\tQ]
Taking sheaf cohomology literally on $M$, the fibre factor $H^j(Y)$ appears extra because $R^j\pi_*\pi^{-1}\neq0$. Since Pal--Poddar's Thm.\ 8.9 and Cor.\ 8.10 identify $H^q(M,(\mathcal G^*M)^p)$ with transversal Dolbeault cohomology, their $H^q$ should be read on the leaf space; we have not checked that the original says so explicitly. Also, vanishing of $\HH^k$ is asserted only in transversal degree $a>0$; the $H^j_{dR}(Y)$ factor survives in each degree $j$.
\end{rem}

\subsection{Examples}\label{p4:sec:examples}
\begin{ex}[Product $Y\times\mathbb P^1$]\tP\tV
$X=\mathbb P^1$, $G=\mathcal O(m)$, $m\ge-1$: by Theorem \ref{p4:thm:L}, $h^0(\mathbb P^1,\mathcal O(m))=m+1$ and $h^0(\mathbb P^1,T\otimes\mathcal O(m))=m+3$,
\[P_{\HH}(t)=P_Y(t)\bigl((m+1)+(m+3)\,t\bigr).\]
For $Y=S^2$: $P_Y=1+t^2$, so $P_{\HH}=(m+1)(1+t^2)+(m+3)(t+t^3)$ ($m=0$: $1,3,1,3$). For $Y=T^2$: $P_Y=1+2t+t^2$, $m=0$: $1,5,7,3$. Both grow linearly in $m$.
\end{ex}
\begin{ex}[Kodaira--Thurston: $T^2\to M\to E$]\tP\tV\tQ
For monodromy the shear $\psi=\begin{pmatrix}1&1\\0&1\end{pmatrix}$: $V^0=V^2=\mathcal O$, $V^1=F_2$ (Atiyah's non-trivial unipotent flat bundle, $h^0=h^1=1$ \tC\ Atiyah, Proc.\ London Math.\ Soc.\ (3) 7 (1957) 414--452). For $G=\mathcal O$,
\[P_{\HH}(t)=(1+t)^2(1+t+t^2)=1+3t+4t^2+3t^3+t^4.\]
Independent check: the Betti numbers of the Lie algebra complex of the Heisenberg algebra $\mathfrak h_3$ ($[X_2,W]=-\tfrac12X_1$) tensored with $\wedge N$ are $(1,3,4,3,1)$ \tV. This check assumes a Nomizu-type theorem (invariant forms compute $\HH^*$), which we have not proved \tQ\ (Theorem \ref{p4:thm:L} itself does not use it). For $\deg G=d>0$, $P_{\HH}(t)=d(1+t)^3$; for $\deg G=0$, $G\not\cong\mathcal O$, $P_{\HH}=0$.
\end{ex}
\begin{ex}[Bieberbach quotient $(T^2\times E)/\mathbb Z_2$]\tP\tV
$(y,x)\mapsto(-y,x+\tau/2)$. $V^1=L_{-1}\oplus L_{-1}$ (non-trivial sign representation) is acyclic for $G=\mathcal O$. $P_{\HH}(t)=(1+t)^2(1+t^2)$ (agrees with the $\mathbb Z_2$-invariant part of the Künneth formula).
\end{ex}
\begin{ex}[$S^2$-bundle with $SO(3)$ holonomy]\tP
The $V^j$ are trivial bundles, so the numbers coincide with the product case.
\end{ex}
\begin{rem}
An SGH bundle with non-trivial flat holonomy character $\chi$ along the leaves has $\chi$ locally constant over a compact base (maximum principle). If $Y$ is a torus and $\chi\neq1$, then $H^*(Y;L_\chi)=0$ and $\HH=0$. \tP
\end{rem}

\section{Contrast with the type-jumping case: polynomial vs periodic}\label{p4:sec:contrast}
For the log-symplectic $\mathbb{CP}^2$ of Parts I--III ($\sigma$ vanishing along a smooth cubic) and $L=\mathcal O(m)$, the dimensions of $\HH^*$ stay bounded and are 3-periodic in $m$.
\begin{center}\small
\begin{tabular}{p{3.6cm}p{4.0cm}p{5.4cm}}\toprule
 & regular type ($Y=S^2$, $X=\mathbb P^1$) & type-jumping ($\mathbb{CP}^2$, smooth cubic)\\\midrule
$\dim\HH^k$ for large $m$ & linear in $m$: $(m+1),(m+3),(m+1),(m+3)$ & bounded: $(1,0,2,0,0)$ for $m\in3\Z_{\ge0}$, otherwise $(0,0,3,0,0)$\\\midrule
$m=0$ & $(1,3,1,3)$ & $(1,0,2,0,0)$\\\midrule
$m=1$ & $(2,4,2,4)$ & $(0,0,3,0,0)$\\\midrule
$m=3$ & $(4,6,4,6)$ & $(1,0,2,0,0)$\\\midrule
total Euler characteristic & $\chi(Y)(-1)\chi(X)=-4$ & $\chi(\mathbb{CP}^2)=3$ (independent of $m$)\\\bottomrule
\end{tabular}
\end{center}
\tV\ The right column is Part I's table (agreeing with direct computation for $m=-15,\dots,15$); the left column is the formula of Corollary \ref{p4:cor:L}.
\begin{obs}\tP\label{p4:obs:contrast}
Both start from $E_1^{a,b}=H^a(\wedge^bT\otimes E)$. In regular type, $N$ acts trivially on SGH bundles, so $d_1=0$ (Theorem \ref{p4:thm:L}). In the type-jumping case $d_1=[\sigma,\cdot]+V\wedge$ is \emph{almost exact} on the $H^0$ terms: although $H^0(\wedge^bT\otimes L^m)$ grows with $m$, the remaining cohomology is a bounded quantity determined by the topology of $U=X\setminus D$. \emph{The difference between ``polynomial in $m$'' and ``periodic'' reduces to vanishing versus (near) exactness of $d_1$.}
\end{obs}

\section{Log-symplectic case in general dimension: a reduction theorem}\label{p4:sec:log}
\subsection{Setting and the quotient complex}\label{p4:sec:Q}
Let $X$ be a compact complex manifold ($\dim X=n$ even), $D=\sum D_i$ a simple normal crossing divisor, $\sigma$ log-symplectic ($\sigma^\sharp:\Omega^1(\log D)\cong T(-\log D)$), $L$ a Poisson module (flat log connection along $D$), $r_i$ the residues (constants on $D_i$) (Part II). Let $C=(\wedge^\bullet T\otimes L,\ d_\sigma+V\wedge)$ (a complex of sheaves of $\mathbb C$-vector spaces, $\HH^k=H^k(X;C)$), $B=\wedge^\bullet T(-\log D)\otimes L\subset C$ (isomorphic to $\mathrm{DR}_{\log}(L)$ via $\wedge\sigma^\sharp$ \tC\ Part II, ncd\_theorem Lemma 3), and $Q=C/B$.

\paragraph{Structure of $Q$.}
In logarithmic coordinates $(z_I,w)$ the $\mathcal O$-basis of $\wedge T$ is $\partial_{z_S}\wedge\partial_w^Y$, and that of $\wedge T(-\log D)$ is $E_S\wedge\partial_w^Y=z^S\partial_{z_S}\wedge\partial_w^Y$. Thus $Q$ is spanned by multivector fields with poles. Let $T(s)$ be the set of branches with poles. Because $\sigma\in\wedge^2T(-\log D)$ and the module vector field $V=\sigma^\sharp(\sum r_i\,\dlog z_i)\in T(-\log D)$, $d$ creates no new poles. Hence $F_k=\{|T(s)|\le k\}$ is an increasing filtration of $Q$ by subcomplexes, with graded pieces
\[\mathrm{gr}_kQ=\bigoplus_{|T|=k}Q_T,\qquad Q_T\cong\bigl(\wedge^\bullet T_{D_T}(-\log D'_T)\otimes L|_{D_T}\bigr)\ \text{(coefficient }z^{-T}E_T\text{)},\]
with differential $d_T=[\sigma_T,\cdot]+\nu_T\wedge$, where $\sigma_T$ is the Poisson structure induced by $\sigma$ on $\wedge^2T_{D_T}(-\log D'_T)$ ($D_T$ is a Poisson submanifold), $\nu_T=\sum_{i\in T}(r_i-1)[H_i]$, $H_i=\sigma^\sharp(\dlog z_i)$. The ``$-1$'' comes from $[\sigma,z_i^{-1}]=-z_i^{-1}H_i$, the ``$r_i$'' from the module structure. $[\sigma_T,\nu_T]=0$, $d_T^2=0$. \tP\ ($n=2$); for general $n$ the derivation is a sketch of the same Schouten-bracket computation, not written line by line \tQ.
The top piece $|T|=n$: $D_T=p$ is a point, $\wedge^\bullet T_{D_T}=\C$, and $\mathrm{gr}_nQ=\bigoplus_{p\in D_{(n)}}\C_p[-n]$ (zero differential). \tP

\begin{defn}[(GEN$'$)]\label{p4:def:GEN}
For every stratum index set $T$ with $1\le|T|\le n-1$, the complex of sheaves $Q_T$ on $D_T$ is acyclic (stalkwise).
\end{defn}
\begin{prop}[Necessary condition for (GEN$'$) and examples]\tP\label{p4:prop:genNec}
(a) $\nu_T\not\equiv0$ is necessary (if $\nu_T\equiv0$ then $Q_T$ is the Lichnerowicz complex of $\sigma_T$ and locally constant Casimirs enter $H^0$). In particular $r_i\neq1$ for $|T|=1$.
(b) For $|T|=1$ at a smooth point of $D$: with $\sigma=E_z\wedge\partial_{w_1}+\sum\partial_{w_{2j}}\wedge\partial_{w_{2j+1}}$ (log Darboux), $\nu=(r-1)\partial_{w_1}$ and $Q$ is the mapping cone of $(r-1)\cdot\mathrm{id}$ into the Lichnerowicz complex of the symplectic part, so acyclic $\iff r\neq1$. The log Darboux theorem is known for real $b$-symplectic structures (Guillemin--Miranda--Pires, Adv.\ Math.\ 264 (2014)); the holomorphic version should follow from the same Moser argument \tQ.
\end{prop}

\subsection{Main theorem}
\begin{thm}[Reduction theorem]\tP\label{p4:thm:R}
In the above setting assume (i) no $r_i$ lies in $\Z_{\ge1}$ and (ii) (GEN$'$). Then
\[\HH^*(L)\ \cong\ H^*(U;L_U)\ \oplus\ \bigoplus_{p\in D_{(n)}}\C[-n].\]
In particular $\chi(X)=\chi(U)+\#D_{(n)}$, i.e.\ $\sum_{0<k<n}\chi(D_k^\circ)=0$ ($D_k^\circ$ the open stratum of $k$-fold points).
\end{thm}
\begin{hyptab}
(i) $r_i\notin\Z_{\ge1}$ & Avoids resonance. & Step (1): $B\simeq Rj_*L_U$ by Deligne.\\\midrule
(ii) (GEN$'$) & The lower strata of $Q$ do not contribute to cohomology. & Step (2): $F_{n-1}Q$ is acyclic.\\\midrule
$D$ SNC, $\sigma$ log-symplectic & $B\cong\mathrm{DR}_{\log}(L)$; graded pieces of the filtration of $Q$ as above. & Steps (1), (2).\\\midrule
$X$ compact & $H^*(X;\cdot)$ finite-dimensional; Part II's Euler characteristic theorem applies. & Step (4).\\\midrule
(V) (vanishing), ampleness & \emph{Not needed.} & Used only in the toric weight computation (original proof of Part II Theorem T).
\end{hyptab}
\begin{proof}
(1) By non-resonance (i), $B\simeq Rj_*L_U$ (Deligne). \tC\ Deligne, \emph{\'Equations diff\'erentielles \`a points singuliers r\'eguliers}, LNM 163, II.3.13; Esnault--Viehweg.
(2) By the filtration and (ii), all graded pieces of $F_{n-1}Q$ are acyclic, so $F_{n-1}Q$ is acyclic. Hence $Q\to Q/F_{n-1}Q=\mathrm{gr}_nQ=\bigoplus_p\C_p[-n]$ is a quasi-isomorphism.
(3) The connecting morphism of the triangle $B\to C\to Q\to B[1]$ lies in $\mathrm{Hom}_{D(X)}(\bigoplus i_{p*}\C[-n],Rj_*L_U[1])=\bigoplus\mathrm{Hom}(j^*i_{p*}\C[-n],L_U[1])=0$ (as $p\in D$, $j^*i_{p*}=0$). Hence $C\simeq B\oplus Q$.
(4) Take $H^*(X;\cdot)$. $H^*(X;Rj_*L_U)=H^*(U;L_U)$. Euler characteristic: by Part II's total Euler characteristic theorem $\chi(\HH)=(-1)^n\chi(X)=\chi(X)$ ($n$ even); the right-hand side is $\chi(U)+\#D_{(n)}$, and $\chi(X)=\sum_k\chi(D_k^\circ)$ ($D_0^\circ=U$, $D_n^\circ=D_{(n)}$).\qedhere
\end{proof}
\begin{cor}[Reformulation of Part II Conjecture C]\tP
Conjecture C ($\HH^*\cong H^*(U;L_U)\oplus\C^{\#D_{(n)}}[-n]$) is equivalent to the vanishing of the contributions of the lower strata of $Q$, i.e.\ to (GEN$'$) ($\Leftarrow$ is Theorem \ref{p4:thm:R}; $\Rightarrow$ was checked in the toric case, where the extra classes correspond exactly to the non-acyclic $Q_T$).
\end{cor}

\subsection{The toric case}\label{p4:sec:toric}
Let $X$ be a toric variety, $\sigma=\tfrac12\sum A_{ij}E_iE_j$ (constant coefficients in the log frame, $A$ a non-degenerate antisymmetric matrix), $\rho\in M_\C$, $r_i=\langle\rho,u_i\rangle-a_i$. The weight-$\chi$ part of $Q_T$ ($T=\{i:\chi_i=-1\}$) is the Koszul complex of $\wedge(N_\C/W_T)$ with differential $v_\chi\wedge$, $v_\chi=A(\chi+\rho)$. It is acyclic $\iff v_\chi\notin W_T$. \tP
\begin{prop}\tP\label{p4:prop:toric}
Under (GEN) ($A_{ij}$ linearly independent over $\mathbb Q(\rho)$) and (R) ($r_i\neq1$), for every $T$ with $1\le|T|\le n-1$ the system $\langle x,u_i\rangle=r_i-1$ $(i\in T)$, $Ax\in W_T$ has no solution, so (GEN$'$) holds. Theorem \ref{p4:thm:R} then applies and $\HH^*=H^*(U;L_U)\oplus\C^{\chi(X)}[-n]$. This reproves Part II Theorem T without the vanishing (V) and ampleness (A).
\end{prop}
\begin{proof}[Sketch]
$\omega(x,m)=0$ $(m\in\mathrm{Ann}(W_T))$ gives $x\parallel m$, hence $r_i=1$, contradicting (R) (same as in the proof of Part II Theorem T).\qedhere
\end{proof}
\paragraph{Verification.}\tV\ \texttt{roundI\_check.py}: for $\mathbb P^2$, $\mathbb P^1\times\mathbb P^1$, $\mathbb P^4$, $\mathbb P^2\times\mathbb P^2$, $(\mathbb P^1)^4$, tested with generic $A$, block $A$, small-integer $A$, and 5 residue choices. In the 48 cases where no stratum obstructing (GEN$'$) has a solution, all agree with the generic prediction (0 counterexamples). In the 25 cases where some stratum has a solution, $\HH\neq$ the generic prediction. In 92 cases a solution exists but is not realized (false alarms).

\begin{ex}[$\mathbb P^2\times\mathbb P^2$, block $A=A_X\oplus A_Y$]\tP\tV
For $\rho=0$, $m=1$: the mixed block $A_{XY}=0$, so on $T=(\text{a node of }X)\times Y$ we have $\nu_T=0$. Then $Q_T$ is the log de Rham complex of $Y$, $\simeq Rj_*\C_{U_Y}$, and each such $T$ contributes $t^2(1+t)^2$. There are $3+3=6$ such $T$, and measured $-$ generic prediction $=[1,4,12,16,16]-[1,4,6,4,10]=[0,0,6,12,6]=6t^2(1+t)^2$ \emph{exactly}. This is a mechanistic explanation of Part III's observation that products of divisors with nodes produce cross terms $t^2(\mu_XH_Y+\mu_YH_X)$.
\end{ex}
\begin{conj}[\tK\ non-acyclic strata]\label{p4:conj:nonacyc}
When (GEN$'$) fails, $E_1^T=H^*(D_T;Q_T)\Rightarrow H^*(X;Q)$. If all $Q_T$ with $\nu_T=0$ are of ``product type'', the extra part is $\sum_Tt^{|T|}H^*(U_T;L_T)$.
\end{conj}

\section{A unification of the two settings?}\label{p4:sec:unified}
In both settings $\HH(E)=H^*(M;C(L;E))$, $L$ is a Lie algebroid, and $N=\ker(\text{anchor})$ has rank $k$ at regular points and jumps at type-jump points. There is a filtration by $N$-degree $b$ with $E_1^{a,b}=H^a(\wedge^bT\text{-type};E)$, and $d_1$ is the action of $N$ (zero for SGH in regular type; $[\sigma,\cdot]+V\wedge$ for $\beta$-transforms).
\begin{center}\small
\begin{tabular}{p{3.2cm}p{5.4cm}p{5.4cm}}\toprule
 & regular type & type-jumping (log-symplectic)\\\midrule
fibrewise de Rham & $H^j(Y)$ of the compact leaf $Y$ (local system $V^j$) & $H^j(U;L_U)$ of the open leaf $U$ \\\midrule
transversal direction & Dolbeault on $X$ & Lichnerowicz-type complexes $Q_T$ on the strata of $D$\\\midrule
treatment of boundary & none ($Y$ compact) & choice of $Rj_*$ / $j_!$ by residues (Deligne)\\\midrule
monodromy & Gauss--Manin & $e^{2\pi ir}$ (around $D$)\\\midrule
degeneration & $d_r=0$ by Lefschetz ($Y$ compact) & $U$ non-compact, no analogue; $d_1\neq0$\\\bottomrule
\end{tabular}
\end{center}
\begin{conj}[\tK\ Unification conjecture (imprecise)]\label{p4:conj:U}
For a compact generalized complex manifold with finitely many types of strata and a module $E$ on which $N$ acts trivially along the regular strata, $\HH^*(E)$ is filtered by stratum-wise contributions of ``fibrewise de Rham on the leaf space $\otimes$ transversal Dolbeault'', and on each stratum whether the type is $(c)$ (compact support) or non-$(c)$ is decided by the indicial roots (residues).
\end{conj}
\noindent This conjecture is not formulated precisely, and has been tested nowhere beyond the two families above (regular type and log-symplectic).

\section{Status ledger and open problems}\label{p4:sec:ledger}
\begin{center}\small
\begin{tabular}{p{5.2cm}p{2.0cm}p{7.2cm}}\toprule
Claim & Status & Remarks\\\midrule
Prop.\ \ref{p4:prop:L1} (independence of $\omega$) & \tP & Pal--Poddar definition numbers \tQ\\
Thm.\ \ref{p4:thm:L} (Leray $+$ Lefschetz) & \tP & needs hard Lefschetz (HL); case without (HL) not studied \tK\\
Cor.\ \ref{p4:cor:L} (degree-wise RR) & \tP & \\
Prop.\ \ref{p4:prop:PPcompare} (agreement with Pal--Poddar) & \tP\ \tQ & reading of Pal--Poddar's $H^q$ unverified\\
Examples (product, Bieberbach, $S^2$-bundle) & \tP\ \tV & \\
Kodaira--Thurston $(1,3,4,3,1)$ & \tV\ \tQ & check assumes a Nomizu-type theorem; Theorem \ref{p4:thm:L} itself does not use it\\
Obs.\ \ref{p4:obs:contrast} (polynomial vs periodic) & \tP\ \tV & \\
Prop.\ \ref{p4:prop:genNec} (a) & \tP & \\
Prop.\ \ref{p4:prop:genNec} (b) & \tP\ \tQ & holomorphic log Darboux only cited (real case)\\
Thm.\ \ref{p4:thm:R} (reduction theorem) & \tP & derivation of the filtration of $Q$ for general $n$ is a sketch \tQ\\
Prop.\ \ref{p4:prop:toric} (toric reproof) & \tP\ \tV & \\
$6t^2(1+t)^2$ for $\mathbb P^2\times\mathbb P^2$ & \tP\ \tV & \\
Conj.\ \ref{p4:conj:nonacyc}, \ref{p4:conj:U} & \tK & \\
\bottomrule
\end{tabular}
\end{center}
\paragraph{Open problems.}(1) Does ``$\nu_T\neq0$ everywhere $\Rightarrow Q_T$ acyclic'' hold for non-linear $\sigma$ (the homotopy $h=\iota_\theta$ needs invertibility of $1+L_{X_\theta}$)? (2) The holomorphic log Darboux (linearization) at crossings (codimension $\ge2$). (3) Non-flat bundles, orbifold bases, and $Y$ not satisfying hard Lefschetz. (4) A precise formulation of Conjecture \ref{p4:conj:U}. (5) Tests on genuinely non-product, non-toric examples of dimension 4.\par
\paragraph{Scope of verification.}Apart from the Kodaira--Thurston example, the computations are exact linear algebra over finite fields as in Parts I--III, or known Künneth-type computations. Properties of Atiyah's $F_2$, the application of Deligne's lemma and the correspondence with Pal--Poddar's definitions are citations of standard facts or sketches, and partly unverified (\tQ).



\clearpage
\part{Riemann--Roch-type formulas, the index conjecture, proof of the local formula}
\setcounter{section}{0}\renewcommand{\thesection}{V.\arabic{section}}
\noindent\textit{Part V of the series; originally published as a separate note (\url{https://hylt-math.github.io/2026/11/22/generalized-geometry-poisson-part5/}). Within this paper, Part-internal references use the numbering V.$\ast$.$\ast$.}\par
%--------------------------------------------------------------
\section{Overview: position of this part and results}
%--------------------------------------------------------------
Parts I--IV studied the generalized Dolbeault cohomology $\HH^\bullet(L)$ of generalized holomorphic line bundles $L$ on the generalized complex structure obtained by a $\beta$-transform from a holomorphic Poisson structure $\sigma$ on a complex manifold $X$, starting from concrete examples. This part collects three things.

\begin{itemize}
\item \textbf{Riemann--Roch-type formulas.} (a) The total Euler characteristic does not depend on $c_1(E)$ (Part II Theorem E, recalled). (b) Its logarithmic version: the Euler characteristic of the log de Rham complex of any $L$ with a flat log connection equals $\chi(U)$ (Theorem \ref{p5:thm:K1}) \tP. (c) For $U$ affine and $n$ even, the whole Poincar\'e polynomial is expressed through $\chi(X)$ and the dimensions $h^k(U;L_U)$ (Theorem \ref{p5:thm:K}) \tP\tV. (d) In regular type the dimension is a polynomial in $m$ in each degree (Part IV Cor.\ L). We contrast these three ``behaviours of dimensions'' (\S\ref{p5:sec:three}).
\item \textbf{Proof of the local formula.} The stalk formula for plane curve germs with quasi-homogeneous isolated singularity, which was only a computational observation and a conjecture in Part II, is proved for $s\neq1$ (Theorem \ref{p5:thm:local}) \tP.
\item \textbf{The index conjecture.} We state precisely a conjecture on a weighted b-calculus on the log tangent algebroid, as the ``index-theorem version'' of Riemann--Roch, and separate what remains unproved (Conjecture \ref{p5:conj:index}) \tK.
\end{itemize}
Finally we summarize, claim by claim, where Parts I--V stand with respect to the original goal (a generalized-geometry version of Riemann--Roch, Serre vanishing and Kodaira embedding) (\S\ref{p5:sec:ledger}). \emph{The index conjecture is not proved.}

%--------------------------------------------------------------
\section{Notation (new or recalled)}\label{p5:sec:notation}
%--------------------------------------------------------------
\begin{longtable}{>{\raggedright\arraybackslash}p{3.8cm}>{\raggedright\arraybackslash}p{11.2cm}}
\toprule Symbol & Meaning\\ \midrule\endhead
$X$, $n$ & compact complex manifold, $\dim_\C X=n$; $n$ is even in the log-symplectic case\\
$\sigma$ & holomorphic Poisson structure on $X$. ``Log-symplectic'' means that the zero divisor of $\sigma^{n/2}$ is a reduced simple normal crossings (SNC) divisor $D$ and $\sigma$ is logarithmically non-degenerate along $D$\\
$D=\bigcup D_i$, $U=X\setminus D$ & SNC divisor and its complement. $D_T=\bigcap_{i\in T}D_i$; $D_{(n)}$ is the set of $n$-fold points\\
$L$, $r_i$ & line bundle which is a Poisson module (generalized holomorphic line bundle). The residue $r_i\in\C$ along $D_i$ is constant (Esnault--Viehweg); $c_1(L)=-\sum r_i[D_i]$ in $H^2(X;\C)$\\
$L_U$ & flat line bundle on $U$ (sheaf of flat sections of $L$); its monodromy around $D_i$ is determined by $r_i$ (Deligne)\\
$T(-\log D)$, $\Omega^p(\log D)$ & log tangent bundle, sheaves of log differential forms\\
$\mathrm{DR}_{\log}(L)$ & log de Rham complex $(\Omega^\bullet(\log D)\otimes L,\nabla)$\\
$\HH^k(L)$ & generalized Dolbeault cohomology: hypercohomology of $(\wedge^\bullet T\otimes L,\ \dbar+[\sigma,\cdot])$\\
$P_{\HH}(t)$ & $\sum_k\dim\HH^k(L)\,t^k$\\
$h^k(U;L_U)$ & $\dim H^k(U;L_U)$\\
$\chi(X)$, $\chi(U)$ & topological Euler characteristics\\
$\partial U$ & boundary of a tubular neighbourhood of $D$ (the end of $U$); a closed real $(2n-1)$-manifold ($\chi(\partial U)=0$)\\
$h$, $(w_z,w_w)$, $N$ & defining function of a quasi-homogeneous isolated singularity, weights, degree; $\delta=N-w_z-w_w$\\
$J=(h_z,h_w)$, $T=\mathcal O/J$, $\mu$ & Jacobian ideal, Milnor algebra, Milnor number $\mu=\dim T$\\
$\mathrm{Sp}$ & spectrum: for each monomial basis element $m$ of $T$, $\alpha=(\mathrm{wt}(m)+w_z+w_w)/N$\\
$s$ & residue (local exponent) of the local model $\mathcal O\,h^s$; $\nabla_s=d+s\,\dlog h$\\
$\chi_h$, $e$ & Hamiltonian vector field $h_w\partial_z-h_z\partial_w$ and Euler field $w_zz\partial_z+w_ww\partial_w$ ($e h=Nh$)\\
\bottomrule
\end{longtable}
(To avoid confusion with $\chi(X)$ the Hamiltonian field is written $\chi_h$.)

\paragraph{Status tags.} \tP\ proved here or in an earlier part; \tC\ result from the literature; \tV\ checked by computer (exact rational/finite-field linear algebra); \tK\ conjecture; \tQ\ unverified (e.g.\ citations from memory). When several tags occur, each refers to the corresponding part.

%--------------------------------------------------------------
\section{Recollection of known results}\label{p5:sec:recall}
%--------------------------------------------------------------
We list the statements from earlier parts that are used here, with their status only (proofs are in the respective parts). Part I: Part I, II: Part II, III: Part III, IV: Part IV.

\begin{thm}[Euler characteristic theorem (Part II, Theorem E) \tP]\label{p5:thm:E}
Let $X$ be a compact complex manifold, $\sigma$ any holomorphic Poisson structure and $E$ a generalized holomorphic vector bundle of rank $\rk E$. Then $\sum_k(-1)^k\dim\HH^k(E)=(-1)^n\,\rk(E)\,\chi(X)$.
\end{thm}
Sketch: filter the complex by the Higgs degree $b$; the graded pieces are the complexes of $H^a(X;\wedge^bT\otimes E)$, and the $\sigma$-part of the differential lowers the filtration, so the Euler characteristic is $\sum_b(-1)^b\chi(\wedge^bT\otimes E)$. By Hirzebruch--Riemann--Roch, $\sum_b(-1)^b\mathrm{ch}(\wedge^bT)=c_n(T)\mathrm{Td}(T)^{-1}$ starts in degree $n$, so higher terms of $\mathrm{ch}(E)$ do not contribute and $\rk E\cdot\int c_n(T_X)$ remains. \emph{Hence it is completely independent of $c_1(E)$.}

\begin{thm}[Reduction theorem (Part IV, Theorem R) \tP\ \tQ]\label{p5:thm:R}
Assume $\sigma$ is log-symplectic along an SNC divisor $D$, no residue $r_i$ lies in $\Z_{\ge1}$, and (GEN$'$) holds (for $0<|T|<n$ the quotient complex $Q_T$ on the stratum $D_T$ is acyclic). Then $\HH^*(L)\cong H^*(U;L_U)\oplus\bigoplus_{p\in D_{(n)}}\C[-n]$.
\end{thm}
For $n=2$ this is the theorem of Part I. The derivation from the filtration of $Q$ in general $n$ is closer to a sketch in Part IV, and we carry the \tQ\ over. In the toric case (GEN$'$) is proved (Part IV, Proposition).

\begin{thm}[Regular type (Part IV, Theorem L and Cor.\ L) \tP\ \tQ]\label{p5:thm:L}
Let $Y\to M\to X$ be a flat symplectic fibre bundle ($Y$ compact connected symplectic satisfying hard Lefschetz, $X$ compact complex) and $E=\pi^*G$ for a holomorphic line bundle $G$ on $X$. Then $\HH^k(M,E)\cong\bigoplus_{a+b+j=k}H^a(X,\wedge^bT_X\otimes V^j\otimes G)$ ($V^j$ the Gauss--Manin bundle of the fibre). In particular for $G=G_0^m$ ($G_0$ ample) the dimension in each degree $k$ is a polynomial in $m$ ($m\gg0$).
\end{thm}
It uses the Leray spectral sequence and Deligne's Lefschetz lemma; the points involving a Nomizu-type theorem are marked \tQ\ in Part IV.

\paragraph{Other earlier results (statements only).} Duality $\HH^k(L)^*\cong\HH^{2n-k}(L^\vee\otimes K^2)$ (Part I, any dimension) \tP; Serre-type vanishing in the transversal degree, failing in total degree (Part I) \tP; impossibility of embedding by generalized holomorphic sections, leaf-space Kodaira, Poisson--Kodaira conjecture in the toric case (Part III) \tP\ \tK.

%--------------------------------------------------------------
\section{Euler characteristic of the log de Rham complex (Theorem K1)}\label{p5:sec:K1}
%--------------------------------------------------------------
\begin{thm}[K1: log version of the Euler characteristic \tP]\label{p5:thm:K1}
Let $X$ be a compact complex manifold, $D\subset X$ an SNC divisor and $L$ a line bundle with a flat logarithmic connection $\nabla:L\to L\otimes\Omega^1_X(\log D)$. Then $\mathbb H^k(X;\mathrm{DR}_{\log}(L))$ is finite dimensional and
\[\sum_k(-1)^k\dim\mathbb H^k(X;\mathrm{DR}_{\log}(L))\;=\;\int_Xc_n\bigl(T_X(-\log D)\bigr)\;=\;\chi(U)\;=\;\sum_{T}(-1)^{|T|}\chi(D_T),\]
independent of $L$ and of the residues $r_i$.
\end{thm}
\begin{hyptab}
$X$ compact & $H^q(X;\Omega^p(\log D)\otimes L)$ finite dimensional; HRR available. & Steps (1), (2).\\\midrule
$D$ SNC & $\Omega^p(\log D)$ locally free, and $\Omega^\bullet(\log D)\simeq Rj_*\C_U$ (log Poincar\'e lemma). & Step (3) (to get the topological value for $L=\mathcal O$).\\\midrule
$\nabla$ flat log connection & makes $\mathrm{DR}_{\log}(L)$ a complex; in rank 1 the residues are constant. & Step (1).\\\midrule
residues $r_i\notin\Z_{\ge1}$ & \emph{not needed.} & Only for the identification $\mathrm{DR}_{\log}(L)\simeq Rj_*L_U$, not for the Euler characteristic.\\\midrule
$\sigma$ (log-symplecticity) & \emph{not needed.} & Only to relate to the generalized Dolbeault complex ($\wedge^\bullet T(-\log D)\otimes L\cong\mathrm{DR}_{\log}(L)$).
\end{hyptab}
\begin{proof}
(1) The stupid filtration of $\mathrm{DR}_{\log}(L)$ has $E_1$-term $H^q(X;\Omega^p(\log D)\otimes L)$, finite dimensional; hence $\chi(\mathrm{DR}_{\log}(L))=\sum_p(-1)^p\chi(X;\Omega^p(\log D)\otimes L)$.

(2) By HRR this equals $\int_X\bigl[\sum_p(-1)^p\mathrm{ch}(\Omega^p(\log D))\bigr]\mathrm{ch}(L)\,\mathrm{Td}(T_X)$. For the rank $n$ bundle $F=\Omega^1(\log D)$ one has $\sum_p(-1)^p\mathrm{ch}(\wedge^pF)=c_n(F^\vee)\,\mathrm{Td}(F^\vee)^{-1}$, which starts in degree $n$. So only the constant term $1$ of $\mathrm{ch}(L)\mathrm{Td}(T_X)$ matters and the result is $\int_Xc_n(T_X(-\log D))$, independent of $L$.

(3) For $L=\mathcal O$ ($\nabla=d$) the log Poincar\'e lemma (Deligne \cite{Del}) $\Omega^\bullet(\log D)\simeq Rj_*\C_U$ shows that the hypercohomology is $H^*(U;\C)$, so the Euler characteristic in (1) is $\chi(U)$; by (2) it equals $\int c_n(T_X(-\log D))$. Finally $\chi(U)=\sum_T(-1)^{|T|}\chi(D_T)$ follows from $\chi_c=\chi$ on the open even-dimensional manifold $U$ (Poincar\'e duality) and additivity of $\chi_c$.
\end{proof}
\paragraph{Check.} \tV\ For $X=\mathbb P^2$, $D$ a smooth curve of degree $d=1,\dots,7$, $L=\mathcal O(a)$, $a=-12,\dots,12$, we computed $\sum_p(-1)^p\chi(\Omega^p(\log D)\otimes L)$ directly: in all 175 cases it equals $3-3d+d^2=\chi(U)$ (script \texttt{cp2/roundK\_check.py}).

\begin{cor}[Wall-crossing identity \tP]\label{p5:cor:wall}
For a rank one local system $L_U$, $\chi_c(U;L_U)=\chi(U;L_U)=\chi(U)$. This follows from $\chi(\partial U)=0$ (closed odd-dimensional manifold) and the Euler characteristic of the triangle $j_!L_U\to Rj_*L_U\to i_*i^*Rj_*L_U$. Hence the Euler characteristic does not change when a residue crosses an integer and $j_!$ and $Rj_*$ exchange roles.
\end{cor}
\paragraph{Relation to Theorem E.} Theorem \ref{p5:thm:E} and the present theorem are two faces of the same phenomenon: ``the Euler characteristic does not see the higher terms of $\mathrm{ch}(L)$''. The total Euler characteristic of $\HH^*$ is $(-1)^n\chi(X)$, that of the comparison subcomplex $B\simeq\mathrm{DR}_{\log}(L)$ is $\chi(U)$, and the difference $(-1)^n\chi(X)-\chi(U)$ is carried by the Euler characteristic of the quotient complex $Q$. For $n$ even and under (GEN$'$) the difference is $\#D_{(n)}$ (Theorem \ref{p5:thm:R}).

%--------------------------------------------------------------
\section{A Riemann--Roch-type formula for the Poincar\'e polynomial (Theorem K)}\label{p5:sec:K}
%--------------------------------------------------------------
\begin{thm}[K \tP\ \tV]\label{p5:thm:K}
Assume the hypotheses of Theorem \ref{p5:thm:R} (SNC, log-symplectic, $r_i\notin\Z_{\ge1}$, (GEN$'$)) and in addition (i) $U$ is affine and (ii) $n$ is even (automatic in the log-symplectic case). Then
\[P_{\HH}(t)\;=\;\chi(X)\,t^n\;+\;\sum_{k=0}^{n-1}h^k(U;L_U)\,\bigl(t^k-(-1)^kt^n\bigr).\]
\end{thm}
\begin{hyptab}
hypotheses of the reduction theorem & give $\HH^*\cong H^*(U;L_U)\oplus\C^{\#D_{(n)}}[-n]$. & Step (1).\\\midrule
$U$ affine & Artin vanishing: $H^k(U;L_U)=0$ for $k>n$ for a local system. & Step (2). Without it the sum could not be confined to $k<n$.\\\midrule
$n$ even & $\sum_k(-1)^kh^k+\#D_{(n)}=\chi(X)$ holds without sign (Theorem \ref{p5:thm:E}). & Step (3).\\\midrule
projectivity / ampleness & \emph{not needed.} & ---
\end{hyptab}
\begin{proof}
(1) By Theorem \ref{p5:thm:R}, $P_{\HH}=\sum_{k\le n}h^kt^k+\#D_{(n)}t^n$ (terms with $k>n$ vanish by (2)). (2) By Artin's vanishing theorem \cite{SGA4}, $h^k=0$ for $k>n$. (3) By Theorem \ref{p5:thm:E} ($n$ even), $\sum_{k\le n}(-1)^kh^k+\#D_{(n)}=\chi(X)$, i.e.\ $h^n+\#D_{(n)}=\chi(X)-\sum_{k<n}(-1)^kh^k$. Substituting into the coefficient of $t^n$ gives the formula.
\end{proof}
\paragraph{Comparison with Part I.} For $n=2$, $r\in\Z_{\le0}$ and $L_U$ trivial, $h^0=1$ and $h^1=b_1$, so $P=\chi(X)t^2+(1-t^2)+b_1(t+t^2)=\chi(X)t^2+\Delta(t)$, the $[r\le0]$ term of the theorem of Part I. For $r\notin\Z$ Part I gives $P=\chi(X)t^2$; Theorem \ref{p5:thm:K} reproduces this only if $H^k(U;L_U)=0$ for $k<n$, which we have not shown for general $n$ \tQ\ (for $n=2$ it was checked in Part I). The side $r\ge2$ follows from duality.

\paragraph{Checks.} \tV\ (a) 188 tables for $n=2$; (b) 61 toric cases for $n=2,4$ ($h^k=\binom nk$ if $\rho\in M$, $0$ otherwise): the measured Poincar\'e polynomials agree with the formula in every case (0 mismatches). Toric $\mathbb P^2$, $\rho\in M$: $P=3t^2+(1-t^2)+2(t+t^2)=1+2t+4t^2$.

\subsection{Three behaviours of the dimensions}\label{p5:sec:three}
Taking $L=L_0^m$ and varying $m$, the cohomology dimensions behave in three qualitatively different ways depending on the setting.

\begin{center}\small
\begin{tabular}{>{\raggedright\arraybackslash}p{3.0cm}>{\raggedright\arraybackslash}p{4.4cm}>{\raggedright\arraybackslash}p{6.4cm}}
\toprule Behaviour & Dimension in each degree $k$ & Example\\\midrule
(A) regular type & \emph{polynomial} in $m$ ($m\gg0$) \tP & $X=\mathbb P^1$, $Y=T^2$ (product), $G=\mathcal O(m)$: from Theorem \ref{p5:thm:L}, $(m+1,\ 3m+5,\ 3m+7,\ m+3)$ ($m\ge-1$)\\\midrule
(B) log-symplectic & \emph{periodic} jumps governed by integrality of the residue $r=-m/3$ \tP\ \tV & $\mathbb P^2$, smooth cubic: $(1,0,2,0,0)$ for $m=3k\ge0$, $(0,0,3,0,0)$ otherwise ($m\ge-3$)\\\midrule
(C) total Euler characteristic & \emph{independent} of $m$ \tP & in (A) $\sum(-1)^kh^k=0$; in (B) $3=\chi(\mathbb P^2)$\\\bottomrule
\end{tabular}
\end{center}
In example (A) the Euler characteristic is $(m+1)-(3m+5)+(3m+7)-(m+3)=0$, consistent with Theorem \ref{p5:thm:E} ($\chi(M)=\chi(Y)\chi(X)=0$). The computation in (A) is the formula of Theorem \ref{p5:thm:L} specialized to a product, using $T=\mathcal O(2)$ on $\mathbb P^1$ and $\dim V^j=(1,2,1)$ (in the range $m\ge-1$ where no $H^1$ terms occur).

\paragraph{Reading.} By (C), the naive Riemann--Roch form ``$\chi(E)=\int\mathrm{ch}(E)\mathrm{Td}$ is a polynomial in $c_1(E)$'' \emph{degenerates} in generalized geometry (it does not see $c_1$); the information lies in the graded dimensions. In (A) these are polynomials degree by degree; in (B) they decompose, as in Theorem \ref{p5:thm:K}, into $\chi(X)$ and topological quantities $h^k(U;L_U)$ depending on the residue only through its integrality. Polynomial behaviour in all degrees occurs only in regular type.

%--------------------------------------------------------------
\section{Proof of the local formula for quasi-homogeneous isolated singularities}\label{p5:sec:local}
%--------------------------------------------------------------
In Part II the following formula was checked by computation for a range of spectra and left as a conjecture in general. We prove it in the quasi-homogeneous case.

\paragraph{Setting.} Let $h\in\C\{z,w\}$ be a quasi-homogeneous polynomial for positive weights $(w_z,w_w)$ and degree $N$ with an isolated singularity at the origin (so $h$ is reduced). Let $\sigma=h\,\partial_z\wedge\partial_w$, $D=\{h=0\}$, $\mathcal O\,h^s$ the module of residue $s$, and $\mathcal C_s=(\mathcal O\to\Theta\to\mathcal O)$ the stalk complex: $d_0f=(hf_w+sfh_w)\partial_z-(hf_z+sfh_z)\partial_w$, $d_1(p,q)=h(p_z+q_w)-(1-s)(h_zp+h_wq)$ (Part I \S I.6.6).

\begin{thm}[Local formula \tP\ \tQ]\label{p5:thm:local}
For $s\ne1$ the Poincar\'e polynomial of the cohomology of $\mathcal C_s$ is
\[\mu\,t^2\;+\;[\,s\in\Z,\ s\le0\,](1+t)\;+\;(t+t^2)\cdot\#\{\alpha\in\mathrm{Sp}:\ s+\alpha\in\Z,\ s+\alpha\le1\}.\]
\end{thm}
\begin{hyptab}
quasi-homogeneous & the Euler field $e$ with $e\,h=Nh$ exists; the complex splits by weights with finite-dimensional pieces. & Steps (2), (3).\\\midrule
isolated singularity & $h_z,h_w$ is a regular sequence (Hilbert series formula); $h$ is reduced. & Steps (5), (4).\\\midrule
$s\ne1$ & for $s=1$ the stalk is infinite dimensional (Part II). & dimension $\mu$ of the quotient $Q$ (step (1)).\\\midrule
$h$ polynomial (positive weights) & each weight piece $\mathcal O_k$ is a finite-dimensional space of polynomials. & Step (4).
\end{hyptab}
\begin{proof}
\emph{(1) Decomposition from earlier parts.} $B_s=(\mathcal O\to\Theta(-\log D)\to h\mathcal O)\cong(\Omega^\bullet(\log D),\,d+s\,\dlog h)$ is a subcomplex, and the quotient $Q=\mathcal C_s/B_s\simeq[J\!\cdot\!\mathcal O_D\to\mathcal O_D]$ has injective differential with $\dim\mathcal O/J=\mu$. So the stalk is $H^*(B_s)\oplus(\mathcal O/J)[-2]$ (as dimensions); this was proved in Part II v2 \tP. It remains to determine $\dim H^*(B_s)$. The sign convention for $s$ was fixed in Part II by matching computed data (see the checks at the end of this section).

\emph{(2) Weight decomposition.} $\mathcal O=\prod_k\mathcal O_k$ ($\mathcal O_k$ = quasi-homogeneous polynomials of weight $k$; since $h$ is quasi-homogeneous, a convergent series is a (convergent) sum of its weight components). Each term of $B_s$ is written with $\mathcal O_k$ and a basis of log forms ($\omega_e,\omega_\chi$: log 1-forms dual to the Euler field $e$ and Hamiltonian field $\chi_h$), and $\nabla_s$ preserves weight. The Lie derivative $L_e=\nabla\iota_e+\iota_e\nabla$ acts on the weight-$k$ part as the constant $k+sN$ (Cartan's formula), so for $k+sN\ne0$ the complex is contractible (homotopy $\iota_e/(k+sN)$), hence acyclic. Therefore $H^*(B_s)$ is concentrated in the weight $k_0=-sN$, and can be non-zero only if $k_0\in\Z_{\ge0}$ (so that $\mathcal O_{k_0}\ne0$).

\emph{(3) The weight-$k_0$ piece.} From the basic formulas
\[\nabla_sf=(ef+sNf)\,\omega_e+\chi_h(f)\,\omega_\chi,\qquad\nabla(a\,\omega_e+b\,\omega_\chi)=\tfrac{\nu}{N}\bigl[\chi_h(a)-(e+sN-\delta)\,b\bigr]\]
($\nu$ the basis 2-form), in weight $k_0$ we get $C^0=\mathcal O_{k_0}$, $C^1=\mathcal O_{k_0}\oplus\mathcal O_{k_0+\delta}$, $C^2=\mathcal O_{k_0+\delta}$, $d^0f=(0,\chi_hf)$, $d^1(a,b)=\chi_h(a)/N$. Put $K=\ker(\chi_h|\mathcal O_{k_0})$ and $\mathrm{Cok}=\mathrm{coker}(\chi_h:\mathcal O_{k_0}\to\mathcal O_{k_0+\delta})$. Then $H^0=K$, $H^1=K\oplus\mathrm{Cok}$, $H^2=\mathrm{Cok}$.

\emph{(4) Determination of $K$.} Let $f\in\mathcal O_{k_0}$ with $\chi_hf=0$. Up to sign $\chi_hf=(df\wedge dh)/(dz\wedge dw)$, so $df\wedge dh=0$. Contracting with $e$ gives $(ef)\,dh-(eh)\,df=k_0f\,dh-Nh\,df=0$, i.e.\ $d(f^N/h^{k_0})=0$ (for $k_0>0$). Hence $f^N=c\,h^{k_0}$ on the connected set $\C^2\setminus D$. Since $h$ is reduced, $N\,v_p(f)=k_0$ for every irreducible factor $p\mid h$. If $f\ne0$ then $N\mid k_0$ and $f/h^{k_0/N}$ is a weight-$0$ (rational) function, hence constant. So $K=\C\,h^{k_0/N}$ for $k_0\in N\Z_{\ge0}$ and $K=0$ otherwise (for $k_0=0$, $f$ is constant). \emph{This is an elementary verification, in the quasi-homogeneous case, of the standard fact that $\chi_hf=0\Rightarrow f\in\C\{h\}$ for isolated singularities.}

\emph{(5) $\dim\mathrm{Cok}$.} Let $d=k_0+\delta$, $a=w_z$, $b=w_w$. Then $\sum_k\dim\mathcal O_kx^k=1/((1-x^a)(1-x^b))$, and since $h_z,h_w$ is a regular sequence of degrees $N-a,N-b$, $\sum\dim T_kx^k=(1-x^{N-a})(1-x^{N-b})/((1-x^a)(1-x^b))$. From the identity $x^\delta+x^N-x^{N-a}-x^{N-b}=x^\delta(1-x^a)(1-x^b)$,
\[\frac{(1-x^{N-a})(1-x^{N-b})}{(1-x^a)(1-x^b)(1-x^N)}-\frac{1-x^\delta}{(1-x^a)(1-x^b)}=\frac{x^\delta}{1-x^N}.\]
Taking the coefficient of $x^d$: $\sum_{r\ge0}\dim T_{d-rN}-(\dim\mathcal O_d-\dim\mathcal O_{k_0})=[\,k_0\in N\Z_{\ge0}\,]=\dim K$. Since $\dim\mathrm{Cok}=\dim\mathcal O_d-\dim\mathcal O_{k_0}+\dim K$, we get $\dim\mathrm{Cok}=\sum_{r\ge0}\dim T_{d-rN}$.

\emph{(6) Translation to the spectrum.} Let $s=-k_0/N$. $\dim T_{d-rN}$ is the number of basis monomials $m$ of weight $\mathrm{wt}(m)=d-rN$. With $\alpha=(\mathrm{wt}(m)+a+b)/N$ we have $s+\alpha=1-r\in\Z_{\le1}$. Conversely each $\alpha$ with $s+\alpha\in\Z_{\le1}$ corresponds to exactly one $r\ge0$. Thus $\dim\mathrm{Cok}=\#\{\alpha\in\mathrm{Sp}:s+\alpha\in\Z,\ s+\alpha\le1\}$ and $\dim K=[\,s\in\Z_{\le0}\,]$. By (3) the Poincar\'e polynomial of $H^*(B_s)$ is $[s\in\Z_{\le0}](1+t)+(t+t^2)\dim\mathrm{Cok}$. Adding $\mu t^2$ from (1) gives the theorem.
\end{proof}
\paragraph{Checks and remaining unverified points.}
\begin{itemize}
\item \tV\ The identities and dimensions of steps (3)--(6) were compared, in exact rational arithmetic, for 14 types ($A_1,A_2,A_3,A_6,D_4,D_5,E_6,E_8$, $z^3+w^6$, $z^4+w^5$, $z^5+w^5$, $z^4+w^4$, $z^3+w^7$, $z^5+w^6$) and about 700 values of $k_0$: 0 mismatches (\texttt{cp2/roundH\_proof\_check.py}).
\item \tQ\ Sign convention of $s$: the sign of $s$ in $B_s\cong(\Omega^\bullet(\log D),d+s\,\dlog h)$ was fixed in Part II to match computed data. The agreement for 14 types $\times$ about 700 values would fail with the opposite sign, so it is a consistency check, but we have not written an independent derivation.
\item \tQ\ Convergence and completion in the weight decomposition: by quasi-homogeneity convergent series are sums of weight components and the complex splits as a product over weights; each weight is finite dimensional and $H^*(B_s)$ is concentrated in one weight, so the finite-dimensional stalk should equal the sum of the parts, but the comparison of formal and convergent series is not written out carefully.
\item Scope: $s\ne1$ and quasi-homogeneous isolated singularities only. Non-quasi-homogeneous isolated singularities, non-isolated singularities and $s=1$ are not covered.
\end{itemize}
Consequently the Part II v2 Observation ``local formula'' can be upgraded for quasi-homogeneous germs with $s\ne1$ from \tV\ \tK\ to \tP\ (apart from the \tQ\ items above); a proposed wording is recorded in \texttt{v3\_corrections.md}.

%--------------------------------------------------------------
\section{The index conjecture}\label{p5:sec:index}
%--------------------------------------------------------------
Hirzebruch--Riemann--Roch is a special case of the Atiyah--Singer index theorem, and a Riemann--Roch theorem in generalized geometry should be formulated as an index theorem for the generalized Dolbeault complex. We restrict to the log-symplectic case and separate what can and cannot be said now. \emph{Everything in this section is conjectural and unproved.}

\paragraph{Why existing index theorems do not apply.} (1) Hu's Kodaira--Nakano vanishing and Riemann--Roch on Lie groupoids (arXiv:2403.13424) \cite{Hu2} treats complexes attached to real Lie algebroids. Our complex $(\wedge^\bullet T\otimes L,\dbar+[\sigma,\cdot])$ is the complex of the complex Lie algebroid $\mathcal L=T^{0,1}\ltimes A^{1,0}$ (no real structure coming from $J$), a different object. (2) The Lie groupoid index of Pflaum--Posthuma--Tang \cite{PPT} requires a compact base, unimodularity, etc.\ \tQ. (3) We found no index theorem for log Dolbeault complexes weighted by the residue of $L$. Nearby literature: geometric quantization of b-symplectic manifolds (Braverman--Loizides--Song \cite{BLS}, boundary term vanishing in the limit), index formulas on boundary groupoids (Qiao--So \cite{QS}, no eta term in codimension $\ge2$), pseudodifferential operators with Lie structure at infinity (Ammann--Lauter--Nistor \cite{ALN}), fibred boundaries (Mazzeo--Melrose \cite{MM}), log Hodge theory (Hausel--Hunsicker--Mazzeo \cite{HHM}). All have different settings and do not apply as they stand.

\paragraph{Framework.} View $A=T_X(-\log D)$ as a (real) Lie algebroid (Gualtieri--Li \cite{GL}: integrable in the log-symplectic case, checked for smooth $D$ \tC\ \tQ) and regard the complex $\mathcal D_L=(\wedge^\bullet T\otimes L,\ \dbar+[\sigma,\cdot])$ on $U$ as a closed operator on weighted Sobolev spaces determined by $A$ (the weight being determined by the residue $r$). The principal symbol in the sense of $A$ is invertible (local computation in Part I \tV), so non-Fredholmness comes from the \emph{indicial roots} at the boundary. In the local model (smooth point of $D$, coordinate $z$) $\nabla z^k=(k-r)z^k\,dz/z$, so the indicial root is $k=r$, which exists iff $r\in\Z$.

\begin{conj}[Index conjecture \tK]\label{p5:conj:index}
Let $X,\sigma,D,L$ be as above (log-symplectic, SNC).
\begin{enumerate}
\item[(I1)] A closed extension of $\mathcal D_L$ on the weighted Sobolev space of weight $r$ is Fredholm if and only if $r_i\notin\Z$ for all $i$.
\item[(I2)] When Fredholm, the index (alternating sum) is $\chi(U)+(-1)^n\#D_{(n)}=(-1)^n\chi(X)$: an Atiyah--Singer-type integral $\int c_n(T(-\log D))$ plus the corner contribution $\#D_{(n)}$.
\item[(I3)] For $r_i\in\Z$, the difference of kernels and cokernels of the two closed extensions of $j_!$-type and $Rj_*$-type is given by a subquotient of $H^*(\partial U;L)$ (an APS-type wall-crossing triangle), and the index does not change since $\chi(\partial U)=0$ (consistent with Corollary \ref{p5:cor:wall}).
\item[(I4)] Graded refinement: the Poincar\'e polynomial $P_{\HH}(t)$ is given by the shape of Theorem \ref{p5:thm:K} ($\chi(X)t^n$ plus jump terms determined by the kernel of the indicial family).
\end{enumerate}
\end{conj}
\paragraph{Relation to what is already proved.} The Euler characteristic value in (I2) is proved by Theorems \ref{p5:thm:E} and \ref{p5:thm:K1} ($\chi(X)$). The triangle in (I3) is used at the topological level, $j_!L_U\to Rj_*L_U\to i_*i^*Rj_*L_U$, in Parts I--IV (fully checked for $n=2$); its index-theoretic counterpart is unproved. (I4) is a theorem for $n=2$ (Part I) \tP\ and, in general dimension, Theorem \ref{p5:thm:K} (under (GEN$'$)) \tP\ \tQ.

\paragraph{What remains to be proved.}
\begin{enumerate}
\item Fredholm theory on the complex Lie algebroid $\mathcal L$ (no $J$): construction of a parametrix and the ``$\Leftrightarrow$'' in (I1). Extending the Ammann--Lauter--Nistor framework to weights is natural but has not been done.
\item The index formula (I2): an analogue of the Atiyah--Singer formula on the b-tangent bundle $^bT$ \cite{BLS} for complexes with $\wedge^\bullet T$ coefficients and weights. The treatment of corners (strata of higher codimension) is the issue. The deformation argument of Qiao--So (no eta term) may help \tQ.
\item The APS-type wall-crossing (I3): the identification with $H^*(\partial U;L)$ beyond $n=2$ with smooth $D$ (Gysin sequence).
\item The graded version (I4): a graded refinement of spectral flow. We found no literature.
\item The algebraic hypothesis (GEN$'$) in general (open problems in \S\ref{p5:sec:ledger}).
\end{enumerate}

%--------------------------------------------------------------
\section{Relation to the original goal and a ledger of the whole series}\label{p5:sec:ledger}
%--------------------------------------------------------------
The original goal is a generalized-geometry version of the three treasures of complex geometry (Riemann--Roch, Serre vanishing, Kodaira embedding). We organize the results by whether each holds in the naive form, holds in a different form, or fails.

\begin{center}\small
\begin{tabular}{>{\raggedright\arraybackslash}p{2.6cm}>{\raggedright\arraybackslash}p{6.2cm}>{\raggedright\arraybackslash}p{6.0cm}}
\toprule Treasure & Form in which it holds & Fails / open\\\midrule
Riemann--Roch & Euler characteristic formula $(-1)^n\rk(E)\chi(X)$ (Thm.\ E), log version $\chi(U)$ (Thm.\ K1), Poincar\'e polynomial formulas (Part I for $n=2$, Thm.\ K), degree-wise polynomials in regular type (Part IV Cor.\ L) & naive HRR depending on $c_1(E)$ fails (Euler characteristic does not see $c_1$). Unification by an index theorem is unproved (Conjecture \ref{p5:conj:index})\\\midrule
Serre vanishing & vanishing in the transversal degree $a>0$ ($m\gg0$, $X$ projective, $\sigma$ arbitrary) (Part I) & vanishing in total degree fails (already for $\sigma=0$)\\\midrule
Kodaira embedding & leaf-space Kodaira; Poisson--Kodaira conjecture in the toric case (Part III) & embedding by generalized holomorphic sections is impossible (Part III). The general Poisson--Kodaira conjecture ($\mathbb P^2$, $m\ge3$, etc.) is unproved\\\bottomrule
\end{tabular}
\end{center}

\begin{center}\small
\begin{tabular}{>{\raggedright\arraybackslash}p{9.0cm}>{\raggedright\arraybackslash}p{1.9cm}>{\raggedright\arraybackslash}p{3.9cm}}
\toprule Claim & Status & Remark\\\midrule
Duality $\HH^k(L)^*\cong\HH^{2n-k}(L^\vee\otimes K^2)$ (any dimension) & \tP\ \tV & identity checked by machine for $n\le4$\\
Serre-type vanishing (transversal degree) & \tP & Part I\\
Surface theorem ($P_r(t)$, nodes only) & \tP\ \tV & Part I; cubic in $\mathbb P^2$, $m=-15,\dots,15$\\
Euler characteristic theorem (Thm.\ E) & \tP & Part II\\
Toric theorem T & \tP\ \tV & Part II; reproved in Part IV\\
Local formula (quasi-homogeneous, $s\ne1$) (Thm.\ \ref{p5:thm:local}) & \tP\ \tQ & this part; sign convention and convergence remain to be checked\\
Local formula (non-quasi-homogeneous, $s=1$) & \tK & unproved\\
Formula for non-reduced $D$ (Part II conjecture) & \tK\ \tV & fit to data\\
Reduction theorem (Thm.\ \ref{p5:thm:R}) & \tP\ \tQ & derivation of the filtration for $n>2$ is a sketch\\
(GEN$'$) in general & \tK & proved in the toric case\\
Regular type Theorem L / Cor.\ L & \tP\ \tQ & includes Nomizu-type hypotheses\\
Theorem K1 (log Euler characteristic) & \tP\ \tV & this part\\
Theorem K (RR-type Poincar\'e polynomial) & \tP\ \tV\ \tQ & this part; inherits \tQ\ from the reduction theorem\\
Impossibility of embedding by generalized holomorphic sections & \tP & Part III\\
Leaf-space Kodaira, toric Poisson--Kodaira & \tP\ \tK & Part III\\
Poisson--Kodaira (general) & \tK & unproved\\
Index conjecture (Conjecture \ref{p5:conj:index}) & \tK & this part; not proved\\
\bottomrule
\end{tabular}
\end{center}
\paragraph{Next steps.} (1) General sufficient conditions for (GEN$'$), and a reference for the holomorphic log-Darboux theorem. (2) Local formulas for non-reduced and non-quasi-homogeneous germs. (3) Poisson--Kodaira for $m\ge3$. (4) For the index conjecture, start with the Fredholm theory (I1). (5) Non-flat bundles and fibres without hard Lefschetz.

\paragraph{Scope of verification.} The computations in Parts I--V are exact linear algebra over finite fields or the rationals, or fits of known formulas. For Hu, Gualtieri--Li, Braverman--Loizides--Song, Qiao--So, Ammann--Lauter--Nistor, Mazzeo--Melrose, Hausel--Hunsicker--Mazzeo and Pflaum--Posthuma--Tang we only checked titles, authors and abstracts; the contents of theorems and journal details are unverified (\tQ).


